% Primitives d'une fonction — Mathématiques Tle A — Cours signé OnBuch+
% Programme officiel MINESEC. Compiler avec : tectonic MA05-primitives-a.tex
\newcommand{\DOCMATIERE}{Mathématiques}
\newcommand{\DOCNIVEAU}{Tle A}
\newcommand{\DOCMODULE}{Relations et opérations fondamentales dans l'ensemble des nombres réels}
\newcommand{\DOCLECON}{A05}
\newcommand{\DOCTITRE}{Primitives d'une fonction}
% OnBuch+ — préambule « cours signés » ENRICHI (compilé avec tectonic / XeLaTeX)
% Système visuel « Pop » d'OnBuch+ : fond crème (jamais blanc), bordures encre
% épaisses, ombres « dures » décalées, titres Archivo Black en capitales.
% Version MATHÉMATIQUES : graphes (pgfplots/tikz), boîtes pédagogiques
% (définition, propriété, méthode, attention, le savais-tu, exercice résolu,
% exercice, TP, bilan), page de garde et signature OnBuch+ ancrée en bas.
%
% Macros attendues dans main.tex AVANT \input{preamble} :
%   \DOCMATIERE  \DOCNIVEAU  \DOCMODULE  \DOCLECON (numéro)  \DOCTITRE
\documentclass[11pt]{article}

\usepackage{fontspec}
\usepackage[french]{babel}
\usepackage[a4paper,margin=2cm,top=3.4cm,bottom=2.8cm,headheight=1.4cm,footskip=1.3cm]{geometry}
\usepackage{amsmath,amssymb,amsfonts}
\usepackage{mathtools} % \xmapsto, \coloneqq... souvent utilisés par les modèles
\usepackage{cancel} % simplifications barrées
% \abs{z} (module, valeur absolue) et \norm{u} (norme), utilisés par les modèles
\providecommand{\abs}{}\let\abs\relax\DeclarePairedDelimiter{\abs}{\lvert}{\rvert}
\providecommand{\norm}{}\let\norm\relax\DeclarePairedDelimiter{\norm}{\lVert}{\rVert}
\usepackage[table]{xcolor} % [table] : fournit aussi \rowcolors (alternance de lignes)
\usepackage{pagecolor}
\usepackage{tikz}
\usetikzlibrary{arrows.meta,positioning,calc,shapes.geometric,shapes.misc,shapes.arrows,decorations.pathmorphing,decorations.markings,patterns,shadows,mindmap,trees,fit,backgrounds,matrix,intersections}
\usepackage{pgfplots}
\usepackage{pgf-pie} % diagrammes circulaires \pie (statistiques)
\pgfplotsset{compat=1.18}
\usepgfplotslibrary{fillbetween,groupplots} % aires entre courbes (calcul intégral)
\usepackage{enumitem}
\usepackage{fancyhdr}
% [most] charge toutes les bibliothèques tcolorbox (listings, minted, raster,
% documentation...) et gonfle inutilement la mémoire TeX sur un document long
% (~300 boîtes) : on ne charge que ce qui est réellement utilisé.
\usepackage[skins,breakable]{tcolorbox}
\usepackage{graphicx}
\usepackage{colortbl}
\usepackage{booktabs}
\usepackage{tabularx}
\usepackage{diagbox} % case d'en-tête diagonale des tableaux à double entrée
% Colonnes extensibles centrée / gauche / droite (C, L, R), souvent utilisées par les modèles
\newcolumntype{C}{>{\centering\arraybackslash}X}
\newcolumntype{L}{>{\raggedright\arraybackslash}X}
\newcolumntype{R}{>{\raggedleft\arraybackslash}X}
\usepackage{array}
\usepackage{multicol}
\usepackage{siunitx}
% parse-numbers=false : accepte aussi \SI{2,5 \times 10^{-3}}{...}
\sisetup{locale=FR,output-decimal-marker={,},group-minimum-digits=4,parse-numbers=false}
\DeclareSIUnit\ohm{\ensuremath{\symup{\Omega}}} % U+2126 absent de la police
\usepackage{qrcode}
% \degree : commande fréquemment utilisée par les modèles pour un angle ou une
% température en dehors de \SI{}{} (non fournie par les paquets chargés).
\providecommand{\degree}{^{\circ}}
% \qed (fin de démonstration), utilisé par les modèles sans amsthm
\providecommand{\qed}{\hfill$\square$}
% \barre : double barre des valeurs interdites dans un tableau de variations (tkz-tab)
\providecommand{\barre}{\ensuremath{\|}}
% environnement proof (amsthm non chargé) utilisé par les modèles
\ifdefined\proof\else\newenvironment{proof}[1][Démonstration]{\par\noindent\textit{#1.}\ }{\qed\par}\fi
\usepackage{lastpage}
\usepackage{hyperref}
\hypersetup{hidelinks}

% ============================================================
% Palette « Pop » OnBuch+ — mêmes hex que la classe P (plus_ui.dart)
% ============================================================
\definecolor{poporange}{HTML}{F59321}
\definecolor{popdark}{HTML}{E07A0C}
\definecolor{popcream}{HTML}{FFF6E8}
\definecolor{popink}{HTML}{1C1714}
\definecolor{poppurple}{HTML}{7A5AE0}
\definecolor{popgreen}{HTML}{2E9E6A}
\definecolor{poppink}{HTML}{E0457B}
\definecolor{popblue}{HTML}{2F7DE1}
\definecolor{popgold}{HTML}{C98A12}
\definecolor{popmuted}{HTML}{8A7F76}
\definecolor{popline}{HTML}{EADFD2}
\definecolor{popgray}{HTML}{8A7F76}
\colorlet{popgrey}{popgray}
% Alias tolérants (le modèle nomme parfois les couleurs différemment)
\colorlet{popwhite}{white}
\colorlet{popblack}{popink}
\colorlet{popred}{poppink}
\colorlet{popyellow}{popgold}
% Teintes claires pour fonds de tableaux / zones
\colorlet{popblueL}{popblue!10!white}
\colorlet{poporangeL}{poporange!14!white}
\colorlet{popgreenL}{popgreen!12!white}
\colorlet{poppurpleL}{poppurple!12!white}
\colorlet{poppinkL}{poppink!12!white}
\colorlet{popgoldL}{popgold!14!white}

\pagecolor{popcream}

% ============================================================
% Typographie
% ============================================================
\defaultfontfeatures{Ligatures=TeX}
\setmainfont{PlusJakartaSans-Medium.ttf}[Path=../../../fonts/,BoldFont=PlusJakartaSans-Bold.ttf]
% Maths en sans-serif (Fira Math) avec chiffres et lettres droites de Plus
% Jakarta, pour que les formules se fondent dans le texte.
\usepackage[math-style=ISO,bold-style=ISO]{unicode-math}
\setmathfont{FiraMath-Regular.otf}[Path=../../../fonts/]
\setmathfont{PlusJakartaSans-Medium.ttf}[Path=../../../fonts/,range={up/num,up/{Latin,latin},\mathup}]
\usepackage{icomma} % virgule décimale sans espace en mode math
% Caractères mathématiques Unicode absents des polices de texte (Plus Jakarta,
% Archivo Black) : rendus par leur équivalent mathématique, en texte comme en
% formule — sinon ils s'affichent en carré vide (ex. « Arithmétique dans ℤ »).
\usepackage{newunicodechar}
\newunicodechar{ℝ}{\ensuremath{\mathbb{R}}}
\newunicodechar{ℤ}{\ensuremath{\mathbb{Z}}}
\newunicodechar{ℂ}{\ensuremath{\mathbb{C}}}
\newunicodechar{ℕ}{\ensuremath{\mathbb{N}}}
\newunicodechar{ℚ}{\ensuremath{\mathbb{Q}}}
\newunicodechar{𝔻}{\ensuremath{\mathbb{D}}}
\newunicodechar{α}{\ensuremath{\alpha}}
\newunicodechar{β}{\ensuremath{\beta}}
\newunicodechar{γ}{\ensuremath{\gamma}}
\newunicodechar{δ}{\ensuremath{\delta}}
\newunicodechar{ε}{\ensuremath{\varepsilon}}
\newunicodechar{θ}{\ensuremath{\theta}}
\newunicodechar{λ}{\ensuremath{\lambda}}
\newunicodechar{μ}{\ensuremath{\mu}}
\newunicodechar{σ}{\ensuremath{\sigma}}
\newunicodechar{τ}{\ensuremath{\tau}}
\newunicodechar{φ}{\ensuremath{\varphi}}
\newunicodechar{ψ}{\ensuremath{\psi}}
\newunicodechar{ω}{\ensuremath{\omega}}
\newunicodechar{Σ}{\ensuremath{\Sigma}}
% majuscules produites par \MakeUppercase dans les titres (α-aminés -> Α)
\newunicodechar{Α}{\ensuremath{\alpha}}
\newunicodechar{Β}{\ensuremath{\beta}}
\newunicodechar{Γ}{\ensuremath{\Gamma}}
\newunicodechar{Δ}{\ensuremath{\Delta}}
\newunicodechar{Ε}{\ensuremath{\varepsilon}}
\newunicodechar{Θ}{\ensuremath{\Theta}}
\newunicodechar{Λ}{\ensuremath{\Lambda}}
\newunicodechar{Μ}{\ensuremath{\mu}}
\newunicodechar{Τ}{\ensuremath{\tau}}
\newunicodechar{Φ}{\ensuremath{\Phi}}
\newunicodechar{Ψ}{\ensuremath{\Psi}}
\newunicodechar{Ω}{\ensuremath{\Omega}}
\newunicodechar{↦}{\ensuremath{\mapsto}}
\newunicodechar{∈}{\ensuremath{\in}}
\newunicodechar{∉}{\ensuremath{\notin}}
\newunicodechar{∧}{\ensuremath{\wedge}}
\newunicodechar{⇔}{\ensuremath{\Leftrightarrow}}
\newunicodechar{⟺}{\ensuremath{\Longleftrightarrow}}
\newunicodechar{⇒}{\ensuremath{\Rightarrow}}
\newunicodechar{⟹}{\ensuremath{\Longrightarrow}}
\newunicodechar{∘}{\ensuremath{\circ}}
\newunicodechar{∪}{\ensuremath{\cup}}
\newunicodechar{∩}{\ensuremath{\cap}}
\newunicodechar{≡}{\ensuremath{\equiv}}
\newunicodechar{⊂}{\ensuremath{\subset}}
\newunicodechar{⊥}{\ensuremath{\perp}}
\newunicodechar{∃}{\ensuremath{\exists}}
\newunicodechar{∀}{\ensuremath{\forall}}
\newunicodechar{∛}{\ensuremath{\sqrt[3]{\,}}}
\newunicodechar{∜}{\ensuremath{\sqrt[4]{\,}}}
\newunicodechar{⃗}{\ensuremath{{}^{\to}}}
\newunicodechar{ⁿ}{\ifmmode^{n}\else\textsuperscript{\ensuremath{n}}\fi}
\newunicodechar{ˣ}{\ifmmode^{x}\else\textsuperscript{\ensuremath{x}}\fi}
\newunicodechar{ᵇ}{\ifmmode^{b}\else\textsuperscript{\ensuremath{b}}\fi}
\newunicodechar{ʳ}{\ifmmode^{r}\else\textsuperscript{\ensuremath{r}}\fi}
\newunicodechar{ᵘ}{\ifmmode^{u}\else\textsuperscript{\ensuremath{u}}\fi}
\newunicodechar{ʸ}{\ifmmode^{y}\else\textsuperscript{\ensuremath{y}}\fi}
\newunicodechar{ᵃ}{\ifmmode^{a}\else\textsuperscript{\ensuremath{a}}\fi}
\newunicodechar{ᵉ}{\ifmmode^{e}\else\textsuperscript{\ensuremath{e}}\fi}
\newunicodechar{ᵏ}{\ifmmode^{k}\else\textsuperscript{\ensuremath{k}}\fi}
\newunicodechar{ᵖ}{\ifmmode^{p}\else\textsuperscript{\ensuremath{p}}\fi}
\newunicodechar{ᴺ}{\ifmmode^{N}\else\textsuperscript{\ensuremath{N}}\fi}
\newunicodechar{ᶜ}{\ifmmode^{c}\else\textsuperscript{\ensuremath{c}}\fi}
\newunicodechar{ʰ}{\ifmmode^{h}\else\textsuperscript{\ensuremath{h}}\fi}
\newunicodechar{ᵠ}{\ifmmode^{q}\else\textsuperscript{\ensuremath{q}}\fi}
\newunicodechar{ₙ}{\ifmmode_{n}\else\textsubscript{\ensuremath{n}}\fi}
\newunicodechar{ₐ}{\ifmmode_{a}\else\textsubscript{\ensuremath{a}}\fi}
\newunicodechar{ᵢ}{\ifmmode_{i}\else\textsubscript{\ensuremath{i}}\fi}
\newunicodechar{ᵣ}{\ifmmode_{r}\else\textsubscript{\ensuremath{r}}\fi}
\newunicodechar{ₖ}{\ifmmode_{k}\else\textsubscript{\ensuremath{k}}\fi}
\newunicodechar{ᵦ}{\ifmmode_{\beta}\else\textsubscript{\ensuremath{\beta}}\fi}
\newunicodechar{ₓ}{\ifmmode_{x}\else\textsubscript{\ensuremath{x}}\fi}
\newfontfamily\popdisplay{ArchivoBlack-Regular.ttf}[Path=../../../fonts/]
\newfontfamily\poplogo{SpaceGrotesk-Bold.ttf}[Path=../../../fonts/]
\newfontfamily\popsemi{PlusJakartaSans-SemiBold.ttf}[Path=../../../fonts/]
\newfontfamily\popxbold{PlusJakartaSans-ExtraBold.ttf}[Path=../../../fonts/]
\newfontfamily\popxboldls{PlusJakartaSans-ExtraBold.ttf}[Path=../../../fonts/,LetterSpace=3.0]

\setlist[enumerate]{leftmargin=1.7em,itemsep=5pt,topsep=4pt}
\setlist[itemize]{leftmargin=1.7em,itemsep=4pt,topsep=4pt,label={\color{poporange}\textbullet}}
\setlength{\parindent}{0pt}
\setlength{\parskip}{5pt}
\renewcommand{\arraystretch}{1.25}

% pgfplots : style Pop par défaut
\pgfplotsset{
  every axis/.append style={axis line style={popink,line width=1pt},tick style={popink},
    grid style={popline,line width=0.5pt},label style={font=\small},tick label style={font=\footnotesize}},
  popcurve/.style={poporange,line width=1.8pt,no markers,smooth},
}
% Même style disponible dans un \draw TikZ simple (hors pgfplots)
\tikzset{popcurve/.style={poporange,line width=1.8pt}}

% ============================================================
% Pill / squiggle
% ============================================================
\newcommand{\poppill}[3]{%
  \tikz[baseline=-0.55ex]\node[draw=popink,line width=1.2pt,rounded corners=2.6pt,%
    fill=#2,inner xsep=6.5pt,inner ysep=3pt]{%
    \popxboldls\fontsize{7}{7}\selectfont\color{#3}\MakeUppercase{#1}};%
}
\newcommand{\popsquiggle}[1]{%
  \tikz[baseline=-0.4ex]{
    \draw[#1,line width=2.6pt,line cap=round]
      plot [smooth] coordinates {(0,0.09) (0.34,0.01) (0.68,0.21) (1.02,0.01) (1.36,0.10)};
  }%
}
% Mot-clé surligné (marqueur orange sous le texte)
\newcommand{\cle}[1]{{\popxbold\color{popdark}#1}}

% ============================================================
% En-tête / pied
% ============================================================
\pagestyle{fancy}
\fancyhf{}
\renewcommand{\headrulewidth}{0pt}
\renewcommand{\footrulewidth}{0pt}
\lhead{\raisebox{-0.15ex}{\poplogo\fontsize{13}{13}\selectfont\color{popink}OnBuch\color{poporange}+}}
\rhead{\poppill{\DOCMATIERE\ \textbullet\ \DOCNIVEAU\ \textbullet\ Leçon \DOCLECON}{white}{popink}}
\lfoot{\popxbold\fontsize{7.6}{7.6}\selectfont\color{popmuted}\MakeUppercase{Cours signé OnBuch+}\ \textbullet\ {\popsemi\DOCTITRE}}
\rfoot{\tikz[baseline=-0.6ex]\node[rounded corners=8.5pt,fill=popink,minimum height=17pt,minimum width=17pt,inner xsep=5pt,inner sep=0]{\popxbold\fontsize{8}{8}\selectfont\color{popcream}\thepage\,/\,\pageref*{LastPage}};}

% ============================================================
% Titres
% ============================================================
\newcommand{\chaptitre}[2]{%
  \begin{tcolorbox}[enhanced,colback=poporange,colframe=popink,boxrule=2.6pt,arc=11pt,
      drop shadow=popink,left=15pt,right=15pt,top=12pt,bottom=11pt,before skip=3pt,after skip=18pt]
    \poppill{#2}{popink}{popcream}\par\vspace{9pt}
    {\raggedright\hyphenpenalty=10000\exhyphenpenalty=10000\popdisplay\fontsize{22}{24}\selectfont\color{popcream}\MakeUppercase{#1}\par}
  \end{tcolorbox}
}
% Section numérotée : pastille orange avec le numéro + titre Archivo
\newcounter{popsec}
\newcommand{\coursec}[1]{%
  \stepcounter{popsec}\setcounter{popsub}{0}%
  \par\vspace{16pt}\noindent
  \begin{minipage}{\linewidth}
  \tikz[baseline=-0.7ex]\node[draw=popink,line width=1.6pt,fill=poporange,rounded corners=4pt,minimum size=22pt,inner sep=0]{\popdisplay\fontsize{12}{12}\selectfont\color{popcream}\thepopsec};%
  \hspace{8pt}{\popdisplay\fontsize{15}{17}\selectfont\color{popink}\MakeUppercase{#1}}\par
  \vspace{4pt}\noindent{\color{poporange}\rule{36pt}{3.6pt}}
  \end{minipage}\par\nopagebreak\vspace{8pt}
}
\newcounter{popsub}
\newcommand{\courssub}[1]{%
  \stepcounter{popsub}%
  \par\vspace{9pt}\noindent{\popxbold\fontsize{11.5}{13}\selectfont\color{popink}{\color{poporange}\thepopsec.\thepopsub}\hspace{6pt}#1}\par\nopagebreak\vspace{4pt}
}

% ============================================================
% Boîtes pédagogiques — même recette : fond blanc, bordure épaisse colorée,
% ombre dure encre, badge plein. Titre optionnel [#1].
% ============================================================
\tcbset{popbase/.style={breakable,enhanced,colback=white,boxrule=2.3pt,arc=9pt,
  shadow={3pt}{-3pt}{0pt}{popink},left=14pt,right=14pt,top=11pt,bottom=11pt,before skip=11pt,after skip=15pt,
  fontupper=\popsemi\fontsize{10.6}{14.5}\selectfont\color{popink},
  fontlower=\popsemi\fontsize{10.6}{14.5}\selectfont\color{popink}}}
% Coche dessinée (le glyphe ✓ n'existe pas dans les polices du document)
\newcommand{\popcheck}[1]{\tikz[baseline=-0.5ex]\draw[#1,line width=1.6pt,line cap=round,line join=round](-0.13,0.0)--(-0.04,-0.09)--(0.13,0.1);}
\newcommand{\popboxhead}[3]{\poppill{#1}{#2}{white}\ifx\relax#3\relax\else\hspace{6pt}{\popxbold\fontsize{10.6}{12}\selectfont\color{popink}#3}\fi\par\vspace{7pt}}

\newtcolorbox{aretenir}[1][]{popbase,colframe=popgreen,before upper={\popboxhead{À retenir}{popgreen}{#1}}}
\newtcolorbox{exemplebox}[1][]{popbase,colframe=poporange,before upper={\popboxhead{Exemple}{poporange}{#1}}}
\newtcolorbox{definition}[1][]{popbase,colframe=popblue,before upper={\popboxhead{Définition}{popblue}{#1}}}
\newtcolorbox{propriete}[1][]{popbase,colframe=poppurple,before upper={\popboxhead{Propriété}{poppurple}{#1}}}
\newtcolorbox{methode}[1][]{popbase,colframe=popgold,before upper={\popboxhead{Méthode}{popgold}{#1}}}
\newtcolorbox{attention}[1][]{popbase,colframe=poppink,before upper={\popboxhead{Attention}{poppink}{#1}}}
\newtcolorbox{experience}[1][]{popbase,colframe=popblue,colback=popblueL,before upper={\popboxhead{Expérience}{popblue}{#1}}}
% « Le savais-tu ? » : carte violette pleine (contexte, culture, Cameroun)
\newtcolorbox{savaistu}[1][]{popbase,colback=poppurple,colframe=popink,
  fontupper=\popsemi\fontsize{10.4}{14.2}\selectfont\color{white},
  before upper={\poppill{Le savais-tu\,?}{white}{poppurple}\ifx\relax#1\relax\else\hspace{6pt}{\popxbold\color{white}#1}\fi\par\vspace{7pt}}}
% Exercice résolu : énoncé en haut, \tcblower puis la solution sur fond crème
\newcounter{popexo}\newcounter{popexr}
\newtcolorbox{exoresolu}[1][]{popbase,colframe=poporange,colbacklower=popcream,
  segmentation style={popink,line width=1.2pt,dashed},
  before upper={\stepcounter{popexr}\popboxhead{Exercice résolu \thepopexr}{poporange}{#1}},
  before lower={\poppill{Solution}{popink}{popcream}\par\vspace{6pt}}}
% Exercice (non résolu en place) — la correction est dans la partie Corrigés
\newtcolorbox{exercice}[1][]{popbase,colframe=popink,
  before upper={\stepcounter{popexo}\popboxhead{Exercice \thepopexo}{popink}{#1}}}
% Corrigé
\newtcolorbox{corrige}[1][]{popbase,colframe=popgreen,colback=popgreenL,
  before upper={\popboxhead{Corrigé}{popgreen}{#1}}}
% Objectifs / prérequis / bilan
\newtcolorbox{objectifs}{popbase,colframe=popink,colback=poporangeL,
  before upper={\popboxhead{Objectifs de la leçon}{popink}{}}}
\newtcolorbox{prerequis}{popbase,colframe=popink,colback=popblueL,
  before upper={\popboxhead{Prérequis}{popblue}{}}}
\newtcolorbox{bilan}{popbase,colframe=popink,colback=white,boxrule=2.8pt,
  before upper={\popboxhead{Fiche bilan}{poporange}{L'essentiel à connaître pour l'examen}}}
% Figure encadrée avec légende
\newtcolorbox{popfigure}[1][]{enhanced,breakable=false,colback=white,colframe=popline,boxrule=1.4pt,arc=8pt,
  left=10pt,right=10pt,top=10pt,bottom=8pt,before skip=10pt,after skip=12pt,halign=center,
  lower separated=false,halign lower=center,
  fontlower=\popsemi\fontsize{9}{11.5}\selectfont\color{popmuted},
  #1}
\newcommand{\legende}[1]{\par\vspace{6pt}{\popsemi\fontsize{9}{11.5}\selectfont\color{popmuted}#1}}

% ============================================================
% Signature OnBuch+
% ============================================================
\newcommand{\poplogoinline}[1]{{\poplogo\fontsize{#1}{#1}\selectfont\color{popink}OnBuch\color{poporange}+}}
\newcommand{\signatureonbuch}{%
  \begin{tcolorbox}[enhanced,colback=popink,colframe=popink,boxrule=2.2pt,arc=10pt,
      drop shadow=poporange,left=14pt,right=14pt,top=10pt,bottom=10pt,before skip=0pt,after skip=0pt]
    \tikz[baseline=-0.7ex]\node[circle,fill=popgreen,draw=popcream,line width=1.3pt,minimum size=22pt,inner sep=0]{\popcheck{white}};%
    \hspace{9pt}\begin{minipage}[c]{0.62\linewidth}
      {\popxbold\fontsize{12}{13}\selectfont\color{popcream}Signé\ \textbullet\ {\poplogo OnBuch\color{poporange}+}}\par\vspace{2pt}
      {\raggedright\popsemi\fontsize{8}{10.5}\selectfont\color{popline}\MakeUppercase{Équipe pédagogique OnBuch+}\\\MakeUppercase{Conforme au programme officiel MINESEC}\par}
    \end{minipage}\hfill
    {\poplogo\fontsize{22}{22}\selectfont\color{popcream}OnBuch\color{poporange}+}
  \end{tcolorbox}%
}

% ============================================================
% Page de garde
% #1 titre · #2 matière · #3 niveau · #4 module · #5 n° leçon · #6 durée ·
% #7 description / objectifs (texte ou itemize)
% ============================================================
\newcommand{\pagegarde}[7]{%
  \thispagestyle{empty}
  \newgeometry{a4paper,margin=1.7cm,top=1.6cm,bottom=1.4cm}%
  \begin{tikzpicture}[remember picture,overlay]
    \fill[poporange!18] ([xshift=-3.2cm,yshift=-3.6cm]current page.north east) circle (4.6cm);
    \fill[poppurple!14] ([xshift=2.4cm,yshift=7.6cm]current page.south west) circle (3.8cm);
  \end{tikzpicture}%
  {\poplogo\fontsize{30}{30}\selectfont\color{popink}OnBuch\color{poporange}+}\hfill
  \raisebox{0.6ex}{\poppill{Cours signé \textbullet\ Premium}{popink}{popcream}}\par
  \vspace{1pt}\popsquiggle{poppurple}\par
  \vspace{8pt}
  \begin{tcolorbox}[enhanced,colback=poporange,colframe=popink,boxrule=2.8pt,arc=13pt,
      drop shadow=popink,left=18pt,right=18pt,top=12pt,bottom=13pt,after skip=10pt]
    \poppill{#2 \textbullet\ #3}{popink}{popcream}\hspace{5pt}\poppill{#4}{white}{popink}\par\vspace{8pt}
    {\popdisplay\fontsize{14}{14}\selectfont\color{popink}LEÇON #5}\par\vspace{4pt}
    {\raggedright\hyphenpenalty=10000\exhyphenpenalty=10000\popdisplay\fontsize{25}{27}\selectfont\color{popcream}\MakeUppercase{#1}\par}
    \vspace{8pt}
    {\popxbold\fontsize{9.5}{9.5}\selectfont\color{popink}\MakeUppercase{Durée conseillée : #6}}
  \end{tcolorbox}
  \begin{tcolorbox}[enhanced,colback=white,colframe=popink,boxrule=2.2pt,arc=10pt,
      drop shadow=popink,left=16pt,right=16pt,top=10pt,bottom=10pt,after skip=8pt]
    \poppill{Dans cette leçon}{poppurple}{white}\par\vspace{5pt}
    \popsemi\fontsize{9.3}{12.6}\selectfont\color{popink}#7
  \end{tcolorbox}
  \noindent\begin{minipage}[t]{0.487\linewidth}
    \begin{tcolorbox}[enhanced,colback=popgreen,colframe=popink,boxrule=2.2pt,arc=10pt,
        drop shadow=popink,left=11pt,right=11pt,top=10pt,bottom=10pt,height=3.1cm,valign=center]
      \tcbox[colback=white,colframe=popink,boxrule=1.3pt,arc=5pt,boxsep=0pt,nobeforeafter,
        left=3pt,right=3pt,top=3pt,bottom=3pt,valign=center]{\qrcode[height=1.9cm]{https://play.google.com/store/apps/details?id=com.luvvix.onbuch&hl=fr}}%
      \hspace{8pt}\begin{minipage}[c]{0.5\linewidth}
        {\popdisplay\fontsize{11}{12.5}\selectfont\color{white}\MakeUppercase{Télécharge OnBuch}}\par\vspace{4pt}
        {\raggedright\popsemi\fontsize{8.4}{11}\selectfont\color{white}Gratuit \textbullet\ Google Play\\Tuteur IA, annales, concours, résultats\par}
      \end{minipage}
    \end{tcolorbox}
  \end{minipage}\hfill
  \begin{minipage}[t]{0.487\linewidth}
    \begin{tcolorbox}[enhanced,colback=poppurple,colframe=popink,boxrule=2.2pt,arc=10pt,
        drop shadow=popink,left=11pt,right=11pt,top=10pt,bottom=10pt,height=3.1cm,valign=center]
      \tikz[baseline=-0.7ex]\node[circle,fill=white,draw=popink,line width=1.3pt,minimum size=1.6cm,inner sep=0]{\popdisplay\fontsize{24}{24}\selectfont\color{poppurple}+};%
      \hspace{8pt}\begin{minipage}[c]{0.6\linewidth}
        {\popdisplay\fontsize{11}{12.5}\selectfont\color{white}\MakeUppercase{Passe à OnBuch+}}\par\vspace{4pt}
        {\raggedright\popsemi\fontsize{8.4}{11}\selectfont\color{white}Tous les cours signés, Tuteur IA illimité, fiches PDF — onglet OnBuch+ de l'app\par}
      \end{minipage}
    \end{tcolorbox}
  \end{minipage}
  \vspace{8pt}
  \popcontenu
  \vfill
  \signatureonbuch
  \restoregeometry
  \newpage
}

% Rangée « Ce cours contient » (page de garde)
\newcommand{\poptile}[3]{% #1 couleur #2 gros texte #3 légende
  \begin{tcolorbox}[enhanced,colback=#1,colframe=popink,boxrule=2pt,arc=9pt,shadow={2.5pt}{-2.5pt}{0pt}{popink},
      left=6pt,right=6pt,top=9pt,bottom=9pt,halign=center,valign=center,height=1.75cm,width=\linewidth,nobeforeafter]
    {\popdisplay\fontsize{12.5}{13}\selectfont\color{white}\MakeUppercase{#2}}\par\vspace{3pt}
    {\centering\popsemi\fontsize{7.8}{9.5}\selectfont\color{white}#3\par}
  \end{tcolorbox}}
\newcommand{\popcontenu}{%
  \par\noindent{\popxboldls\fontsize{7.5}{7.5}\selectfont\color{popmuted}CE COURS SIGNÉ CONTIENT}\par\vspace{7pt}
  \noindent\begin{minipage}[t]{0.235\linewidth}\poptile{popblue}{Cours}{complet, illustré, pas à pas}\end{minipage}\hfill
  \begin{minipage}[t]{0.235\linewidth}\poptile{poporange}{Méthodes}{et exercices résolus}\end{minipage}\hfill
  \begin{minipage}[t]{0.235\linewidth}\poptile{poppink}{Exercices}{type examen + corrigés}\end{minipage}\hfill
  \begin{minipage}[t]{0.235\linewidth}\poptile{popgreen}{Fiche bilan}{l'essentiel à retenir}\end{minipage}\par}

% Sommaire
\newtcolorbox{sommaire}{enhanced,colback=white,colframe=popink,boxrule=2.2pt,arc=10pt,
  drop shadow=popink,left=18pt,right=18pt,top=13pt,bottom=13pt,before skip=6pt,after skip=20pt,
  before upper={\poppill{Sommaire}{poppurple}{white}\par\vspace{9pt}\popsemi\fontsize{10.6}{14.5}\selectfont\color{popink}}}

% Signature de fin de cours — poussée tout en bas de la dernière page
\newcommand{\findecours}{%
  \par\vfill\nopagebreak
  \begin{center}\popsquiggle{poporange}\end{center}\vspace{4pt}
  \signatureonbuch
}
\AtBeginDocument{\def\llbracket{\mathopen{[\mkern-3.5mu[}}\def\rrbracket{\mathclose{]\mkern-3.5mu]}}} % intervalles d'entiers (glyphe ⟦ absent de la police)
% Glyphes absents de Fira Math : redéfinis à partir de glyphes disponibles
\AtBeginDocument{%
  \def\setminus{\mathbin{\backslash}}\let\smallsetminus\setminus
  \def\vdots{\mathord{\vbox{\baselineskip3.2pt\lineskiplimit0pt\kern4pt\hbox{.}\hbox{.}\hbox{.}}}}%
  \def\ddots{\mathinner{\mkern1mu\raise6.5pt\hbox{.}\mkern2mu\raise3.6pt\hbox{.}\mkern2mu\raise0.7pt\hbox{.}\mkern1mu}}%
  \def\triangle{\mathord{\scalebox{0.9}{$\symup{\Delta}$}}}%
  \def\nabla{\mathord{\raisebox{\depth}{\scalebox{1}[-1]{$\symup{\Delta}$}}}}%
  \def\checkmark{\popcheck{popdark}}%
  \def\star{\mathbin{\ast}}\def\bigstar{\mathord{\ast}}%
  \def\diamond{\mathbin{\scalebox{0.6}{$\Diamond$}}}%
}

%% FIGFIX-BEGIN v3 — correctifs globaux des figures (docs/onbuch-generation/tools/figfix)
\usepackage{adjustbox}\usepackage{etoolbox}
% toute figure TikZ/pgfplots plus large que la ligne est réduite à la largeur disponible
\BeforeBeginEnvironment{tikzpicture}{\begin{adjustbox}{max width=\linewidth}}
\AfterEndEnvironment{tikzpicture}{\end{adjustbox}}
% légendes pgfplots lisibles par-dessus les courbes
\pgfplotsset{every axis/.append style={legend style={fill=white,fill opacity=0.92,text opacity=1,draw=black!15,inner sep=3pt}}}
% étiquettes d'arêtes (syntaxe edge["..."]) avec fond blanc
\tikzset{every edge quotes/.append style={fill=white,inner sep=1.2pt}}
% bulles de cartes mentales : texte à la ligne dans la bulle, bulle un peu plus grande
\tikzset{every concept/.append style={text width=2.0cm,align=center,minimum size=2.3cm,inner sep=2pt}}
%% FIGFIX-END
\begin{document}

\pagegarde{Primitives d'une fonction}{Mathématiques}{Tle A}{Relations et opérations fondamentales dans l'ensemble des nombres réels}{A05}{4 h}{Cette leçon introduit la notion de primitive d'une fonction sur un intervalle, opération inverse de la dérivation. On y apprend à déterminer les primitives des fonctions usuelles par lecture inverse du tableau des dérivées et à trouver la primitive unique satisfaisant une condition initiale F(a) = b.
\vspace{4pt}
\begin{multicols}{2}\small
\begin{itemize}
\item À la fin de cette leçon, je sais définir une primitive d'une fonction sur un intervalle
\item À la fin de cette leçon, je sais vérifier qu'une fonction F est une primitive d'une fonction f en calculant F'
\item À la fin de cette leçon, je sais énoncer et utiliser le fait que deux primitives d'une même fonction diffèrent d'une constante
\item À la fin de cette leçon, je sais déterminer une primitive d'une fonction usuelle par lecture inverse du tableau des dérivées
\item À la fin de cette leçon, je sais déterminer une primitive d'une somme ou d'un produit par une constante
\item À la fin de cette leçon, je sais déterminer la primitive qui prend une valeur donnée en un point (condition initiale)
\end{itemize}\end{multicols}}

\begin{sommaire}
\begin{multicols}{2}
\begin{enumerate}
\item Notion de primitive
\item Ensemble des primitives d'une fonction
\item Primitives des fonctions usuelles
\item Opérations sur les primitives
\item Primitive vérifiant une condition initiale
\item Applications et synthèse
\item Activité d'intégration
\item Méthodes et exercices résolus
\item Exercices
\item Corrigés des exercices
\item Fiche bilan
\end{enumerate}
\end{multicols}
\end{sommaire}

\begin{objectifs}
\begin{itemize}
\item À la fin de cette leçon, je sais définir une primitive d'une fonction sur un intervalle
\item À la fin de cette leçon, je sais vérifier qu'une fonction F est une primitive d'une fonction f en calculant F'
\item À la fin de cette leçon, je sais énoncer et utiliser le fait que deux primitives d'une même fonction diffèrent d'une constante
\item À la fin de cette leçon, je sais déterminer une primitive d'une fonction usuelle par lecture inverse du tableau des dérivées
\item À la fin de cette leçon, je sais déterminer une primitive d'une somme ou d'un produit par une constante
\item À la fin de cette leçon, je sais déterminer la primitive qui prend une valeur donnée en un point (condition initiale)
\item À la fin de cette leçon, je sais modéliser un problème concret (vitesse/distance, coût marginal/coût total) par une recherche de primitive
\end{itemize}
\end{objectifs}

\begin{prerequis}
\begin{itemize}
\item Définition du nombre dérivé et de la fonction dérivée (Première)
\item Tableau des dérivées des fonctions usuelles (puissances, 1/x, racine carrée, cos, sin)
\item Règles de dérivation : somme, produit par une constante, composée avec ax+b
\item Lecture et interprétation graphique d'une fonction et de sa dérivée
\item Résolution d'équations simples dans ℝ
\end{itemize}
\end{prerequis}

\newpage

\coursec{Notion de primitive}

\courssub{Une situation concrète : remonter de la vitesse à la distance}

Imagine un chauffeur de taxi-brousse qui assure la liaison Yaoundé--Douala. Son compteur de vitesse fonctionne parfaitement : à chaque instant $t$, il peut lire sa vitesse $v(t)$. Malheureusement, son carnet de bord kilométrique a été perdu, et il ne sait plus quelle distance $d(t)$ il a parcourue depuis le départ. Peut-il reconstituer la distance parcourue à partir de la seule connaissance de sa vitesse ?

De même, un commerçant du marché Mokolo note chaque jour son bénéfice journalier, c'est-à-dire le \emph{taux de variation} de sa caisse. À la fin du mois, il veut retrouver son bénéfice total. Dans les deux cas, le problème est le même : on connaît la \emph{dérivée} d'une grandeur (la vitesse est la dérivée de la distance, le bénéfice journalier est la dérivée du bénéfice cumulé), et on veut remonter à la grandeur elle-même.

Rappelons un fait essentiel du cours de dérivation : si $d(t)$ désigne la distance parcourue à l'instant $t$, alors la vitesse instantanée est
\[ v(t) = d'(t). \]
Notre problème consiste donc à faire le chemin \emph{en sens inverse} : connaissant $v$, trouver une fonction $d$ dont la dérivée est $v$, c'est-à-dire une fonction $d$ telle que $d'(t) = v(t)$.

\begin{experience}[Découverte : remonter de la vitesse à la distance]
Supposons que le chauffeur roule à vitesse variable et que sa vitesse (en km/h) soit modélisée par $v(t) = 2t$ pendant les premières heures du trajet. Cherchons une distance $d(t)$ telle que $d'(t) = 2t$.
\begin{enumerate}
\item On sait que la dérivée de $t^2$ est $2t$. Donc $d(t) = t^2$ convient : en effet, $d'(t) = 2t = v(t)$.
\item Mais $d(t) = t^2 + 50$ convient aussi, car la dérivée d'une constante est nulle : $(t^2 + 50)' = 2t$.
\item De même, $d(t) = t^2 - 7$ convient encore.
\end{enumerate}
Conclusion de l'activité : la fonction $v(t) = 2t$ admet \emph{plusieurs} « fonctions d'origine » possibles, qui diffèrent toutes d'une constante. Pour choisir la bonne, il faudra une information supplémentaire, par exemple la distance déjà parcourue au départ (une \emph{condition initiale}).
\end{experience}

Cette opération inverse de la dérivation porte un nom : c'est la recherche d'une \cle{primitive}. C'est l'objet de cette leçon, et plus précisément de cette première section où nous allons définir rigoureusement cette notion et apprendre à la vérifier.

\courssub{Définition d'une primitive}

L'activité précédente nous a montré le sens de l'opération : au lieu de partir d'une fonction pour calculer sa dérivée, on part d'une fonction $f$ et on cherche une fonction $F$ dont $f$ serait la dérivée. Formalisons.

\begin{definition}[Primitive d'une fonction sur un intervalle]
Soit $f$ une fonction définie sur un intervalle $I$. On appelle \cle{primitive} de $f$ sur $I$ toute fonction $F$ vérifiant les deux conditions suivantes :
\begin{enumerate}
\item $F$ est \textbf{dérivable} sur $I$ ;
\item pour tout réel $x$ de $I$, $F'(x) = f(x)$, c'est-à-dire $F' = f$ sur $I$.
\end{enumerate}
\end{definition}

\begin{aretenir}[L'idée à retenir]
Dériver et « primitiver » sont deux opérations inverses l'une de l'autre :
\begin{itemize}
\item \textbf{Dériver} : on part de $F$ et on calcule $F' = f$ (sens « descendant ») ;
\item \textbf{Primitiver} : on part de $f$ et on cherche $F$ telle que $F' = f$ (sens « remontant »).
\end{itemize}
\end{aretenir}

\begin{popfigure}
\begin{center}
\begin{tikzpicture}[node distance=4.5cm]
\node[draw=popblue, fill=popblueL, rounded corners, very thick, minimum width=2.6cm, minimum height=1.2cm, font=\large\bfseries] (F) at (0,0) {$F$};
\node[draw=poporange, fill=poporangeL, rounded corners, very thick, minimum width=2.6cm, minimum height=1.2cm, font=\large\bfseries] (f) at (5.5,0) {$f$};
\draw[-{Stealth[length=3mm]}, very thick, popdark] (F.north) to[bend left=25] node[above, font=\small\bfseries, text=popdark]{on dérive : $F' = f$} (f.north);
\draw[-{Stealth[length=3mm]}, very thick, popgreen] (f.south) to[bend left=25] node[below, font=\small\bfseries, text=popgreen!60!black]{on primitive : $F$ est une primitive de $f$} (F.south);
\end{tikzpicture}
\end{center}
\legende{Les deux sens de l'opération : dériver fait « descendre » de $F$ vers $f$ ; primitiver fait « remonter » de $f$ vers $F$.}
\end{popfigure}

\begin{attention}[Bien respecter le sens de l'opération]
L'erreur la plus fréquente est de confondre le sens : dire « $F$ est une primitive de $f$ » signifie que \textbf{$F' = f$}, et certainement pas que $f' = F$. Retiens l'ordre alphabétique comme moyen mnémotechnique : dans l'écriture $F' = f$, la fonction $F$ (avec majuscule, la « grande » fonction) est celle que l'on dérive, et $f$ (minuscule) est le résultat de la dérivation. Avant d'écrire quoi que ce soit sur ta copie, demande-toi toujours : \emph{quelle fonction dois-je dériver pour retomber sur l'autre ?}
\end{attention}

\courssub{Vérification pratique : comment prouver qu'une fonction est une primitive}

La définition fournit immédiatement une méthode de vérification, très souvent demandée au Baccalauréat sous la forme « Montrer que $F$ est une primitive de $f$ sur $I$ ».

\begin{methode}[Vérifier que $F$ est une primitive de $f$ sur $I$]
\begin{enumerate}
\item Justifier que $F$ est dérivable sur $I$ (somme, produit, quotient ou composée de fonctions dérivables).
\item Calculer la dérivée $F'(x)$ pour $x$ dans $I$.
\item Simplifier l'expression obtenue et la comparer à $f(x)$.
\item Conclure : si $F'(x) = f(x)$ pour tout $x$ de $I$, alors $F$ est une primitive de $f$ sur $I$.
\end{enumerate}
\end{methode}

\begin{exemplebox}[Premiers exemples fondamentaux]
\textbf{Exemple 1.} Montrons que $F(x) = x^2$ est une primitive de $f(x) = 2x$ sur $\mathbb{R}$.

La fonction $F$ est une fonction polynôme, donc dérivable sur $\mathbb{R}$, et pour tout réel $x$ :
\[ F'(x) = 2x = f(x). \]
Donc $F$ est bien une primitive de $f$ sur $\mathbb{R}$.

\textbf{Exemple 2.} Montrons que $F(x) = \dfrac{1}{x}$ est une primitive de $f(x) = -\dfrac{1}{x^2}$ sur $]0\,;\,+\infty[$.

La fonction $F$ est dérivable sur $]0\,;\,+\infty[$ (fonction inverse, dérivable sur tout intervalle ne contenant pas $0$), et pour tout $x > 0$ :
\[ F'(x) = -\frac{1}{x^2} = f(x). \]
Donc $F$ est une primitive de $f$ sur $]0\,;\,+\infty[$. Remarque au passage : l'intervalle $I$ fait partie intégrante de la définition ; ici, on ne peut pas parler de primitive sur $\mathbb{R}$ tout entier car $f$ et $F$ ne sont pas définies en $0$.
\end{exemplebox}

\begin{exoresolu}[Vérification complète d'une primitive]
Énoncé. On considère les fonctions définies sur $\mathbb{R}$ par $F(x) = \dfrac{x^3}{3} + 5$ et $f(x) = x^2$. Montrer que $F$ est une primitive de $f$ sur $\mathbb{R}$.
\tcblower
Solution détaillée.
\begin{enumerate}
\item \textbf{Dérivabilité.} $F$ est une fonction polynôme (somme de $x \mapsto \dfrac{x^3}{3}$ et de la constante $5$), donc $F$ est dérivable sur $\mathbb{R}$.
\item \textbf{Calcul de la dérivée.} Pour tout réel $x$, en utilisant $(x^3)' = 3x^2$ et le fait que la dérivée d'une constante est nulle :
\begin{align*}
F'(x) &= \frac{1}{3}\times 3x^2 + 0 \\
&= x^2.
\end{align*}
\item \textbf{Comparaison.} Pour tout réel $x$, $F'(x) = x^2 = f(x)$.
\item \textbf{Conclusion.} $F$ est dérivable sur $\mathbb{R}$ et $F' = f$ sur $\mathbb{R}$, donc $F$ est une primitive de $f$ sur $\mathbb{R}$.
\end{enumerate}
Remarque : la constante $5$ a disparu dans la dérivation. C'est exactement ce phénomène qui explique la non-unicité des primitives, que nous allons maintenant observer.
\end{exoresolu}

\courssub{Non-unicité des primitives : « une » primitive, pas « la » primitive}

Reprenons la fonction $f(x) = 2x$ sur $\mathbb{R}$. Nous avons vu que $F(x) = x^2$ est une primitive de $f$. Mais considérons aussi $G(x) = x^2 + 3$ et $H(x) = x^2 - 7$. Ces deux fonctions sont dérivables sur $\mathbb{R}$ et :
\[ G'(x) = 2x = f(x) \qquad \text{et} \qquad H'(x) = 2x = f(x), \]
car la dérivée d'une constante est nulle. Ainsi $G$ et $H$ sont \emph{aussi} des primitives de $f$. Plus généralement, pour toute constante réelle $k$, la fonction $x \mapsto x^2 + k$ est une primitive de $f$.

Géométriquement, toutes ces courbes s'obtiennent l'une de l'autre par translation verticale : elles ont donc, en tout point de même abscisse, des tangentes \emph{parallèles}, de même pente $f(x) = 2x$.

\begin{popfigure}
\begin{center}
\begin{tikzpicture}
\begin{axis}[
  width=0.85\linewidth, height=7cm,
  axis lines=middle,
  xlabel={$x$}, ylabel={$y$},
  xmin=-2.6, xmax=2.6, ymin=-2.5, ymax=6.5,
  xtick={-2,-1,1,2}, ytick={-1,1,2,3,4,5},
  legend style={at={(0.02,0.98)}, anchor=north west, font=\small, draw=popline},
  samples=100
]
\addplot[popcurve, domain=-2.3:2.3] {x^2};
\addlegendentry{$y = x^2$}
\addplot[popcurve, poporange, domain=-2.1:2.1] {x^2+2};
\addlegendentry{$y = x^2 + 2$}
\addplot[popcurve, popgreen!70!black, domain=-2.4:2.4] {x^2-1};
\addlegendentry{$y = x^2 - 1$}
% tangentes en x = 1, pente 2
\addplot[popdark, dashed, thick, domain=0.1:1.9] {2*(x-1)+1};
\addplot[popdark, dashed, thick, domain=0.1:1.9] {2*(x-1)+3};
\addplot[popdark, dashed, thick, domain=0.1:1.9] {2*(x-1)};
\addplot[only marks, mark=*, popink] coordinates {(1,1) (1,3) (1,0)};
\end{axis}
\end{tikzpicture}
\end{center}
\legende{Trois primitives de $f(x) = 2x$ : les courbes de $x^2$, $x^2+2$ et $x^2-1$ sont translatées verticales l'une de l'autre. En $x = 1$, leurs tangentes (en pointillés) sont parallèles, de pente commune $f(1) = 2$.}
\end{popfigure}

\begin{aretenir}[Vocabulaire : « une » primitive]
Une fonction qui admet une primitive en admet en réalité une infinité : si $F$ est une primitive de $f$ sur $I$, alors pour toute constante réelle $k$, la fonction $F + k$ est aussi une primitive de $f$ sur $I$, puisque $(F + k)' = F' + 0 = f$.

C'est pourquoi on dit toujours \textbf{« une primitive »} de $f$, et jamais \textbf{« la primitive »} de $f$. Nous verrons dans la section suivante que, réciproquement, toutes les primitives d'une même fonction sur un intervalle diffèrent d'une constante, et que l'on peut en choisir une unique en imposant une condition du type $F(a) = b$.
\end{aretenir}

\begin{attention}[Pièges de rédaction à éviter]
\begin{itemize}
\item Ne jamais écrire « la primitive de $f$ est... » sans condition supplémentaire : écris « \emph{une} primitive de $f$ est... ».
\item Ne pas oublier de préciser l'\textbf{intervalle} $I$ : une primitive est toujours définie \emph{sur un intervalle}. Par exemple, $x \mapsto \dfrac{1}{x}$ est une primitive de $x \mapsto -\dfrac{1}{x^2}$ sur $]0\,;\,+\infty[$ (ou sur $]-\infty\,;\,0[$), mais pas sur $\mathbb{R}$.
\item Ne pas oublier de vérifier (au moins mentalement) la \textbf{dérivabilité} de $F$ sur $I$ : c'est la première condition de la définition.
\end{itemize}
\end{attention}

\begin{savaistu}[D'où vient le mot « primitive » ?]
Le mot « primitive » vient du latin \emph{primitivus}, qui signifie « qui vient en premier, originel ». Le vocabulaire est bien choisi : chercher une primitive, c'est littéralement \emph{remonter à la fonction d'origine}, celle qui existait « avant » la dérivation. Cette idée de marche arrière est au cœur des mathématiques : soustraire est l'opération inverse d'additionner, diviser celle de multiplier, extraire une racine carrée celle d'élever au carré... et primitiver est l'opération inverse de dériver. Au chapitre suivant, tu découvriras que cette remontée permet de calculer des aires et des distances parcourues : c'est le célèbre lien entre primitives et intégrales, établi par Newton et Leibniz au XVII\textsuperscript{e} siècle, qui a révolutionné la physique... et qui permettrait à notre chauffeur de taxi-brousse de reconstituer son carnet de bord kilométrique !
\end{savaistu}

\begin{exercice}[À toi de jouer]
\begin{enumerate}
\item Montrer que $F(x) = \dfrac{x^4}{4} - 2$ est une primitive de $f(x) = x^3$ sur $\mathbb{R}$.
\item Donner deux autres primitives de $f(x) = x^3$ sur $\mathbb{R}$.
\item Montrer que $G(x) = 2\sqrt{x}$ est une primitive de $g(x) = \dfrac{1}{\sqrt{x}}$ sur $]0\,;\,+\infty[$.
\end{enumerate}
\end{exercice}

\begin{corrige}[Exercice : À toi de jouer]
\begin{enumerate}
\item $F$ est une fonction polynôme, donc dérivable sur $\mathbb{R}$, et pour tout réel $x$ : $F'(x) = \dfrac{4x^3}{4} - 0 = x^3 = f(x)$. Donc $F$ est une primitive de $f$ sur $\mathbb{R}$.
\item Pour toute constante $k$, $x \mapsto \dfrac{x^4}{4} + k$ convient. Par exemple $x \mapsto \dfrac{x^4}{4}$ et $x \mapsto \dfrac{x^4}{4} + 10$ sont deux autres primitives de $f$ sur $\mathbb{R}$.
\item $G$ est dérivable sur $]0\,;\,+\infty[$ comme produit de la fonction racine carrée (dérivable sur $]0\,;\,+\infty[$) par $2$. Pour tout $x > 0$ : $G'(x) = 2 \times \dfrac{1}{2\sqrt{x}} = \dfrac{1}{\sqrt{x}} = g(x)$. Donc $G$ est une primitive de $g$ sur $]0\,;\,+\infty[$.
\end{enumerate}
\end{corrige}

\coursec{Ensemble des primitives d'une fonction}

Dans la section précédente, nous avons découvert la notion de \cle{primitive} : une fonction $F$ est une primitive de $f$ sur un intervalle $I$ lorsque $F' = f$ sur $I$. Nous avons aussi remarqué, presque en passant, un phénomène troublant : la fonction $x \longmapsto x^2$ est une primitive de $x \longmapsto 2x$, mais $x \longmapsto x^2 + 5$ aussi, et $x \longmapsto x^2 - 7$ également... Autrement dit, \textbf{une fonction qui admet une primitive en admet une infinité}. Cette section répond précisément à deux questions : d'où viennent toutes ces primitives ? Et surtout : une fois qu'on en connaît une, les a-t-on \emph{toutes} ? La réponse, élégante et fondamentale, est l'un des piliers du calcul intégral que tu rencontreras bientôt.

\courssub{Une première observation : décaler verticalement ne change pas la dérivée}

Commençons par une intuition géométrique. La dérivée d'une fonction en un point mesure la \emph{pente de la tangente} à sa courbe. Or, si l'on translate une courbe verticalement (vers le haut ou vers le bas), on ne la déforme pas : en chaque point, la tangente se déplace avec la courbe, mais \textbf{sa pente reste exactement la même}. On s'attend donc à ce qu'ajouter une constante à une fonction ne change pas sa dérivée. C'est exactement ce que confirme le calcul.

\begin{propriete}[Stabilité par ajout d'une constante]
Soit $f$ une fonction définie sur un intervalle $I$ et $F$ une primitive de $f$ sur $I$. Alors, pour tout réel $k$, la fonction $F + k$ définie par $(F+k)(x) = F(x) + k$ est également une primitive de $f$ sur $I$.
\end{propriete}

\begin{experience}[Démonstration de la propriété]
Montrons que $F + k$ est une primitive de $f$. Il suffit de vérifier que sa dérivée vaut $f$.
\begin{enumerate}
\item La fonction constante $x \longmapsto k$ est dérivable sur $I$ et sa dérivée est nulle : $k' = 0$.
\item $F$ est dérivable sur $I$ (c'est une primitive de $f$), donc la somme $F + k$ est dérivable sur $I$ et :
\[ (F + k)'(x) = F'(x) + k' = F'(x) + 0 = f(x). \]
\item Ainsi $(F+k)' = f$ sur $I$ : la fonction $F + k$ est bien une primitive de $f$ sur $I$. \hfill $\blacksquare$
\end{enumerate}
\end{experience}

\begin{exemplebox}[Les primitives de $f(x) = 3x^2$ foisonnent]
Posons $f(x) = 3x^2$ sur $\mathbb{R}$. La fonction $F(x) = x^3$ est une primitive de $f$, car $F'(x) = 3x^2 = f(x)$. D'après la propriété précédente, chacune des fonctions suivantes est aussi une primitive de $f$ :
\[ x^3 + 1, \qquad x^3 - 4, \qquad x^3 + \sqrt{2}, \qquad x^3 - \num{2024}. \]
Vérifions-en une : si $G(x) = x^3 - 4$, alors $G'(x) = 3x^2 - 0 = 3x^2 = f(x)$. On peut fabriquer ainsi \textbf{une infinité} de primitives, une pour chaque valeur de $k$.
\end{exemplebox}

\courssub{Le théorème fondamental : toutes les primitives sont connues dès qu'on en connaît une}

La propriété précédente nous donne \emph{beaucoup} de primitives : toutes les fonctions $F + k$. Mais une question plus subtile demeure : \textbf{existe-t-il d'autres primitives}, qui ne seraient pas de cette forme ? Par exemple, $x \longmapsto 2x$ admettrait-elle une primitive « exotique » qui ne s'écrirait pas $x^2 + k$ ? Le théorème suivant, l'un des plus importants du chapitre, répond par la négative : \emph{non, on a tout}.

\begin{propriete}[Théorème fondamental — ensemble des primitives]
Soit $f$ une fonction admettant une primitive $F$ sur un \textbf{intervalle} $I$. Alors les primitives de $f$ sur $I$ sont \textbf{exactement} les fonctions de la forme
\[ F + k, \quad k \in \mathbb{R}. \]
Autrement dit, deux primitives quelconques de $f$ sur $I$ diffèrent d'une constante, et l'ensemble des primitives de $f$ sur $I$ est :
\[ \{\, F + k \; ; \; k \in \mathbb{R} \,\}. \]
\end{propriete}

\begin{experience}[Démonstration guidée du théorème fondamental]
Nous devons prouver deux choses : (a) toute fonction $F + k$ est une primitive de $f$ ; (b) toute primitive de $f$ est de la forme $F + k$.

\textbf{Étape 1 : le sens facile.} Le point (a) est exactement la propriété démontrée plus haut : pour tout réel $k$, $(F+k)' = F' = f$, donc $F + k$ est une primitive de $f$.

\textbf{Étape 2 : le sens difficile.} Soit $G$ une primitive \emph{quelconque} de $f$ sur $I$. Montrons que $G$ s'écrit forcément $F + k$.
\begin{enumerate}
\item Considérons la fonction différence $H = G - F$, définie sur $I$.
\item $F$ et $G$ sont dérivables sur $I$, donc $H$ est dérivable sur $I$ et :
\[ H'(x) = (G - F)'(x) = G'(x) - F'(x) = f(x) - f(x) = 0. \]
\item Ainsi $H$ est une fonction de dérivée nulle \textbf{sur un intervalle}. Nous admettons le résultat suivant (conséquence du théorème des accroissements finis, vu en classe de Première) : \emph{une fonction dont la dérivée est nulle sur un intervalle est constante sur cet intervalle}.
\item Il existe donc un réel $k$ tel que $H(x) = k$ pour tout $x \in I$, c'est-à-dire $G(x) - F(x) = k$, soit :
\[ G(x) = F(x) + k. \]
\end{enumerate}
Toute primitive $G$ de $f$ est donc bien de la forme $F + k$. Les étapes 1 et 2 réunies prouvent le théorème. \hfill $\blacksquare$
\end{experience}

\begin{popfigure}
\begin{center}
\begin{tikzpicture}[node distance=1.6cm and 2.2cm, every node/.style={font=\small}]
\node[draw=popblue, fill=popblueL, rounded corners, thick, align=center, inner sep=6pt] (hyp) {$F$ et $G$ primitives\\ de $f$ sur $I$};
\node[draw=poporange, fill=poporangeL, rounded corners, thick, align=center, inner sep=6pt, right=of hyp] (der) {$(G-F)' = G' - F'$\\ $= f - f = 0$ sur $I$};
\node[draw=popgreen!60!black, fill=popgreenL, rounded corners, thick, align=center, inner sep=6pt, right=of der] (cst) {$G - F$ constante\\ (dérivée nulle sur\\ un intervalle)};
\node[draw=popink, fill=poppinkL, rounded corners, thick, align=center, inner sep=6pt, right=of cst] (ccl) {$G = F + k$,\\ $k \in \mathbb{R}$};
\draw[-{Stealth[length=3mm]}, very thick, popdark] (hyp) -- (der);
\draw[-{Stealth[length=3mm]}, very thick, popdark] (der) -- (cst);
\draw[-{Stealth[length=3mm]}, very thick, popdark] (cst) -- (ccl);
\end{tikzpicture}
\end{center}
\legende{Schéma de la démonstration : la différence de deux primitives a une dérivée nulle, donc elle est constante.}
\end{popfigure}

\begin{aretenir}[Le réflexe à acquérir]
Dès que tu as trouvé \textbf{une} primitive $F$ de $f$ sur un intervalle $I$, tu connais \textbf{toutes} les primitives de $f$ sur $I$ : ce sont les fonctions $F + k$, $k \in \mathbb{R}$. On écrit souvent : « les primitives de $f$ sont les $F + k$, $k \in \mathbb{R}$ ».
\end{aretenir}

\begin{exoresolu}[Toutes les primitives de $f(x) = 3x^2$]
Énoncé. Déterminer l'ensemble des primitives de la fonction $f$ définie sur $\mathbb{R}$ par $f(x) = 3x^2$.
\tcblower
\textbf{Solution détaillée.}
\begin{enumerate}
\item \emph{On cherche une primitive.} Par lecture inverse du tableau des dérivées, on sait que la dérivée de $x^3$ est $3x^2$. Donc $F(x) = x^3$ est une primitive de $f$ sur $\mathbb{R}$.
\item \emph{On applique le théorème fondamental.} $\mathbb{R}$ est un intervalle, donc toutes les primitives de $f$ sur $\mathbb{R}$ sont les fonctions :
\[ F_k(x) = x^3 + k, \quad k \in \mathbb{R}. \]
\item \emph{Vérification.} Pour tout réel $k$, $F_k'(x) = 3x^2 + 0 = f(x)$ : chaque fonction proposée convient, et le théorème garantit qu'il n'y en a pas d'autre.
\end{enumerate}
L'ensemble des primitives de $f$ est donc $\{\, x \longmapsto x^3 + k \; ; \; k \in \mathbb{R} \,\}$.
\end{exoresolu}

\courssub{Interprétation graphique : une famille de courbes par translation verticale}

Le théorème fondamental possède une traduction géométrique très parlante. Si $F$ est une primitive de $f$, la courbe de $F + k$ s'obtient en translatant celle de $F$ verticalement de $k$ unités (vers le haut si $k > 0$, vers le bas si $k < 0$). Les courbes de toutes les primitives de $f$ forment donc une \cle{famille de courbes} toutes déduites de l'une d'entre elles par des translations verticales : elles sont « empilées » les unes au-dessus des autres, et en tout point d'abscisse $x$, leurs tangentes sont \textbf{parallèles} (même coefficient directeur $f(x)$).

\begin{popfigure}
\begin{center}
\begin{tikzpicture}
\begin{axis}[
  width=0.9\linewidth, height=7cm,
  axis lines=middle,
  xlabel={$x$}, ylabel={$y$},
  xmin=-2.6, xmax=2.6, ymin=-3.5, ymax=6.5,
  xtick={-2,-1,1,2}, ytick={-2,2,4},
  grid=both, grid style={popline!40, very thin},
  legend style={at={(0.02,0.98)}, anchor=north west, font=\small, draw=popline}
]
\addplot[popcurve, popblue, domain=-2.4:2.4, samples=100]{x^2 - 2};
\addlegendentry{$y = x^2 - 2$}
\addplot[popcurve, poporange, domain=-2.4:2.4, samples=100]{x^2};
\addlegendentry{$y = x^2$}
\addplot[popcurve, popgreen!60!black, domain=-2.2:2.2, samples=100]{x^2 + 2};
\addlegendentry{$y = x^2 + 2$}
\draw[-{Stealth[length=2.5mm]}, very thick, popink] (axis cs:1.6,2.56) -- (axis cs:1.6,4.56) node[midway, right, font=\small]{ $+2$};
\draw[-{Stealth[length=2.5mm]}, very thick, popink] (axis cs:1.6,0.56) -- (axis cs:1.6,2.56) node[midway, right, font=\small]{ $+2$};
\end{axis}
\end{tikzpicture}
\end{center}
\legende{Trois primitives de $f(x) = 2x$ : les paraboles $y = x^2 + k$ pour $k \in \{-2\,;\,0\,;\,2\}$ se déduisent l'une de l'autre par translation verticale. En chaque abscisse, les tangentes sont parallèles.}
\end{popfigure}

\courssub{Une condition essentielle : le domaine doit être un intervalle}

Dans l'énoncé du théorème fondamental, l'hypothèse « sur un \textbf{intervalle} » n'est pas une précaution de langage : elle est absolument indispensable. C'est l'un des pièges classiques du Baccalauréat.

\begin{attention}[Le théorème est faux hors d'un intervalle]
Considérons la fonction $f(x) = \dfrac{1}{x}$, définie sur $\mathbb{R}^* = \left]-\infty\,;\,0\right[ \cup \left]0\,;\,+\infty\right[$, qui \textbf{n'est pas un intervalle} (c'est une réunion de deux intervalles disjoints). Les fonctions
\[ G(x) = \begin{cases} \ln(x) + 3 & \text{si } x > 0,\\ \ln(-x) - 7 & \text{si } x < 0, \end{cases} \qquad \text{et} \qquad H(x) = \ln|x| \]
sont toutes deux des primitives de $f$ sur $\mathbb{R}^*$ (vérifie que $G'(x) = H'(x) = \dfrac{1}{x}$ sur chaque morceau). Pourtant leur différence vaut $3$ sur $\left]0\,;\,+\infty\right[$ et $-7$ sur $\left]-\infty\,;\,0\right[$ : elle n'est \textbf{pas constante}. Sur un domaine « en deux morceaux », on peut choisir une constante différente sur chaque morceau ! Retiens donc : le théorème fondamental s'applique uniquement sur un intervalle, et c'est pourquoi la définition d'une primitive elle-même exige un intervalle.
\end{attention}

\courssub{Existence des primitives : un résultat admis (pour l'instant)}

Jusqu'ici, nous avons raisonné ainsi : « \emph{si} $f$ admet une primitive $F$, alors... ». Mais quelles fonctions admettent effectivement des primitives ? Le résultat suivant, que nous admettons, lève toute inquiétude pour les fonctions rencontrées au lycée.

\begin{propriete}[Existence de primitives — admise]
Toute fonction \textbf{continue} sur un intervalle $I$ admet des primitives sur $I$.
\end{propriete}

Ce théorème sera pleinement justifié au chapitre sur l'\cle{intégration} : on y montrera que si $f$ est continue sur $I$ et $a \in I$, alors la fonction « aire sous la courbe » définie par $F(x) = \displaystyle\int_a^x f(t)\,\mathrm{d}t$ est dérivable sur $I$ et vérifie $F' = f$ : c'est le \emph{théorème fondamental de l'analyse}. En attendant, retiens que toutes les fonctions usuelles (polynômes, quotients, racines, fonctions trigonométriques, exponentielle, logarithme) sont continues sur leur ensemble de définition, donc y admettent des primitives.

\begin{attention}[Deux confusions à bannir de tes copies]
\begin{itemize}
\item Ne confonds pas « $f$ admet des primitives » et « $f$ est dérivable » : ce sont deux notions distinctes. La continuité suffit pour garantir l'existence de primitives, sans que $f$ soit dérivable (par exemple $x \longmapsto |x|$ est continue sur $\mathbb{R}$, donc admet des primitives, sans être dérivable en $0$).
\item N'écris jamais « \textbf{la} primitive de $f$ » sans précaution : il y en a une infinité. Écris « \textbf{une} primitive de $f$ est... » ou « les primitives de $f$ sont les $F + k$, $k \in \mathbb{R}$ ». La prochaine section montrera comment une condition du type $F(a) = b$ permet de sélectionner \emph{une unique} primitive.
\end{itemize}
\end{attention}

\begin{savaistu}[La fameuse « constante d'intégration »]
Le réel $k$ qui apparaît dans l'ensemble des primitives porte un nom : on l'appelle la \cle{constante d'intégration}. Loin d'être un détail gênant, elle porte une information physique précieuse. En mécanique, si $v(t)$ est la vitesse d'un taxi-brousse sur la route Yaoundé--Douala, alors une primitive de $v$ donne sa position $x(t)$. La constante $k$ correspond alors à la \textbf{position initiale} du véhicule : deux passagers partis de gares différentes avec la même vitesse à chaque instant ont des positions qui diffèrent d'une constante — exactement comme deux primitives ! C'est pour déterminer cette constante qu'on se donne une « condition initiale », objet de la fin de cette leçon. C'est aussi cette idée qui, au XVIII\textsuperscript{e} siècle, a permis à Newton et Leibniz de fonder le calcul intégral, outil indispensable aujourd'hui des ingénieurs de CAMTEL comme des météorologues.
\end{savaistu}

\begin{exercice}[Pour vérifier que tu as compris]
\begin{enumerate}
\item Montrer que $F(x) = x^3 - 2x$ et $G(x) = x^3 - 2x + 9$ sont deux primitives d'une même fonction $f$ sur $\mathbb{R}$, que l'on précisera.
\item Déterminer l'ensemble des primitives de $g(x) = 2x - 5$ sur $\mathbb{R}$.
\item Les fonctions $A(x) = \dfrac{x}{x+1}$ et $B(x) = \dfrac{-1}{x+1}$ sont-elles deux primitives d'une même fonction sur $\left]-1\,;\,+\infty\right[$ ? Justifier en calculant $A(x) - B(x)$.
\end{enumerate}
\end{exercice}

\begin{corrige}[Exercice]
\begin{enumerate}
\item $F'(x) = 3x^2 - 2$ et $G'(x) = 3x^2 - 2$ : $F$ et $G$ sont deux primitives de $f(x) = 3x^2 - 2$ sur $\mathbb{R}$. On remarque que $G - F = 9$, constante, conformément au théorème fondamental.
\item Une primitive de $2x$ est $x^2$, une primitive de $-5$ est $-5x$. Donc $G_0(x) = x^2 - 5x$ est une primitive de $g$ sur l'intervalle $\mathbb{R}$. L'ensemble des primitives de $g$ est $\{\, x \longmapsto x^2 - 5x + k \; ; \; k \in \mathbb{R} \,\}$.
\item Sur l'intervalle $\left]-1\,;\,+\infty\right[$, on calcule :
\[ A(x) - B(x) = \frac{x}{x+1} - \frac{-1}{x+1} = \frac{x+1}{x+1} = 1. \]
La différence $A - B$ est constante (égale à $1$) sur cet intervalle. De plus $A'(x) = \dfrac{1}{(x+1)^2}$ et $B'(x) = \dfrac{1}{(x+1)^2}$ : $A$ et $B$ sont bien deux primitives de $x \longmapsto \dfrac{1}{(x+1)^2}$ sur $\left]-1\,;\,+\infty\right[$, différant de la constante $1$.
\end{enumerate}
\end{corrige}

En résumé : une fonction continue sur un intervalle admet une infinité de primitives, mais cette infinité est parfaitement organisée — il suffit d'en connaître une seule, $F$, pour toutes les décrire : ce sont les $F + k$, $k \in \mathbb{R}$. Reste alors une question naturelle : comment choisir, parmi toutes ces primitives, celle qui répond à un problème concret ? C'est l'objet de la section « Primitive vérifiant une condition initiale ».

\coursec{Primitives des fonctions usuelles}

Dans la section précédente, nous avons établi un résultat fondamental : si une fonction $f$ admet une primitive $F$ sur un intervalle $I$, alors elle en admet une infinité, toutes de la forme $F + k$ où $k$ est une constante réelle. Autrement dit, il suffit de connaître \emph{une seule} primitive pour les connaître toutes. Reste donc la question pratique : \textbf{comment trouver concrètement une primitive d'une fonction donnée ?}

La réponse est d'une simplicité déconcertante : puisque dériver et primitiver sont deux opérations inverses l'une de l'autre, il suffit de lire le tableau des dérivées... \textbf{de droite à gauche} ! C'est toute l'ambition de cette section : construire, justifier et apprendre à utiliser le tableau des primitives usuelles.

\courssub{Le principe : la lecture inverse du tableau des dérivées}

\textbf{Intuition : remonter le courant}\par

Quand tu dérives, tu pars d'une fonction $F$ et tu calcules $F'$. Quand tu cherches une primitive, tu fais le chemin inverse : on te donne $f$ et tu cherches $F$ telle que $F' = f$. C'est comme remonter une rivière à contre-courant : le paysage est le même, mais le sens de lecture change.

Prenons un exemple simple. Tu sais que $(x^2)' = 2x$. Lis cette égalité à l'envers : la fonction $2x$ admet pour primitive la fonction $x^2$. Voilà, tu viens de trouver ta première primitive par lecture inverse !

\begin{popfigure}
\begin{center}
\begin{tikzpicture}[node distance=0.4cm, every node/.style={align=center}]
\node[draw=popblue, fill=popblueL, rounded corners, minimum width=3.2cm, minimum height=1cm, font=\bfseries] (F) at (0,0) {$F(x) = x^2$};
\node[draw=poporange, fill=poporangeL, rounded corners, minimum width=3.2cm, minimum height=1cm, font=\bfseries] (f) at (7,0) {$f(x) = 2x$};
\draw[-{Stealth[length=3mm]}, line width=1.4pt, popblue] (F) to[bend left=25] node[midway, above, font=\small\bfseries, text=popblue]{dérivation} (f);
\draw[-{Stealth[length=3mm]}, line width=1.4pt, poporange] (f) to[bend left=25] node[midway, below, font=\small\bfseries, text=poporange]{primitivation} (F);
\end{tikzpicture}
\end{center}
\legende{Dérivation et primitivation sont deux opérations réciproques : le tableau des dérivées, lu de droite à gauche, donne le tableau des primitives.}
\end{popfigure}

\begin{methode}[Lecture inverse du tableau des dérivées]
Pour trouver une primitive d'une fonction usuelle $f$ sur un intervalle $I$ :
\begin{enumerate}
\item Identifier dans la colonne « dérivée » du tableau des dérivées la fonction $f$ (ou une forme qui s'en rapproche).
\item Lire la ligne à l'envers : la fonction $F$ de la colonne « fonction » est une primitive de $f$.
\item \textbf{Vérifier systématiquement} en redérivant : on doit obtenir $F' = f$.
\item Ne pas oublier la constante : toutes les primitives sont les $F + k$, $k \in \mathbb{R}$.
\end{enumerate}
\end{methode}

\begin{attention}[La vérification n'est pas optionnelle]
La recherche de primitives est un art où l'on « devine » puis l'on \textbf{vérifie}. La seule preuve qu'une fonction $F$ est bien une primitive de $f$ est le calcul de $F'$. Au Baccalauréat, une vérification rédigée en une ligne ($F'(x) = \ldots = f(x)$) sécurise ta réponse et te fait gagner des points de rigueur.
\end{attention}

\courssub{Le tableau des primitives usuelles}

Voici le tableau central de cette leçon. Chaque ligne s'obtient par lecture inverse d'une ligne du tableau des dérivées, et nous la justifierons ensuite une par une.

\begin{propriete}[Primitives des fonctions usuelles]
Pour chaque ligne du tableau ci-dessous, la fonction $F$ est une primitive de la fonction $f$ sur l'intervalle $I$ indiqué. L'ensemble des primitives de $f$ sur $I$ est l'ensemble des fonctions $F + k$, $k \in \mathbb{R}$.
\begin{center}
\begin{tabularx}{\linewidth}{|c|X|X|c|}
\hline
\rowcolor{popblueL} \textbf{Fonction $f(x)$} & \textbf{Une primitive $F(x)$} & \textbf{Intervalle $I$} & \textbf{Condition} \\
\hline
$a$ (constante) & $ax$ & $\mathbb{R}$ & $a \in \mathbb{R}$ \\
\hline
$x^n$ & $\dfrac{x^{n+1}}{n+1}$ & $\mathbb{R}$ si $n \geq 0$ ; $]-\infty;0[$ ou $]0;+\infty[$ si $n \leq -2$ & $n \in \mathbb{Z}$, $n \neq -1$ \\
\hline
$\dfrac{1}{x^2}$ & $-\dfrac{1}{x}$ & $]-\infty;0[$ ou $]0;+\infty[$ & \\
\hline
$\dfrac{1}{\sqrt{x}}$ & $2\sqrt{x}$ & $]0;+\infty[$ & \\
\hline
$\cos x$ & $\sin x$ & $\mathbb{R}$ & \\
\hline
$\sin x$ & $-\cos x$ & $\mathbb{R}$ & \\
\hline
\end{tabularx}
\end{center}
\end{propriete}

\begin{attention}[Le cas $n = -1$ est exclu !]
La formule $x^n \to \dfrac{x^{n+1}}{n+1}$ ne s'applique \textbf{pas} à $f(x) = x^{-1} = \dfrac{1}{x}$, car on diviserait par $n+1 = 0$. La fonction $\dfrac{1}{x}$ possède bien des primitives (la fonction logarithme népérien, que tu rencontreras plus tard), mais elle échappe à ce tableau. Retenir : \cle{aucune puissance de $x$ ne se dérive en $\dfrac{1}{x}$}.
\end{attention}

\courssub{Justification de chaque ligne par dérivation}

Chaque ligne du tableau se \textbf{prouve} en dérivant la primitive proposée. Ces vérifications sont formatrices : elles t'entraînent à la fois à dériver et à comprendre pourquoi les coefficients apparaissent.

\begin{experience}[Vérification systématique du tableau]
\textbf{Ligne 1 : $f(x) = a$.} Dérivons $F(x) = ax$ sur $\mathbb{R}$ : $F'(x) = a = f(x)$. ✓ En particulier, une primitive de la fonction constante $1$ est $x$.

\textbf{Ligne 2 : $f(x) = x^n$, $n \in \mathbb{Z}$, $n \neq -1$.} Dérivons $F(x) = \dfrac{x^{n+1}}{n+1}$. Comme $n+1 \neq 0$, on peut dériver la puissance :
\[ F'(x) = \frac{1}{n+1} \times (n+1)\, x^{n} = x^n = f(x). \]
✓ C'est ici qu'on comprend \emph{d'où vient} le coefficient $\dfrac{1}{n+1}$ : il compense exactement le facteur $n+1$ qui « descend » lors de la dérivation de $x^{n+1}$.

\textbf{Ligne 3 : $f(x) = \dfrac{1}{x^2}$.} Dérivons $F(x) = -\dfrac{1}{x}$ sur $]0;+\infty[$ (ou $]-\infty;0[$) :
\[ F'(x) = -\left(-\frac{1}{x^2}\right) = \frac{1}{x^2} = f(x). \]
✓ Remarque que c'est cohérent avec la ligne 2 : $\dfrac{1}{x^2} = x^{-2}$, et la formule donne $\dfrac{x^{-1}}{-1} = -\dfrac{1}{x}$.

\textbf{Ligne 4 : $f(x) = \dfrac{1}{\sqrt{x}}$.} Dérivons $F(x) = 2\sqrt{x}$ sur $]0;+\infty[$ :
\[ F'(x) = 2 \times \frac{1}{2\sqrt{x}} = \frac{1}{\sqrt{x}} = f(x). \]
✓ Là encore, c'est la ligne 2 avec $n = -\tfrac{1}{2}$ : $\dfrac{x^{1/2}}{1/2} = 2\sqrt{x}$.

\textbf{Ligne 5 : $f(x) = \cos x$.} Dérivons $F(x) = \sin x$ : $F'(x) = \cos x = f(x)$. ✓

\textbf{Ligne 6 : $f(x) = \sin x$.} Dérivons $F(x) = -\cos x$ :
\[ F'(x) = -(-\sin x) = \sin x = f(x). \]
✓ Le signe moins est indispensable : c'est lui qui compense le $-\sin x$ issu de la dérivation de $\cos x$.
\end{experience}

\begin{attention}[Le piège du signe avec les fonctions trigonométriques]
Les deux lignes trigonométriques ne sont \textbf{pas} symétriques :
\[ \text{primitive de } \cos x \;=\; \sin x \qquad \text{mais} \qquad \text{primitive de } \sin x \;=\; -\cos x. \]
Beaucoup d'élèves écrivent que la primitive de $\sin x$ est $\cos x$ par analogie avec $(\sin x)' = \cos x$. C'est faux : $(\cos x)' = -\sin x \neq \sin x$. Astuce de mémorisation : \cle{quand on primitive, le signe moins migre du cosinus vers le sinus}. En cas de doute, redérive !
\end{attention}

\courssub{Les cas particuliers fondamentaux : les puissances de $x$}

La formule $x^n \to \dfrac{x^{n+1}}{n+1}$ est de loin la plus utilisée au Baccalauréat. Détaillons ses cas les plus fréquents, pour $n \geq 0$ (valides alors sur $\mathbb{R}$ tout entier) :

\begin{center}
\begin{tabularx}{\linewidth}{|c|X|X|}
\hline
\rowcolor{popblueL} \textbf{Fonction $f(x)$} & \textbf{Une primitive $F(x)$} & \textbf{Vérification $F'(x)$} \\
\hline
$1$ & $x$ & $1$ \\
\hline
$x$ & $\dfrac{x^2}{2}$ & $\dfrac{2x}{2} = x$ \\
\hline
$x^2$ & $\dfrac{x^3}{3}$ & $\dfrac{3x^2}{3} = x^2$ \\
\hline
$x^3$ & $\dfrac{x^4}{4}$ & $\dfrac{4x^3}{4} = x^3$ \\
\hline
$x^n$ ($n \geq 0$) & $\dfrac{x^{n+1}}{n+1}$ & $x^n$ \\
\hline
\end{tabularx}
\end{center}

\begin{methode}[Primitiver une puissance de $x$ : la règle des deux gestes]
Pour trouver une primitive de $x^n$ (avec $n \neq -1$) :
\begin{enumerate}
\item \textbf{Augmenter l'exposant de 1} : $x^n$ devient $x^{n+1}$.
\item \textbf{Diviser par le nouvel exposant} : on obtient $\dfrac{x^{n+1}}{n+1}$.
\end{enumerate}
En résumé : \cle{monter l'exposant, puis diviser par lui}. Puis vérifier en redérivant.
\end{methode}

\begin{attention}[L'erreur n°1 : oublier le coefficient]
L'erreur la plus fréquente au Baccalauréat est d'écrire que la primitive de $x^2$ est $x^3$. C'est faux : $(x^3)' = 3x^2 \neq x^2$. La bonne réponse est $\dfrac{x^3}{3}$. Chaque fois que tu « montes » l'exposant, le facteur qui descendra à la dérivation doit être neutralisé par une division. \textbf{Réflexe de survie : après avoir trouvé une primitive, redérive-la mentalement.}
\end{attention}

\begin{exemplebox}[Primitive de $f(x) = x^4$]
Cherchons les primitives de $f(x) = x^4$ sur $\mathbb{R}$.
\begin{enumerate}
\item On augmente l'exposant de 1 : $x^4 \to x^5$.
\item On divise par le nouvel exposant : $F(x) = \dfrac{x^5}{5}$.
\item Vérification : $F'(x) = \dfrac{5x^4}{5} = x^4 = f(x)$. ✓
\end{enumerate}
L'ensemble des primitives de $f$ sur $\mathbb{R}$ est donc l'ensemble des fonctions
\[ F_k(x) = \frac{x^5}{5} + k, \quad k \in \mathbb{R}. \]
\end{exemplebox}

\courssub{Intervalles de validité : une précision qui compte}

Une primitive n'existe que sur un \textbf{intervalle} où la fonction est définie (et, en Terminale, continue). Il faut donc toujours préciser l'intervalle $I$ :

\begin{itemize}
\item Les fonctions $x \mapsto x^n$ avec $n \geq 0$, $\cos$ et $\sin$ sont définies sur $\mathbb{R}$ : leurs primitives sont valables sur $\mathbb{R}$.
\item La fonction $x \mapsto \dfrac{1}{x^2}$ n'est pas définie en $0$. Sa primitive $-\dfrac{1}{x}$ est valable sur $]-\infty;0[$ \textbf{ou} sur $]0;+\infty[$, mais jamais sur un intervalle contenant $0$. Attention : sur chacun de ces deux intervalles, la constante $k$ peut être différente !
\item La fonction $x \mapsto \dfrac{1}{\sqrt{x}}$ n'est définie que sur $]0;+\infty[$ : sa primitive $2\sqrt{x}$ n'a de sens que là.
\item Pour $x^n$ avec $n \leq -2$ (par exemple $x^{-3} = \dfrac{1}{x^3}$), la fonction n'est pas définie en $0$ : la formule s'applique sur $]-\infty;0[$ ou $]0;+\infty[$.
\end{itemize}

\begin{attention}[Ne jamais « traverser » une valeur interdite]
Écrire « une primitive de $\dfrac{1}{x^2}$ sur $\mathbb{R}$ est $-\dfrac{1}{x}$ » est une faute grave : ni $f$ ni $F$ ne sont définies en $0$. Au Baccalauréat, précise toujours l'intervalle, surtout quand l'énoncé en impose un.
\end{attention}

\courssub{Illustration graphique : la relation pente/valeur}

Visualisons le lien entre une fonction et sa primitive sur l'exemple le plus simple : $f(x) = 2x$ et $F(x) = x^2$. En tout point d'abscisse $a$, la \textbf{pente de la tangente} à la courbe de $F$ vaut exactement la \textbf{valeur} $f(a)$. Par exemple, en $a = 1{,}5$ : la tangente à la parabole a pour pente $f(1{,}5) = 3$.

\begin{popfigure}
\begin{center}
\begin{tikzpicture}
\begin{axis}[
    width=0.95\linewidth, height=7.5cm,
    axis lines=middle,
    xlabel={$x$}, ylabel={$y$},
    xmin=-0.4, xmax=2.6, ymin=-0.6, ymax=5.2,
    xtick={0,0.5,1,1.5,2,2.5}, ytick={1,2,3,4,5},
    grid=both, grid style={popline!40},
    legend style={at={(0.03,0.97)}, anchor=north west, font=\small},
    clip=true
]
\addplot[popcurve, popblue, domain=0:2.2, samples=200]{x^2};
\addlegendentry{$F(x) = x^2$ (primitive)}
\addplot[popcurve, poporange, domain=0:2.4, samples=200]{2*x};
\addlegendentry{$f(x) = 2x$ (fonction)}
\addplot[popcurve, poppurple, dashed, domain=0.4:2.4, samples=50]{3*x - 2.25};
\addlegendentry{tangente à $F$ en $a = 1{,}5$, de pente $f(1{,}5) = 3$}
\addplot[only marks, mark=*, mark size=2pt, popdark] coordinates {(1.5,2.25) (1.5,3)};
\draw[popmuted, dashed] (axis cs:1.5,0) -- (axis cs:1.5,3);
\end{axis}
\end{tikzpicture}
\end{center}
\legende{La valeur de $f$ en $a$ donne la pente de la tangente à la courbe de $F$ au point d'abscisse $a$ : ici $F'(1{,}5) = f(1{,}5) = 3$.}
\end{popfigure}

\begin{savaistu}[Pourquoi les primitives sont-elles si importantes ?]
Les primitives sont la clé du \emph{calcul intégral}, que tu découvriras bientôt : l'aire sous la courbe d'une fonction positive $f$ entre $a$ et $b$ se calcule par $F(b) - F(a)$. Ce résultat, le \textbf{théorème fondamental de l'analyse}, dû à Newton et Leibniz à la fin du XVII\textsuperscript{e} siècle, a révolutionné les sciences : c'est grâce à lui qu'on calcule des distances parcourues à partir de vitesses, des volumes, des probabilités... et même la consommation d'électricité d'un quartier de Yaoundé à partir de la puissance instantanée mesurée par ENEO !
\end{savaistu}

\courssub{Entraînement à la lecture inverse rapide}

\begin{exoresolu}[Lecture inverse du tableau : entraînement chronométré]
Déterminer une primitive de chacune des fonctions suivantes, en précisant l'intervalle, puis vérifier par dérivation :
\[ f_1(x) = x^5 \;;\qquad f_2(x) = \frac{1}{x^2} \;;\qquad f_3(x) = \frac{1}{\sqrt{x}} \;;\qquad f_4(x) = \cos x \;;\qquad f_5(x) = \sin x \;;\qquad f_6(x) = 7. \]
\tcblower
\textbf{Solution détaillée.}

\textbf{1.} $f_1(x) = x^5$ sur $\mathbb{R}$ : on monte l'exposant ($x^6$) et on divise par 6 : $F_1(x) = \dfrac{x^6}{6}$. Vérification : $F_1'(x) = \dfrac{6x^5}{6} = x^5$. ✓

\textbf{2.} $f_2(x) = \dfrac{1}{x^2}$ sur $]0;+\infty[$ (par exemple) : lecture directe du tableau, $F_2(x) = -\dfrac{1}{x}$. Vérification : $F_2'(x) = -\left(-\dfrac{1}{x^2}\right) = \dfrac{1}{x^2}$. ✓

\textbf{3.} $f_3(x) = \dfrac{1}{\sqrt{x}}$ sur $]0;+\infty[$ : $F_3(x) = 2\sqrt{x}$. Vérification : $F_3'(x) = 2 \times \dfrac{1}{2\sqrt{x}} = \dfrac{1}{\sqrt{x}}$. ✓

\textbf{4.} $f_4(x) = \cos x$ sur $\mathbb{R}$ : $F_4(x) = \sin x$. Vérification : $(\sin x)' = \cos x$. ✓

\textbf{5.} $f_5(x) = \sin x$ sur $\mathbb{R}$ : attention au signe, $F_5(x) = -\cos x$. Vérification : $(-\cos x)' = -(-\sin x) = \sin x$. ✓

\textbf{6.} $f_6(x) = 7$ sur $\mathbb{R}$ : c'est une constante, $F_6(x) = 7x$. Vérification : $(7x)' = 7$. ✓
\end{exoresolu}

\begin{exercice}[À toi de jouer]
Déterminer une primitive de chacune des fonctions suivantes sur l'intervalle indiqué, puis vérifier :
\begin{enumerate}
\item $f(x) = x^7$ sur $\mathbb{R}$ ;
\item $g(x) = \dfrac{1}{x^3}$ sur $]0;+\infty[$ (indice : écrire $g(x) = x^{-3}$) ;
\item $h(x) = 0$ sur $\mathbb{R}$ ;
\item $s(x) = \sin x$ sur $\mathbb{R}$, puis donner celle qui s'annule en $0$.
\end{enumerate}
\end{exercice}

\begin{corrige}[Exercice : À toi de jouer]
\textbf{1.} $F(x) = \dfrac{x^8}{8}$. Vérification : $F'(x) = \dfrac{8x^7}{8} = x^7$. ✓

\textbf{2.} $g(x) = x^{-3}$ avec $-3 \neq -1$ : la formule donne $G(x) = \dfrac{x^{-2}}{-2} = -\dfrac{1}{2x^2}$. Vérification : $G'(x) = -\dfrac{1}{2} \times (-2)x^{-3} = x^{-3} = \dfrac{1}{x^3}$. ✓

\textbf{3.} $H(x) = 0 \times x = 0$, ou plus généralement toute constante : les primitives de la fonction nulle sont les fonctions constantes. Vérification : $(k)' = 0$. ✓

\textbf{4.} Les primitives de $s(x) = \sin x$ sont les $S_k(x) = -\cos x + k$. On veut $S_k(0) = 0$, c'est-à-dire $-\cos 0 + k = -1 + k = 0$, donc $k = 1$. La primitive cherchée est $S(x) = 1 - \cos x$. Vérification : $S'(x) = \sin x$ et $S(0) = 1 - 1 = 0$. ✓
\end{corrige}

\begin{aretenir}[L'essentiel de la section]
\begin{itemize}
\item Chercher une primitive, c'est \cle{lire le tableau des dérivées de droite à gauche}.
\item Règle des puissances ($n \in \mathbb{Z}$, $n \neq -1$) : \textbf{monter l'exposant, diviser par le nouvel exposant} :
\[ x^n \;\longrightarrow\; \frac{x^{n+1}}{n+1}. \]
En particulier : $1 \to x$, $\;x \to \dfrac{x^2}{2}$, $\;x^2 \to \dfrac{x^3}{3}$.
\item Lignes à connaître par cœur : $\dfrac{1}{x^2} \to -\dfrac{1}{x}$, $\;\dfrac{1}{\sqrt{x}} \to 2\sqrt{x}$, $\;\cos x \to \sin x$, $\;\sin x \to -\cos x$ (signe moins !).
\item Toujours préciser l'\textbf{intervalle} de validité, et ne jamais oublier la constante $k$ quand on cherche \emph{toutes} les primitives.
\item Réflexe absolu : \textbf{vérifier en redérivant}.
\end{itemize}
\end{aretenir}

\coursec{Opérations sur les primitives}

Dans la section précédente, nous avons établi le tableau des \cle{primitives usuelles}, obtenu par lecture inverse du tableau des dérivées : une primitive de $x^n$ est $\dfrac{x^{n+1}}{n+1}$, une primitive de $\dfrac{1}{x^2}$ est $-\dfrac{1}{x}$, une primitive de $\cos$ est $\sin$, etc. C'est un beau résultat, mais il a une limite évidente : la plupart des fonctions que l'on rencontre au Baccalauréat ne figurent pas dans ce tableau ! Comment primitiver $f(x) = 3x^2 - 4x + 7$ ? Comment primitiver $g(x) = \dfrac{6}{x^2} + 5x$ ?

La bonne nouvelle, c'est qu'il suffit de \emph{décomposer} ces fonctions en morceaux usuels, grâce à deux opérations simples : l'\cle{addition} et la \cle{multiplication par un réel}. C'est exactement ce que nous allons faire dans cette section. L'idée directrice est la suivante : \emph{la dérivation se comporte bien avec les sommes et les produits par une constante ; donc, par lecture inverse, la primitivation aussi}.

\courssub{Linéarité de la primitivation}

\textbf{Intuition : tout repose sur la dérivation}\par

Rappelle-toi deux règles de dérivation apprises en Première :
\[ (u + v)' = u' + v' \qquad \text{et} \qquad (\lambda u)' = \lambda u' \quad (\lambda \in \mathbb{R}). \]
Autrement dit : \og la dérivée d'une somme est la somme des dérivées \fg{} et \og on peut sortir une constante de la dérivation \fg. Or, chercher une primitive, c'est faire le chemin \emph{en sens inverse} : on cherche une fonction dont la dérivée est connue. Il est donc naturel de deviner que ces deux règles se \og retournent \fg{} : la primitive d'une somme devrait être la somme des primitives, et on devrait pouvoir sortir une constante de la primitivation. C'est exactement ce qu'affirme la propriété suivante.

\begin{propriete}[Linéarité des primitives]
Soient $f$ et $g$ deux fonctions admettant des primitives $F$ et $G$ sur un intervalle $I$, et soit $\lambda$ un nombre réel. Alors :
\begin{enumerate}
\item la fonction $F + G$ est une primitive de $f + g$ sur $I$ ;
\item la fonction $\lambda F$ est une primitive de $\lambda f$ sur $I$.
\end{enumerate}
En résumé : \textbf{une primitive de $f+g$ est $F+G$} et \textbf{une primitive de $\lambda f$ est $\lambda F$}.
\end{propriete}

\begin{experience}[Démonstration (simple et formatrice)]
La démonstration est immédiate : il suffit de \emph{dériver} les candidats proposés et de vérifier qu'on retombe sur la bonne fonction. C'est la méthode universelle pour prouver qu'une fonction est une primitive d'une autre.

\textbf{Point 1 : primitive de la somme.} Par hypothèse, $F' = f$ et $G' = g$ sur $I$. La fonction $F + G$ est dérivable sur $I$ comme somme de fonctions dérivables, et :
\[ (F + G)' = F' + G' = f + g. \]
Donc $F + G$ a pour dérivée $f + g$ : c'est bien une primitive de $f + g$ sur $I$.

\textbf{Point 2 : primitive de $\lambda f$.} La fonction $\lambda F$ est dérivable sur $I$ et :
\[ (\lambda F)' = \lambda F' = \lambda f. \]
Donc $\lambda F$ est bien une primitive de $\lambda f$ sur $I$. \hfill$\blacksquare$
\end{experience}

\begin{aretenir}[Ce qu'il faut retenir]
Primitiver est une opération \cle{linéaire} : on peut primitiver \emph{terme à terme} et \emph{sortir les coefficients}. Concrètement :
\[ \text{primitive de } \big(\lambda f + \mu g\big) = \lambda \times (\text{primitive de } f) + \mu \times (\text{primitive de } g). \]
C'est cette propriété qui va nous permettre de primitiver tous les polynômes et bien d'autres fonctions.
\end{aretenir}

\courssub{Application aux polynômes : primitiver terme à terme}

Un polynôme est une somme de monômes, et chaque monôme $ax^n$ est le produit d'une constante $a$ par la fonction usuelle $x^n$. La linéarité s'applique donc parfaitement.

\begin{methode}[Primitiver un polynôme]
Pour déterminer une primitive d'un polynôme sur $\mathbb{R}$ :
\begin{enumerate}
\item \textbf{Décomposer} le polynôme en monômes (c'est déjà fait par écriture).
\item \textbf{Primitiver chaque monôme séparément} : pour $ax^n$, utiliser la lecture inverse du tableau des dérivées : une primitive de $x^n$ est $\dfrac{x^{n+1}}{n+1}$, donc une primitive de $ax^n$ est $a\,\dfrac{x^{n+1}}{n+1}$.
\item \textbf{Additionner} les résultats obtenus.
\item \textbf{Ajouter la constante} $k$ (un seul $k$ pour toute la fonction !) si l'on veut \emph{toutes} les primitives.
\item \textbf{Vérifier mentalement} en redérivant le résultat : on doit retrouver le polynôme de départ.
\end{enumerate}
\end{methode}

\begin{exemplebox}[Primitives de $f(x) = 3x^2 - 4x + 7$]
Déterminons toutes les primitives de $f(x) = 3x^2 - 4x + 7$ sur $\mathbb{R}$.

On primitive terme à terme :
\begin{itemize}
\item une primitive de $3x^2$ est $3 \times \dfrac{x^3}{3} = x^3$ ;
\item une primitive de $-4x$ est $-4 \times \dfrac{x^2}{2} = -2x^2$ ;
\item une primitive de $7$ (c'est-à-dire $7x^0$) est $7x$.
\end{itemize}
On additionne et on ajoute la constante :
\[ F(x) = x^3 - 2x^2 + 7x + k, \quad k \in \mathbb{R}. \]
\textbf{Vérification} : $F'(x) = 3x^2 - 4x + 7 = f(x)$. \checkmark
\end{exemplebox}

\begin{attention}[Un seul $k$, et à la fin seulement !]
Deux erreurs reviennent sans cesse dans les copies :
\begin{itemize}
\item Écrire une constante \emph{à chaque terme} : $x^3 + k_1 - 2x^2 + k_2 + 7x + k_3$. C'est inutile et lourd : la somme $k_1 + k_2 + k_3$ est elle-même une constante quelconque, qu'on note simplement $k$.
\item Oublier complètement la constante. Sans le \og $+\,k$ \fg, tu n'as écrit qu'\emph{une} primitive, pas \emph{toutes} les primitives. Au Bac, si la question demande \og déterminer les primitives de $f$ \fg, l'absence de $k$ coûte des points.
\end{itemize}
\end{attention}

\begin{exoresolu}[{Primitives de $f(x) = 6x^2 + \dfrac{2}{x^2}$ sur $]0\,;+\infty[$}]
\textbf{Énoncé.} Déterminer toutes les primitives de la fonction $f$ définie sur $]0\,;+\infty[$ par $f(x) = 6x^2 + \dfrac{2}{x^2}$.
\tcblower
\textbf{Solution détaillée.} La fonction $f$ est la somme de deux termes que l'on sait primitiver grâce au tableau des primitives usuelles.

\textbf{Premier terme.} Une primitive de $x^2$ est $\dfrac{x^3}{3}$, donc par linéarité une primitive de $6x^2$ est :
\[ 6 \times \frac{x^3}{3} = 2x^3. \]

\textbf{Second terme.} Une primitive de $\dfrac{1}{x^2}$ est $-\dfrac{1}{x}$ (vérification : $\left(-\dfrac{1}{x}\right)' = \dfrac{1}{x^2}$). Donc une primitive de $\dfrac{2}{x^2}$ est :
\[ 2 \times \left(-\frac{1}{x}\right) = -\frac{2}{x}. \]

\textbf{Conclusion.} Par linéarité, les primitives de $f$ sur $]0\,;+\infty[$ sont les fonctions :
\[ F(x) = 2x^3 - \frac{2}{x} + k, \quad k \in \mathbb{R}. \]
\textbf{Vérification} : $F'(x) = 6x^2 - 2 \times \left(-\dfrac{1}{x^2}\right) = 6x^2 + \dfrac{2}{x^2} = f(x)$. \checkmark
\end{exoresolu}

\begin{popfigure}
\begin{center}
\begin{tikzpicture}[
  node distance=6mm and 10mm,
  terme/.style={rectangle, rounded corners=2pt, draw=popblue, fill=popblueL, inner sep=4pt, font=\small},
  prim/.style={rectangle, rounded corners=2pt, draw=popgreen, fill=popgreenL, inner sep=4pt, font=\small},
  res/.style={rectangle, rounded corners=2pt, draw=poporange, fill=poporangeL, inner sep=5pt, font=\small\bfseries},
  fleche/.style={-{Stealth[length=2.5mm]}, thick, popdark}
]
\node[res] (f) {$f(x) = 6x^2 + \dfrac{2}{x^2}$};
\node[terme, below left=8mm and 2mm of f] (t1) {terme : $6x^2$};
\node[terme, below right=8mm and 2mm of f] (t2) {terme : $\dfrac{2}{x^2}$};
\node[prim, below=7mm of t1] (p1) {primitive : $6 \times \dfrac{x^3}{3} = 2x^3$};
\node[prim, below=7mm of t2] (p2) {primitive : $2 \times \left(-\dfrac{1}{x}\right) = -\dfrac{2}{x}$};
\node[res, below=20mm of f] (r) {$F(x) = 2x^3 - \dfrac{2}{x} + k$};
\draw[fleche] (f) -- (t1);
\draw[fleche] (f) -- (t2);
\draw[fleche] (t1) -- (p1);
\draw[fleche] (t2) -- (p2);
\draw[fleche] (p1.south) |- ($(r.north)+(0,2mm)$) -- (r.north);
\draw[fleche] (p2.south) |- ($(r.north)+(0,2mm)$) -- (r.north);
\end{tikzpicture}
\end{center}
\legende{Arbre de calcul : on décompose $f$ en termes usuels, on primitive chaque terme par lecture inverse du tableau des dérivées, puis on somme et on ajoute la constante $k$.}
\end{popfigure}

\courssub{Ce que la linéarité ne permet PAS de faire}

La linéarité concerne les \emph{sommes} et les \emph{produits par une constante}. Il est tentant de croire qu'elle s'étend à tous les produits. C'est faux, et c'est l'une des erreurs les plus graves de ce chapitre.

\begin{attention}[La primitive d'un produit n'est PAS le produit des primitives]
Il n'existe \textbf{pas de règle simple} pour primitiver un produit $f \times g$ ou un quotient $\dfrac{f}{g}$. Montrons-le sur un exemple très simple : $f(x) = x \cdot x = x^2$.

\begin{itemize}
\item Une primitive de $x$ est $\dfrac{x^2}{2}$. Le \og produit des primitives \fg{} donnerait :
\[ \frac{x^2}{2} \times \frac{x^2}{2} = \frac{x^4}{4}. \]
\item Or la dérivée de $\dfrac{x^4}{4}$ est $x^3$, et non pas $x^2$ ! La bonne primitive de $x \cdot x = x^2$ est $\dfrac{x^3}{3}$.
\end{itemize}
Conclusion : $\dfrac{x^4}{4} \neq \dfrac{x^3}{3}$, le produit des primitives ne convient pas. De même, le quotient des primitives ne convient pas pour un quotient. Face à un produit ou un quotient, il faut d'abord \emph{transformer l'écriture} (développer, simplifier, décomposer) pour se ramener à des sommes de termes usuels.
\end{attention}

\begin{exemplebox}[Se ramener à une somme avant de primitiver]
Pour primitiver $h(x) = (2x+1)(x-3)$ sur $\mathbb{R}$, on ne primitive surtout pas chaque facteur ! On \textbf{développe} d'abord :
\[ h(x) = 2x^2 - 6x + x - 3 = 2x^2 - 5x - 3, \]
puis on primitive terme à terme :
\[ H(x) = \frac{2x^3}{3} - \frac{5x^2}{2} - 3x + k, \quad k \in \mathbb{R}. \]
Vérification : $H'(x) = 2x^2 - 5x - 3 = h(x)$. \checkmark
\end{exemplebox}

\courssub{Complément utile au Bac : primitive de $f(ax+b)$}

\begin{savaistu}[Hors programme strict, mais redoutablement efficace]
Ce qui suit n'est pas exigible dans le programme officiel de Terminale A, mais il apparaît très souvent dans les sujets de Baccalauréat (notamment avec les fonctions trigonométriques) et te fera gagner un temps précieux. Retiens-le comme un \emph{outil de champion}.
\end{savaistu}

On rencontre souvent des fonctions de la forme $x \mapsto f(ax+b)$, par exemple $\cos(2x)$, $(3x-1)^2$ ou $\dfrac{1}{2x+5}$. Comment les primitiver ?

\begin{propriete}[Primitive d'une fonction composée avec une fonction affine]
Soit $f$ une fonction admettant une primitive $F_0$ sur un intervalle $J$, et soient $a$ et $b$ deux réels avec $a \neq 0$. Alors, sur tout intervalle $I$ tel que $ax+b \in J$ pour tout $x \in I$, la fonction définie par
\[ F(x) = \frac{1}{a}\, F_0(ax+b) \]
est une primitive de $x \mapsto f(ax+b)$ sur $I$.
\end{propriete}

\begin{experience}[Justification par dérivation composée]
Dérivons le candidat $F(x) = \dfrac{1}{a} F_0(ax+b)$. La dérivée d'une composée $x \mapsto u(ax+b)$ est $a \times u'(ax+b)$ (règle de dérivation des fonctions composées, la fonction intérieure $ax+b$ ayant pour dérivée $a$). Donc :
\[ F'(x) = \frac{1}{a} \times a \times F_0'(ax+b) = F_0'(ax+b) = f(ax+b), \]
car $F_0' = f$. Le facteur $\dfrac{1}{a}$ compense exactement le facteur $a$ qui apparaît en dérivant la composition. \hfill$\blacksquare$
\end{experience}

\begin{exemplebox}[Deux applications directes]
\textbf{Exemple 1.} Primitiver $f(x) = \cos(2x)$ sur $\mathbb{R}$.\\
Une primitive de $\cos$ est $\sin$, donc avec $a = 2$ :
\[ F(x) = \frac{1}{2}\sin(2x) + k, \quad k \in \mathbb{R}. \]
Vérification : $F'(x) = \dfrac{1}{2} \times 2\cos(2x) = \cos(2x)$. \checkmark

\textbf{Exemple 2.} Primitiver $g(x) = (3x-1)^2$ sur $\mathbb{R}$.\\
Une primitive de $u^2$ est $\dfrac{u^3}{3}$, donc avec $a = 3$ :
\[ G(x) = \frac{1}{3} \times \frac{(3x-1)^3}{3} + k = \frac{(3x-1)^3}{9} + k, \quad k \in \mathbb{R}. \]
Vérification : $G'(x) = \dfrac{1}{9} \times 3 \times 3(3x-1)^2 = (3x-1)^2$. \checkmark
\end{exemplebox}

\begin{attention}[N'oublie jamais le facteur $\dfrac{1}{a}$]
L'erreur classique consiste à écrire qu'une primitive de $\cos(2x)$ est $\sin(2x)$. C'est faux : la dérivée de $\sin(2x)$ est $2\cos(2x)$, il y a un facteur $2$ en trop. Le réflexe de sécurité est toujours le même : \textbf{redérive ton résultat} pour vérifier que tu retombes exactement sur la fonction de départ. Cette vérification prend dix secondes et peut sauver plusieurs points au Bac.
\end{attention}

\courssub{Illustration graphique : de $f$ vers $F$}

Pour bien visualiser le lien entre une fonction et sa primitive, traçons sur un même graphique la droite $f(x) = 2x+1$ et la parabole $F(x) = x^2 + x$, qui est une primitive de $f$ (en effet, $F'(x) = 2x+1$).

\begin{popfigure}
\begin{center}
\begin{tikzpicture}
\begin{axis}[
  width=0.85\linewidth, height=6.5cm,
  axis lines=middle,
  xlabel={$x$}, ylabel={$y$},
  xmin=-2.6, xmax=2.6, ymin=-1.5, ymax=6.5,
  xtick={-2,-1,1,2}, ytick={1,2,3,4,5,6},
  grid=both, grid style={popline, very thin},
  legend style={at={(0.02,0.98)}, anchor=north west, font=\small, draw=popline}
]
\addplot[popcurve, popblue, domain=-2.3:2.3, samples=100]{2*x+1};
\addlegendentry{$f(x)=2x+1$ (dérivée)}
\addplot[popcurve, poporange, domain=-2.3:2.1, samples=150]{x^2+x};
\addlegendentry{$F(x)=x^2+x$ (primitive)}
\end{axis}
\end{tikzpicture}
\end{center}
\legende{La courbe orange $F$ et sa dérivée bleue $f$. Observe : là où $f$ est négative ($x < -\tfrac12$), $F$ décroît ; là où $f$ est positive ($x > -\tfrac12$), $F$ croît ; et $F$ atteint son minimum exactement là où $f$ s'annule (en $x = -\tfrac12$). Le signe de $f$ donne le sens de variation de $F$ : c'est le cœur du lien entre une fonction et ses primitives.}
\end{popfigure}

\begin{aretenir}[Bilan de la section]
\begin{itemize}
\item \textbf{Linéarité} : une primitive de $f+g$ est $F+G$ ; une primitive de $\lambda f$ est $\lambda F$. On primitive donc \emph{terme à terme}.
\item \textbf{Polynômes} : primitiver chaque monôme $ax^n$ en $a\dfrac{x^{n+1}}{n+1}$, additionner, puis ajouter \emph{une seule} constante $k$.
\item \textbf{Piège majeur} : pas de règle pour un produit ou un quotient — transformer d'abord l'écriture (développer, décomposer).
\item \textbf{Complément Bac} : une primitive de $f(ax+b)$ est $\dfrac{1}{a}F_0(ax+b)$, où $F_0$ est une primitive de $f$.
\item \textbf{Réflexe permanent} : vérifier le résultat en le redérivant.
\end{itemize}
\end{aretenir}

\begin{exercice}[À toi de jouer]
\begin{enumerate}
\item Déterminer toutes les primitives de $f(x) = 4x^3 - 6x + \dfrac{3}{x^2}$ sur $]0\,;+\infty[$.
\item Déterminer toutes les primitives de $g(x) = (x+2)(3x-1)$ sur $\mathbb{R}$ (attention au piège !).
\item (Complément Bac) Déterminer une primitive de $h(x) = \sin(3x) + \dfrac{1}{(2x+1)^2}$ sur $\left]-\dfrac12\,;+\infty\right[$.
\end{enumerate}
\end{exercice}

\begin{corrige}[Exercice : À toi de jouer]
\textbf{1.} On primitive terme à terme : une primitive de $4x^3$ est $x^4$ ; une primitive de $-6x$ est $-3x^2$ ; une primitive de $\dfrac{3}{x^2}$ est $-\dfrac{3}{x}$. Donc :
\[ F(x) = x^4 - 3x^2 - \frac{3}{x} + k, \quad k \in \mathbb{R}. \]
Vérification : $F'(x) = 4x^3 - 6x + \dfrac{3}{x^2} = f(x)$. \checkmark

\textbf{2.} On développe d'abord : $g(x) = 3x^2 - x + 6x - 2 = 3x^2 + 5x - 2$. Puis :
\[ G(x) = x^3 + \frac{5x^2}{2} - 2x + k, \quad k \in \mathbb{R}. \]

\textbf{3.} Une primitive de $\sin$ est $-\cos$, donc une primitive de $\sin(3x)$ est $-\dfrac{1}{3}\cos(3x)$. Une primitive de $\dfrac{1}{u^2}$ est $-\dfrac{1}{u}$, donc une primitive de $\dfrac{1}{(2x+1)^2}$ est $\dfrac{1}{2}\times\left(-\dfrac{1}{2x+1}\right) = -\dfrac{1}{2(2x+1)}$. Finalement :
\[ H(x) = -\frac{1}{3}\cos(3x) - \frac{1}{2(2x+1)} + k, \quad k \in \mathbb{R}. \]
\end{corrige}

\coursec{Primitive vérifiant une condition initiale}

Dans la section précédente, nous avons appris à déterminer les primitives d'une fonction grâce aux opérations : primitive d'une somme, d'un produit par une constante, des formes composées $u'u^n$, $\dfrac{u'}{u^2}$, etc. Mais un point essentiel demeure : une fonction qui admet une primitive en admet en réalité une \textbf{infinité}, toutes de la forme $F_0 + k$ où $k$ est une constante réelle. Comment choisir, parmi cette infinité de candidats, \emph{celle} qui répond à un problème concret ? C'est toute la question de la \cle{condition initiale}, que nous allons maintenant étudier.

\courssub{Le problème : choisir une primitive parmi une infinité}

\textbf{Une situation concrète pour comprendre}\par

Imaginons un chauffeur de taxi-moto (bendskin) à Douala. On connaît sa vitesse à chaque instant : $v(t) = 2t - 3$ (en unités adaptées). La vitesse est la dérivée de la position. Pour retrouver la position $x(t)$, il faut donc trouver une primitive de $v$. Or les primitives de $v$ sont toutes les fonctions
\[ F(t) = t^2 - 3t + k, \quad k \in \mathbb{R}. \]
Chaque valeur de $k$ donne une position différente ! Pourtant, le chauffeur, lui, occupe \textbf{une seule} position à chaque instant. Pour lever l'ambiguïté, il suffit de connaître sa position à un instant donné, par exemple sa position de départ : « à l'instant $t = 1$, il se trouve au point d'abscisse $0$ ». Cette information supplémentaire, appelée \cle{condition initiale}, va permettre de déterminer $k$ de manière unique.

\begin{aretenir}[L'idée fondamentale]
Les primitives d'une fonction $f$ sur un intervalle $I$ forment une \textbf{famille infinie} de fonctions $F_0 + k$ ($k \in \mathbb{R}$). Pour désigner \textbf{une seule} primitive, il faut et il suffit d'imposer une condition du type
\[ F(a) = b, \]
c'est-à-dire exiger que la courbe de la primitive passe par un point donné $(a\,;\,b)$.
\end{aretenir}

\courssub{Le théorème d'existence et d'unicité}

\begin{propriete}[Théorème — Primitive vérifiant une condition initiale]
Soit $f$ une fonction admettant des primitives sur un intervalle $I$, $a$ un point de $I$ et $b$ un réel. Alors il existe une \textbf{unique} primitive $F$ de $f$ sur $I$ telle que
\[ F(a) = b. \]
\end{propriete}

Ce théorème affirme deux choses distinctes, qu'il ne faut pas confondre :
\begin{itemize}
\item l'\textbf{existence} : il y a \emph{au moins} une primitive qui convient ;
\item l'\textbf{unicité} : il n'y en a \emph{qu'une seule}.
\end{itemize}

\begin{experience}[Démonstration guidée du théorème]
Cette démonstration est courte et très formatrice : suis chaque étape attentivement.

\textbf{Hypothèses.} $f$ admet des primitives sur $I$ ; on note $F_0$ l'une d'entre elles. On se donne $a \in I$ et $b \in \mathbb{R}$.

\textbf{Étape 1 : forme générale des primitives.}
D'après le théorème vu sur l'ensemble des primitives, toute primitive $G$ de $f$ sur $I$ s'écrit
\[ G(x) = F_0(x) + k, \quad \text{où } k \in \mathbb{R}. \]

\textbf{Étape 2 : traduction de la condition.}
On cherche $G$ telle que $G(a) = b$. En remplaçant $x$ par $a$ dans l'expression ci-dessus :
\[ G(a) = b \iff F_0(a) + k = b \iff k = b - F_0(a). \]

\textbf{Étape 3 : conclusion.}
\begin{itemize}
\item \emph{Existence} : en posant $k = b - F_0(a)$, la fonction $G(x) = F_0(x) + b - F_0(a)$ est bien une primitive de $f$ (c'est $F_0$ plus une constante), et elle vérifie $G(a) = F_0(a) + b - F_0(a) = b$. Il existe donc bien une primitive convenable.
\item \emph{Unicité} : le calcul de l'étape 2 montre que la condition $G(a) = b$ \textbf{force} $k$ à valoir $b - F_0(a)$. Aucune autre valeur de $k$ ne convient. La primitive est donc unique. $\blacksquare$
\end{itemize}
\end{experience}

\begin{aretenir}[Formule à retenir]
Si $F_0$ est une primitive de $f$ sur $I$, alors l'unique primitive $F$ de $f$ vérifiant $F(a) = b$ est
\[ F(x) = F_0(x) + \bigl(b - F_0(a)\bigr). \]
En pratique, on ne mémorise pas cette formule : on refait le raisonnement en trois étapes (méthode ci-dessous).
\end{aretenir}

\courssub{Méthode : déterminer la primitive vérifiant $F(a) = b$}

\begin{methode}[Les trois étapes incontournables]
Pour déterminer la primitive $F$ de $f$ telle que $F(a) = b$ :
\begin{enumerate}
\item \textbf{Trouver une primitive $F_0$} de $f$ sur $I$ (lecture inverse du tableau des dérivées, opérations sur les primitives) et écrire la forme générale : $F(x) = F_0(x) + k$, $k \in \mathbb{R}$.
\item \textbf{Traduire la condition} : remplacer $x$ par $a$ et égaler à $b$, c'est-à-dire écrire l'équation $F_0(a) + k = b$.
\item \textbf{Résoudre en $k$} : $k = b - F_0(a)$, substituer dans $F$, puis \textbf{vérifier} que $F(a) = b$ et que $F' = f$.
\end{enumerate}
\end{methode}

\begin{popfigure}
\begin{tikzpicture}[node distance=0.55cm and 0.9cm, every node/.style={font=\small}]
\node[draw=popblue, fill=popblueL, rounded corners, text width=0.28\linewidth, align=center] (e1) {\textbf{Étape 1}\\ Primitive générale\\ $F(x) = F_0(x) + k$};
\node[draw=poporange, fill=poporangeL, rounded corners, text width=0.28\linewidth, align=center, right=of e1] (e2) {\textbf{Étape 2}\\ Condition initiale\\ $F_0(a) + k = b$};
\node[draw=popgreen, fill=popgreenL, rounded corners, text width=0.28\linewidth, align=center, right=of e2] (e3) {\textbf{Étape 3}\\ Résolution : $k = b - F_0(a)$\\ puis vérification};
\draw[-{Stealth}, thick, popdark] (e1) -- (e2);
\draw[-{Stealth}, thick, popdark] (e2) -- (e3);
\end{tikzpicture}
\legende{Organigramme de la méthode en trois étapes pour déterminer la primitive vérifiant $F(a) = b$.}
\end{popfigure}

\courssub{Exemples résolus}

\begin{exemplebox}[Exemple fondateur : primitive de $f(x) = 2x - 3$ telle que $F(1) = 0$]
\textbf{Étape 1.} Une primitive de $2x$ est $x^2$, une primitive de $-3$ est $-3x$. Donc les primitives de $f$ sont
\[ F(x) = x^2 - 3x + k, \quad k \in \mathbb{R}. \]
\textbf{Étape 2.} La condition $F(1) = 0$ s'écrit :
\[ 1^2 - 3 \times 1 + k = 0 \iff -2 + k = 0. \]
\textbf{Étape 3.} On obtient $k = 2$. D'où
\[ F(x) = x^2 - 3x + 2. \]
\emph{Vérification} : $F'(x) = 2x - 3 = f(x)$ et $F(1) = 1 - 3 + 2 = 0$. La réponse est correcte.
\end{exemplebox}

\begin{exoresolu}[Primitive de $f(x) = 3x^2 + 1$ qui s'annule en $2$]
Déterminer la primitive $F$ de la fonction $f$ définie sur $\mathbb{R}$ par $f(x) = 3x^2 + 1$ qui s'annule en $2$.
\tcblower
\textbf{Solution détaillée.}

\textbf{Étape 1 : primitive générale.} Une primitive de $3x^2$ est $x^3$ (car $(x^3)' = 3x^2$) et une primitive de $1$ est $x$. Donc
\[ F(x) = x^3 + x + k, \quad k \in \mathbb{R}. \]

\textbf{Étape 2 : traduction de la condition.} « $F$ s'annule en $2$ » signifie $F(2) = 0$ :
\[ 2^3 + 2 + k = 0 \iff 8 + 2 + k = 0 \iff 10 + k = 0. \]

\textbf{Étape 3 : résolution et conclusion.} $k = -10$, donc
\[ F(x) = x^3 + x - 10. \]
\emph{Vérification} : $F'(x) = 3x^2 + 1 = f(x)$ et $F(2) = 8 + 2 - 10 = 0$. $\blacksquare$
\end{exoresolu}

\begin{exoresolu}[Avec une fonction composée]
Déterminer la primitive $G$ de $g(x) = \dfrac{2x}{(x^2+1)^2}$ sur $\mathbb{R}$ telle que $G(0) = 3$.
\tcblower
\textbf{Solution détaillée.}

\textbf{Étape 1.} On reconnaît la forme $\dfrac{u'}{u^2}$ avec $u(x) = x^2 + 1$, donc $u'(x) = 2x$. Or une primitive de $\dfrac{u'}{u^2}$ est $-\dfrac{1}{u}$. Ainsi :
\[ G(x) = -\frac{1}{x^2+1} + k, \quad k \in \mathbb{R}. \]

\textbf{Étape 2.} $G(0) = 3 \iff -\dfrac{1}{0+1} + k = 3 \iff -1 + k = 3$.

\textbf{Étape 3.} $k = 4$, donc
\[ G(x) = 4 - \frac{1}{x^2+1}. \]
\emph{Vérification} : $G(0) = 4 - 1 = 3$ et $G'(x) = \dfrac{2x}{(x^2+1)^2} = g(x)$. $\blacksquare$
\end{exoresolu}

\courssub{Interprétation graphique : une seule courbe passe par le point $(a\,;\,b)$}

Géométriquement, les courbes des fonctions $F_0 + k$ se déduisent les unes des autres par des \textbf{translations verticales} : elles forment une « famille de courbes parallèles ». Par un point donné $(a\,;\,b)$ du plan (avec $a \in I$), il passe \textbf{exactement une} courbe de cette famille. C'est précisément le contenu du théorème d'existence et d'unicité.

\begin{popfigure}
\begin{center}
\begin{tikzpicture}
\begin{axis}[
    width=0.85\linewidth, height=7cm,
    axis lines=middle,
    xlabel={$x$}, ylabel={$y$},
    xmin=-2.2, xmax=3.2, ymin=-4.5, ymax=6.5,
    xtick={-2,-1,1,2,3}, ytick={-4,-2,2,4,6},
    grid=both, grid style={popline!40},
    legend style={at={(0.02,0.98)}, anchor=north west, font=\scriptsize}
]
\addplot[popcurve, domain=-1.6:2.8, samples=100, popblue!50]{x^2 - 3*x - 1};
\addplot[popcurve, domain=-1.8:3, samples=100, popblue!50]{x^2 - 3*x};
\addplot[popcurve, domain=-2:3.1, samples=100, very thick, popink]{x^2 - 3*x + 2};
\addplot[popcurve, domain=-2.1:3.1, samples=100, popblue!50]{x^2 - 3*x + 4};
\addplot[only marks, mark=*, mark size=2.5pt, poporange] coordinates {(1,0)};
\node[poporange, font=\small\bfseries, anchor=south west] at (axis cs:1,0) {$(1\,;\,0)$};
\legend{$k=-1$, $k=0$, $F(x)=x^2-3x+2$, $k=4$}
\end{axis}
\end{tikzpicture}
\end{center}
\legende{Famille de primitives de $f(x) = 2x - 3$. Parmi toutes les courbes $y = x^2 - 3x + k$, une seule passe par le point $(1\,;\,0)$ : celle d'équation $y = x^2 - 3x + 2$.}
\end{popfigure}

\begin{attention}[Erreurs fréquentes à éviter]
\begin{itemize}
\item \textbf{Oublier la constante $k$ avant d'appliquer la condition.} Si tu écris directement $F(x) = x^2 - 3x$ puis que tu vérifies $F(1) = -2 \neq 0$, tu risques de conclure à tort qu'il n'y a pas de solution. La constante $k$ doit apparaître \emph{dès la première ligne} : $F(x) = F_0(x) + k$.
\item \textbf{Appliquer la condition à $f$ au lieu de $F$.} La condition porte sur la \emph{primitive} : c'est $F(a) = b$, et non $f(a) = b$. Confondre les deux est une faute classique au Baccalauréat.
\item \textbf{Conclure sans vérifier.} Après avoir trouvé $k$, prends dix secondes pour vérifier que $F(a) = b$ et que $F' = f$. Cette vérification est appréciée du correcteur et te protège des erreurs de calcul.
\item \textbf{Se tromper de signe.} De $F_0(a) + k = b$, on tire $k = b - F_0(a)$, et non $k = F_0(a) - b$. Isole $k$ proprement.
\end{itemize}
\end{attention}

\courssub{Lien avec les situations concrètes}

La condition initiale n'est pas un artifice mathématique : elle traduit une \textbf{donnée physique ou économique réelle}.

\begin{itemize}
\item \textbf{Position initiale (cinématique).} La vitesse $v(t)$ étant la dérivée de la position $x(t)$, retrouver la position revient à primitiver $v$. La condition $x(t_0) = x_0$ (« le véhicule se trouvait au kilomètre $x_0$ à l'instant $t_0$ ») fixe la constante. C'est exactement la situation du bendskin évoquée en introduction.
\item \textbf{Capital initial (économie).} Si l'on connaît le taux d'accroissement d'un capital placé dans une microfinance à Yaoundé, le capital $C(t)$ est une primitive de ce taux ; la condition $C(0) = C_0$ (mise de départ, par exemple $\num{50000}$ FCFA) détermine entièrement l'évolution du capital.
\item \textbf{Coût fixe (production).} En économie, le \emph{coût marginal} est la dérivée du coût total. Primitiver le coût marginal redonne le coût total à une constante près : cette constante est précisément le \textbf{coût fixe} (loyer de l'atelier, machines), que l'on détermine par la condition $C(0) = $ coût fixe.
\end{itemize}

\begin{savaistu}[Pourquoi « condition initiale » ?]
En physique et en ingénierie, les phénomènes sont souvent décrits par des \emph{équations différentielles} : des relations entre une fonction et ses dérivées (par exemple $y' = 2t - 3$). Résoudre une équation différentielle, c'est primitiver ! Mais la solution n'est unique que si l'on connaît l'\textbf{état initial} du système : position de départ d'un mobile, température initiale d'un four, capital de départ. C'est le principe de déterminisme : le présent (la condition initiale) détermine entièrement le futur (la solution unique). Tu retrouveras cette idée en Terminale C et D avec les équations différentielles, et elle est au cœur de la météorologie, de la balistique et des prévisions de CAMTEL sur le trafic réseau.
\end{savaistu}

\begin{exercice}[À toi de jouer]
\begin{enumerate}
\item Déterminer la primitive $F$ de $f(x) = 4x - 5$ sur $\mathbb{R}$ telle que $F(2) = 3$.
\item Déterminer la primitive $G$ de $g(x) = \dfrac{1}{x^2}$ sur $]0\,;\,+\infty[$ telle que $G(1) = 0$.
\item Une entreprise de transformation de manioc à Bafoussam a un coût marginal $C_m(q) = 6q + 10$ (en milliers de FCFA par tonne). Sachant que le coût fixe est de $30$ milliers de FCFA (c'est-à-dire $C(0) = 30$), déterminer le coût total $C(q)$.
\end{enumerate}
\end{exercice}

\begin{corrige}[Exercices]
\begin{enumerate}
\item Primitive générale : $F(x) = 2x^2 - 5x + k$. Condition : $F(2) = 8 - 10 + k = 3$, donc $k = 5$. D'où $F(x) = 2x^2 - 5x + 5$. Vérification : $F(2) = 8 - 10 + 5 = 3$ et $F'(x) = 4x - 5$. 
\item Une primitive de $\dfrac{1}{x^2}$ est $-\dfrac{1}{x}$. Donc $G(x) = -\dfrac{1}{x} + k$. Condition : $G(1) = -1 + k = 0$, donc $k = 1$. D'où $G(x) = 1 - \dfrac{1}{x}$. Vérification : $G(1) = 0$ et $G'(x) = \dfrac{1}{x^2}$. 
\item $C$ est une primitive de $C_m$ : $C(q) = 3q^2 + 10q + k$. Condition : $C(0) = k = 30$. D'où $C(q) = 3q^2 + 10q + 30$ (en milliers de FCFA). La constante $k$ est bien le coût fixe. 
\end{enumerate}
\end{corrige}

\begin{aretenir}[L'essentiel de la section]
\begin{itemize}
\item Si $f$ admet des primitives sur $I$, alors pour tout $a \in I$ et tout $b \in \mathbb{R}$, il existe une \textbf{unique} primitive $F$ de $f$ telle que $F(a) = b$.
\item Méthode : (1) écrire $F(x) = F_0(x) + k$ ; (2) traduire $F_0(a) + k = b$ ; (3) résoudre $k = b - F_0(a)$ et vérifier.
\item Graphiquement : parmi toutes les courbes $y = F_0(x) + k$, une seule passe par le point $(a\,;\,b)$.
\item Concrètement : la constante $k$ porte l'information initiale du problème (position initiale, capital initial, coût fixe).
\end{itemize}
\end{aretenir}

\coursec{Applications et synthèse}

Dans la section précédente, nous avons appris à déterminer \emph{la} primitive d'une fonction qui prend une valeur donnée en un point : la condition initiale $F(x_0) = y_0$ fixe la constante $k$ et sélectionne une unique fonction parmi toutes les primitives. C'est exactement cet outil qui donne tout son sens concret aux primitives. En effet, dans la vie réelle — physique, économie, transport, hydraulique — on connaît souvent un \cle{taux de variation} (une vitesse, un coût marginal, un débit, une croissance démographique) et l'on cherche à retrouver la \emph{quantité} elle-même (la distance, le coût total, le volume, la population). Primitiver, c'est remonter du taux à la quantité. Cette dernière section te montre comment, puis fait le bilan complet de la leçon.

\courssub{Application physique : de la vitesse à la distance}

\textbf{Intuition.} Tu es dans un taxi-moto (« benskin ») à Douala. Le compteur affiche ta \emph{vitesse} à chaque instant : c'est le taux de variation de ta position. Si tu connais la vitesse $v(t)$ à tout instant $t$, peux-tu retrouver la distance parcourue $d(t)$ ? Oui : la vitesse est la dérivée de la distance, donc la distance est une \emph{primitive} de la vitesse.

\begin{propriete}[Vitesse et distance]
Si un mobile se déplace avec une vitesse $v(t)$, fonction continue du temps $t$ sur un intervalle $I$, alors la distance parcourue $d(t)$ est une primitive de $v$ sur $I$ :
\[ d'(t) = v(t). \]
La condition initiale $d(t_0) = d_0$ (position de départ) détermine alors $d$ de manière unique.
\end{propriete}

\begin{exemplebox}[La moto-taxi de Douala]
Une moto-taxi démarre dans le quartier Deïdo avec une vitesse qui augmente régulièrement :
\[ v(t) = 2t + 10 \quad (\text{en \SI{}{km/h}}, t \text{ en \SI{}{h}}), \qquad d(0) = 0. \]
Cherchons la distance parcourue. On primitive terme à terme :
\[ d(t) = t^2 + 10t + k. \]
La condition $d(0) = 0$ donne $k = 0$, donc $d(t) = t^2 + 10t$. Au bout de \SI{2}{h} :
\[ d(2) = 2^2 + 10 \times 2 = 4 + 20 = \SI{24}{km}. \]
La moto a parcouru \SI{24}{km} : de Deïdo, elle a largement dépassé l'aéroport !
\end{exemplebox}

\begin{popfigure}
\begin{tikzpicture}
\begin{axis}[
  width=0.95\linewidth, height=7cm,
  xlabel={$t$ (\SI{}{h})}, ylabel={},
  xmin=0, xmax=2.2, ymin=0, ymax=30,
  grid=both, grid style={popline!40},
  legend style={at={(0.03,0.97)}, anchor=north west, font=\small, fill=popcream, draw=popmuted},
  axis lines=left, tick label style={font=\small}]
\addplot[popcurve, popblue, thick, domain=0:2, samples=50]{2*x+10};
\addlegendentry{$v(t)=2t+10$ (\SI{}{km/h})}
\addplot[popcurve, poporange, thick, domain=0:2, samples=50]{x^2+10*x};
\addlegendentry{$d(t)=t^2+10t$ (\SI{}{km})}
\addplot[mark=*, popink, mark size=2.5] coordinates {(2,24)};
\node[popink, anchor=south west, font=\small] at (axis cs:2,24) {$d(2)=24$};
\end{axis}
\end{tikzpicture}
\legende{La vitesse $v$ (bleu) est la dérivée de la distance $d$ (orange) : la pente de la courbe de $d$ en $t$ vaut $v(t)$. Le point rose marque $d(2)=\SI{24}{km}$.}
\end{popfigure}

Observe la figure : la courbe de $d$ devient de plus en plus pentue, ce qui traduit visuellement que la vitesse augmente. La pente de la tangente à la courbe de $d$ à l'instant $t$ est exactement $v(t)$ : c'est le lien géométrique entre dérivation et primitivation.

\courssub{Application économique : du coût marginal au coût total}

\textbf{Intuition.} Un commerçant au marché Mokolo produit des sacs d'arachides grillées. Le \cle{coût marginal} $C_m(q)$ est le coût de production d'une unité supplémentaire quand on en produit déjà $q$ : c'est le taux de variation du coût total. En économie mathématique, on pose donc :
\[ C_m(q) = C'(q), \]
autrement dit le coût total $C$ est une primitive du coût marginal $C_m$. La condition initiale est ici le \cle{coût fixe} $C(0) = C_0$ : ce que l'on dépense même sans rien produire (loyer de la boutique, licence, transport, électricité).

\begin{exoresolu}[Coût total d'un atelier à Yaoundé]
\textbf{Énoncé.} Un petit atelier de couture a un coût marginal, en milliers de FCFA, donné par
\[ C_m(q) = 6q + 4, \]
où $q$ est le nombre de pagnes confectionnés. Le coût fixe (loyer, machine) est de \num{20} milliers de FCFA. Déterminer le coût total $C(q)$, puis le coût de production de \num{10} pagnes.
\tcblower
\textbf{Solution.} Le coût total est une primitive du coût marginal :
\[ C(q) = 6 \times \frac{q^2}{2} + 4q + k = 3q^2 + 4q + k. \]
La condition initiale $C(0) = 20$ (coût fixe) donne $k = 20$. Donc :
\[ C(q) = 3q^2 + 4q + 20 \quad (\text{en milliers de FCFA}). \]
Pour $q = 10$ pagnes :
\[ C(10) = 3 \times 100 + 4 \times 10 + 20 = 300 + 40 + 20 = 360. \]
Produire \num{10} pagnes coûte \num{360000} FCFA (soit \num{360} milliers de FCFA). Vérification : $C'(q) = 6q + 4 = C_m(q)$ et $C(0) = 20$. Tout est cohérent.
\end{exoresolu}

\courssub{Méthode générale de modélisation}

Les deux exemples précédents suivent exactement la même démarche. Retiens-la : elle structure la plupart des problèmes concrets du Bac faisant intervenir des primitives.

\begin{methode}[Résoudre un problème concret en 4 temps]
\begin{enumerate}
  \item \textbf{Identifier $f$} : repérer dans l'énoncé le taux de variation (vitesse, coût marginal, débit, taux d'accroissement, croissance). C'est la fonction à primitiver.
  \item \textbf{Primitiver} : trouver toutes les primitives $F(x) + k$ par lecture inverse du tableau des dérivées.
  \item \textbf{Utiliser la condition initiale} : l'énoncé donne toujours une valeur connue (position de départ, coût fixe, stock initial, population initiale) qui permet de calculer $k$.
  \item \textbf{Interpréter le résultat} : répondre à la question posée avec une phrase complète et les bonnes unités (\SI{}{km}, FCFA, \SI{}{L}, individus...).
\end{enumerate}
\end{methode}

\begin{attention}[Unités et condition initiale : les deux contrôles obligatoires]
Dans un problème concret, deux vérifications sont indispensables :
\begin{itemize}
  \item \textbf{Cohérence des unités} : si $v$ est en \SI{}{km/h} et $t$ en \SI{}{h}, alors $d$ est en \SI{}{km}. Si $C_m$ est en milliers de FCFA par unité, alors $C$ est en milliers de FCFA. Une unité incohérente signale presque toujours une erreur de primitivation (exposant mal géré, facteur oublié).
  \item \textbf{Condition initiale} : sans elle, tu ne peux pas donner de réponse chiffrée. Oublier la constante $k$, ou oublier de la calculer, est l'erreur la plus fréquente — et la plus sanctionnée au Bac.
\end{itemize}
\end{attention}

\courssub{Schéma-bilan : dérivation et primitivation, deux opérations réciproques}

Toute la leçon tient en un schéma : dériver et primitiver sont deux opérations \emph{inverses} l'une de l'autre, à la constante $k$ près — et c'est la condition initiale qui lève cette indétermination.

\begin{popfigure}
\begin{tikzpicture}[node distance=4.5cm, >=Stealth]
\node[draw=popblue, fill=popblueL, rounded corners, thick, minimum width=2.8cm, minimum height=1.3cm, font=\bfseries\large] (f) {$f$};
\node[draw=poporange, fill=poporangeL, rounded corners, thick, minimum width=2.8cm, minimum height=1.3cm, font=\bfseries\large, right=of f] (F) {$F$};
\draw[->, very thick, poporange] (f.north) to[bend left=30] node[above, font=\small, align=center]{On primitive\\ $F' = f$} (F.north);
\draw[->, very thick, popblue] (F.south) to[bend left=30] node[below, font=\small, align=center]{On dérive\\ $F' = f$} (f.south);
\node[below=1.1cm of F, draw=poppurple, fill=poppurpleL, rounded corners, font=\small, inner sep=8pt, align=center] (cond) {Condition initiale\\ $F(x_0)=y_0$ : fixe $k$};
\draw[->, dashed, poppurple, thick] (cond) -- (F.south east);
\node[above=0.2cm of f, font=\small, text=popmuted, align=center] {Fonction connue\\(taux de variation)};
\node[above=0.2cm of F, font=\small, text=popmuted, align=center] {Primitive recherchée\\(quantité)};
\end{tikzpicture}
\legende{De $f$ à $F$ par primitivation (infinité de primitives $F+k$), de $F$ à $f$ par dérivation (résultat unique). La condition initiale sélectionne une unique primitive.}
\end{popfigure}

\begin{aretenir}[Les trois erreurs à ne plus jamais commettre]
\begin{itemize}
  \item \textbf{Confondre le sens de l'opération} : primitiver $x^n$ donne $\dfrac{x^{n+1}}{n+1}$ (l'exposant \emph{augmente}), alors que dériver le fait \emph{diminuer}. En cas de doute, dérive ta réponse : tu dois retomber sur $f$.
  \item \textbf{Oublier la constante} : « une primitive » s'écrit toujours avec $+k$ (ou $+C$) tant qu'aucune condition initiale n'est utilisée.
  \item \textbf{Oublier l'intervalle} : une primitive est définie \emph{sur un intervalle}. Par exemple, les primitives de $\dfrac{1}{x}$ sont $\ln|x| + k$ sur $]0\,;+\infty[$ \emph{ou} sur $]-\infty\,;0[$, mais jamais sur $\mathbb{R}$ tout entier (la constante $k$ peut être différente de part et d'autre de $0$).
\end{itemize}
\end{aretenir}

\courssub{Tableau récapitulatif à mémoriser pour le Bac}

Voici le tableau final dérivées $\leftrightarrow$ primitives, conforme au programme officiel de Terminale. Lis-le dans les deux sens : de gauche à droite pour primitiver, de droite à gauche pour vérifier. Les formules composées (ligne 10 à 15) sont des \textbf{lectures directes} de la dérivation chainée : elles évitent les substitutions lourdes et sont très attendues au Bac.

\begin{center}
\small
\begin{tabularx}{\linewidth}{|c|X|X|c|}
\hline
\rowcolor{popblueL} Fonction $f$ & Primitive $F$ (à $k$ près) & Dérivée de $F$ (vérification) & Intervalle usuel \\
\hline
$a$ (constante) & $ax$ & $a$ & $\mathbb{R}$ \\
\hline
$x^n$ ($n \in \mathbb{N}$) & $\dfrac{x^{n+1}}{n+1}$ & $x^n$ & $\mathbb{R}$ \\
\hline
$\dfrac{1}{x^2}$ & $-\dfrac{1}{x}$ & $\dfrac{1}{x^2}$ & $]0\,;+\infty[$ ou $]-\infty\,;0[$ \\
\hline
$\dfrac{1}{x}$ & $\ln|x|$ & $\dfrac{1}{x}$ & $]0\,;+\infty[$ ou $]-\infty\,;0[$ \\
\hline
$\dfrac{1}{\sqrt{x}}$ & $2\sqrt{x}$ & $\dfrac{1}{\sqrt{x}}$ & $]0\,;+\infty[$ \\
\hline
$\cos x$ & $\sin x$ & $\cos x$ & $\mathbb{R}$ \\
\hline
$\sin x$ & $-\cos x$ & $\sin x$ & $\mathbb{R}$ \\
\hline
$e^x$ & $e^x$ & $e^x$ & $\mathbb{R}$ \\
\hline
$\dfrac{1}{1+x^2}$ & $\arctan x$ & $\dfrac{1}{1+x^2}$ & $\mathbb{R}$ \\
\hline
$\dfrac{1}{\sqrt{1-x^2}}$ & $\arcsin x$ & $\dfrac{1}{\sqrt{1-x^2}}$ & $]-1\,;1[$ \\
\hline
\rowcolor{popgreenL} \multicolumn{4}{|c|}{\textbf{Formules composées (très utiles au Bac)}} \\
\hline
$u' \cdot u^n$ ($n \neq -1$) & $\dfrac{u^{n+1}}{n+1}$ & $u' u^n$ & selon $u$ \\
\hline
$\dfrac{u'}{u}$ ($u>0$ ou $u<0$) & $\ln|u|$ & $\dfrac{u'}{u}$ & selon $u$ \\
\hline
$\dfrac{u'}{2\sqrt{u}}$ ($u>0$) & $\sqrt{u}$ & $\dfrac{u'}{2\sqrt{u}}$ & selon $u$ \\
\hline
$u' e^u$ & $e^u$ & $u' e^u$ & selon $u$ \\
\hline
$\dfrac{u'}{1+u^2}$ & $\arctan u$ & $\dfrac{u'}{1+u^2}$ & selon $u$ \\
\hline
$\dfrac{u'}{\sqrt{1-u^2}}$ ($|u|<1$) & $\arcsin u$ & $\dfrac{u'}{\sqrt{1-u^2}}$ & selon $u$ \\
\hline
\rowcolor{popblueL} \multicolumn{4}{|c|}{\textbf{Linéarité}} \\
\hline
$u' + v'$ & $u + v$ & $u'+v'$ & intersection des intervalles \\
\hline
$a\,u'$ ($a \in \mathbb{R}$) & $a\,u$ & $a\,u'$ & selon $u$ \\
\hline
\end{tabularx}
\end{center}

\begin{exercice}[À toi de jouer]
Un robinet remplit un réservoir du village de Bafut avec un débit $q(t) = 3t + 2$ (en \SI{}{L/min}). Le réservoir contient déjà \SI{10}{L} à l'instant $t = 0$.
\begin{enumerate}
  \item Déterminer le volume $V(t)$ d'eau dans le réservoir à l'instant $t$.
  \item Quel volume contient-il au bout de \SI{5}{min} ?
\end{enumerate}
\end{exercice}

\begin{corrige}[Exercice : le réservoir de Bafut]
\begin{enumerate}
  \item Le débit est la dérivée du volume : $V'(t) = q(t) = 3t + 2$. On primitive :
  \[ V(t) = \frac{3}{2}t^2 + 2t + k. \]
  La condition initiale $V(0) = 10$ donne $k = 10$, donc $V(t) = \dfrac{3}{2}t^2 + 2t + 10$ (en \SI{}{L}).
  \item $V(5) = \dfrac{3}{2}\times 25 + 2\times 5 + 10 = 37{,}5 + 10 + 10 = \SI{57,5}{L}$.
\end{enumerate}
Vérification : $V'(t) = 3t + 2 = q(t)$ et $V(0) = 10$. Cohérence des unités : débit en \SI{}{L/min} multiplié par un temps en \SI{}{min} donne bien des \SI{}{L}.
\end{corrige}

\begin{savaistu}[Newton, Leibniz et la naissance du calcul différentiel]
Au XVII\textsuperscript{e} siècle, Isaac Newton (en Angleterre) et Gottfried Wilhelm Leibniz (en Allemagne) construisent indépendamment le calcul différentiel et intégral. Newton voulait précisément résoudre le problème que tu viens d'étudier : connaissant la vitesse d'un corps à chaque instant (une planète, une pomme qui tombe), retrouver sa position. C'est sa célèbre « méthode des fluxions ». Leibniz, lui, invente les notations que tu utilises encore : le signe $\int$ et l'écriture $\mathrm{d}x$. Leur querelle de priorité a divisé l'Europe savante pendant des décennies, mais leur découverte commune est devenue le fondement de toute la physique moderne : mécanique, électricité, météorologie... et même les prévisions de trafic utilisées par les applications de transport à Douala !
\end{savaistu}

\courssub{Ouverture : vers le calcul d'aires}

Tu te demandes peut-être pourquoi le programme insiste tant sur les primitives. La réponse arrive dans le prochain chapitre, « Intégration ». Voici l'idée en une phrase : l'\cle{aire} du domaine situé sous la courbe d'une fonction positive $f$, entre les abscisses $a$ et $b$, se calcule à l'aide d'une primitive $F$ de $f$ par la formule fondamentale :
\[ \text{Aire} = F(b) - F(a). \]
Regarde la figure de la moto-taxi : l'aire sous la droite $v(t) = 2t + 10$ entre $0$ et $2$ vaut $d(2) - d(0) = 24$ — c'est exactement la distance parcourue ! Ce n'est pas un hasard : c'est le théorème fondamental de l'analyse, que tu découvriras bientôt. Tout le travail accompli dans cette leçon — tableau des primitives, constante, condition initiale — sera ton outil principal pour calculer des aires, des volumes et des moyennes. Tu es prêt : l'intégration n'attend plus que toi.

\newpage

\coursec{Activité d'intégration}

\coursec{Primitives d'une fonction}

\courssub{Objectifs de l'activité}

Cette activité d'intégration te place dans une situation économique concrète et authentique : la gestion d'un commerce au marché Mokolo, l'un des plus grands marchés de Yaoundé. À travers cette situation, tu vas mobiliser l'ensemble des savoirs et savoir-faire de la leçon sur les primitives. À l'issue de cette activité, tu seras capable de :

\begin{itemize}
\item reconnaître qu'une fonction donnée est la dérivée d'une grandeur économique (lien entre \cle{bénéfice marginal} et \cle{bénéfice total}) ;
\item déterminer une \cle{primitive} d'une fonction polynôme par \cle{lecture inverse} du tableau des dérivées ;
\item déterminer \textbf{la} primitive vérifiant une \cle{condition initiale} du type $B(0) = -50$ ;
\item exploiter la primitive obtenue pour calculer une valeur concrète et prendre une décision économique (quantité optimale de vente) ;
\item vérifier systématiquement un résultat par \cle{dérivation}, réflexe essentiel du mathématicien et du bachelier.
\end{itemize}

Cette activité se réalise en groupes de trois ou quatre élèves, avec une restitution orale au tableau. Durée conseillée : 45 minutes de travail de groupe, 15 minutes de restitution et de synthèse.

\courssub{Situation}

Monsieur Etoundi est commerçant au marché Mokolo à Yaoundé. Chaque semaine, il achète des sacs d'oignons en gros à Bafia et les revend au détail. Rigoureux dans sa comptabilité, il a remarqué que son bénéfice dépend du nombre de sacs vendus, mais pas de façon simple : plus il vend de sacs, plus il doit casser ses prix pour écouler son stock avant que les oignons ne s'abîment, ce qui réduit le gain supplémentaire apporté par chaque sac vendu.

Un jeune neveu, élève en Terminale A, lui explique que les économistes appellent \emph{bénéfice marginal} la variation du bénéfice provoquée par la vente d'un sac supplémentaire, et que cette grandeur n'est autre que la \cle{dérivée} du bénéfice total. Après plusieurs semaines d'observation, on a modélisé la situation de la façon suivante.

On note $q$ le nombre de sacs d'oignons vendus par semaine ($q$ est un réel positif, exprimé en sacs) et $B(q)$ le bénéfice hebdomadaire de M. Etoundi, exprimé en \textbf{milliers de FCFA}. Le bénéfice marginal, exprimé en milliers de FCFA par sac, est donné par :
\[ B_m(q) = -2q + 30. \]
De plus, M. Etoundi supporte des frais fixes (location de l'étal au marché, transport Yaoundé--Bafia, taxes municipales) : même lorsqu'il ne vend aucun sac, il perd \SI{50000}{FCFA} par semaine. On a donc :
\[ B(0) = -50 \quad \text{(en milliers de FCFA)}. \]

M. Etoundi se pose trois questions pratiques :
\begin{enumerate}
\item Quel sera son bénéfice s'il parvient à vendre 10 sacs dans la semaine ?
\item Combien de sacs doit-il viser pour que son bénéfice soit le plus grand possible ?
\item Son neveu affirme qu'au-delà d'un certain nombre de sacs, vendre davantage \emph{réduit} le bénéfice. Est-ce vrai ?
\end{enumerate}

Ton groupe est missionné pour répondre à ces questions à l'aide des outils mathématiques de la leçon.

\courssub{Consignes}

\begin{enumerate}
\item \textbf{Comprendre le lien dérivée--primitive.} On admet que le bénéfice marginal est la dérivée du bénéfice total : $B_m = B'$. Expliquer, avec vos mots, ce que signifie la phrase : \og chercher $B$, c'est chercher une primitive de $B_m$ \fg.
\item \textbf{Déterminer les primitives de $B_m$.} Par lecture inverse du tableau des dérivées, déterminer toutes les primitives de la fonction $q \longmapsto -2q + 30$ sur $[0\,;\,+\infty[$. On écrira $B(q) = \ldots + k$, où $k$ est une constante réelle.
\item \textbf{Exploiter la condition initiale.} En utilisant le fait que $B(0) = -50$, déterminer la valeur de la constante $k$ et donner l'expression complète de $B(q)$.
\item \textbf{Vérifier.} Dériver l'expression obtenue et contrôler que l'on retrouve bien $B_m(q)$, puis vérifier que $B(0) = -50$. Pourquoi cette double vérification est-elle indispensable ?
\item \textbf{Calculer le bénéfice pour 10 sacs.} Calculer $B(10)$ et interpréter le résultat en FCFA pour M. Etoundi.
\item \textbf{Optimiser.} Étudier le signe de $B'(q) = -2q + 30$ sur $[0\,;\,+\infty[$, dresser le tableau de variations de $B$, puis en déduire le nombre de sacs qui maximise le bénéfice et la valeur de ce bénéfice maximal. Répondre alors à la troisième question de M. Etoundi.
\item \textbf{Restitution.} Présenter vos résultats au tableau en rédigeant chaque étape avec soin, comme vous le feriez à l'épreuve de mathématiques du Baccalauréat.
\end{enumerate}

\courssub{Ressources à mobiliser}

\begin{aretenir}[Boîte à outils de la leçon]
\begin{itemize}
\item \textbf{Définition.} $F$ est une \cle{primitive} de $f$ sur un intervalle $I$ si $F$ est dérivable sur $I$ et si $F'(x) = f(x)$ pour tout $x$ de $I$.
\item \textbf{Lecture inverse du tableau des dérivées.} Une primitive de $x \longmapsto x^n$ ($n \in \mathbb{N}$) est $x \longmapsto \dfrac{x^{n+1}}{n+1}$. En particulier, une primitive de $x \longmapsto x$ est $x \longmapsto \dfrac{x^2}{2}$, et une primitive de $x \longmapsto a$ (constante) est $x \longmapsto ax$.
\item \textbf{Linéarité.} Une primitive d'une somme est la somme des primitives ; une primitive de $\alpha f$ est $\alpha F$.
\item \textbf{Ensemble des primitives.} Deux primitives d'une même fonction sur un intervalle diffèrent d'une \cle{constante} : toutes les primitives s'écrivent $F + k$, $k \in \mathbb{R}$.
\item \textbf{Condition initiale.} Il existe une \textbf{unique} primitive $F$ vérifiant $F(a) = b$ : on détermine $k$ en résolvant $F(a) + k = b$.
\item \textbf{Lien signe de $f'$ et variations de $f$.} Si $f' \geq 0$ sur un intervalle, $f$ y est croissante ; si $f' \leq 0$, $f$ y est décroissante.
\item \textbf{Réflexe de vérification.} Toute primitive trouvée se contrôle en la \cle{dérivant} : on doit retrouver la fonction de départ.
\end{itemize}
\end{aretenir}

\courssub{Démarche et corrigé commenté}

\begin{exoresolu}[Le bénéfice du commerçant de Mokolo — corrigé détaillé]
\textbf{Énoncé.} Le bénéfice marginal de M. Etoundi est $B_m(q) = -2q + 30$ (en milliers de FCFA par sac) et $B(0) = -50$. (1) Justifier que $B$ est une primitive de $B_m$. (2) Déterminer $B(q)$. (3) Calculer $B(10)$. (4) Déterminer le nombre de sacs maximisant le bénéfice.

\tcblower

\textbf{Question 1 — Le lien dérivée--primitive.}

Par définition, le bénéfice marginal est la dérivée du bénéfice total : $B' = B_m$. Dire que $B' = B_m$ signifie exactement, d'après la définition du cours, que $B$ est \textbf{une primitive de} $B_m$ sur $[0\,;\,+\infty[$. Chercher le bénéfice total à partir du bénéfice marginal, c'est donc \og remonter \fg{} de la dérivée à la fonction : c'est le geste inverse de la dérivation.

\medskip

\textbf{Question 2 — Détermination de $B(q)$.}

\emph{Étape a : primitives de $B_m$.} On procède par lecture inverse du tableau des dérivées, terme à terme :
\begin{itemize}
\item une primitive de $q \longmapsto -2q$ est $q \longmapsto -2 \times \dfrac{q^2}{2} = -q^2$ ;
\item une primitive de $q \longmapsto 30$ est $q \longmapsto 30q$.
\end{itemize}
Par linéarité (primitive d'une somme), toutes les primitives de $B_m$ sur $[0\,;\,+\infty[$ sont les fonctions :
\[ q \longmapsto -q^2 + 30q + k, \quad k \in \mathbb{R}. \]

\emph{Étape b : condition initiale.} Le bénéfice cherché est la primitive vérifiant $B(0) = -50$. Or :
\[ B(0) = -0^2 + 30 \times 0 + k = k. \]
La condition $B(0) = -50$ impose donc $k = -50$. D'où l'unique primitive solution :
\[ \boxed{B(q) = -q^2 + 30q - 50.} \]

\emph{Vérification systématique (question 4 des consignes).} On dérive : $B'(q) = -2q + 30 = B_m(q)$ : c'est bien le bénéfice marginal annoncé. De plus $B(0) = -50$ : la condition initiale est respectée. Cette double vérification est indispensable car une erreur de calcul sur un coefficient ou sur la constante est immédiatement détectée par la dérivation ; au Baccalauréat, elle sécurise tous les points de la question.

\medskip

\textbf{Question 3 — Bénéfice pour 10 sacs vendus.}
\[ B(10) = -10^2 + 30 \times 10 - 50 = -100 + 300 - 50 = 150. \]
Le bénéfice est de $150$ milliers de FCFA, c'est-à-dire \SI{150000}{FCFA}. Si M. Etoundi vend 10 sacs dans la semaine, il gagne \SI{150000}{FCFA} nets, après paiement de ses frais fixes.

\medskip

\textbf{Question 4 — Optimisation du bénéfice.}

On étudie le signe de la dérivée $B'(q) = -2q + 30$ sur $[0\,;\,+\infty[$ :
\[ -2q + 30 \geq 0 \iff q \leq 15. \]
Donc $B'(q) \geq 0$ sur $[0\,;\,15]$ (la fonction $B$ y est croissante) et $B'(q) \leq 0$ sur $[15\,;\,+\infty[$ (la fonction $B$ y est décroissante). Tableau de variations :

\begin{center}
\begin{tabular}{|c|ccccc|}
\hline
\rowcolor{popblueL} $q$ & $0$ & & $15$ & & $+\infty$ \\
\hline
Signe de $B'(q)$ & & $+$ & $0$ & $-$ & \\
\hline
Variations de $B$ & $-50$ & $\nearrow$ & $175$ & $\searrow$ & \\
\hline
\end{tabular}
\end{center}

Le bénéfice maximal est atteint pour $q = 15$ sacs et vaut :
\[ B(15) = -15^2 + 30 \times 15 - 50 = -225 + 450 - 50 = 175, \]
soit \SI{175000}{FCFA} par semaine.

\medskip

\textbf{Interprétation économique.} M. Etoundi a intérêt à viser la vente de 15 sacs par semaine : son bénéfice est alors maximal, égal à \SI{175000}{FCFA}. Au-delà de 15 sacs, la dérivée devient négative : chaque sac supplémentaire vendu oblige à baisser tellement les prix que le bénéfice total \emph{diminue}. Le neveu avait donc raison : vendre plus ne signifie pas toujours gagner plus. Par exemple, $B(20) = -400 + 600 - 50 = 150$ : vendre 20 sacs rapporte autant que vendre 10 sacs, avec deux fois plus de travail et de risques.

\begin{popfigure}
\begin{tikzpicture}
\begin{axis}[
    width=\linewidth, height=0.55\linewidth,
    axis lines=middle,
    xlabel={$q$ (sacs)}, ylabel={Valeurs (kFCFA)},
    domain=0:22, samples=200,
    xmin=0, xmax=22, ymin=-60, ymax=200,
    xtick={0,5,10,15,20}, ytick={-50,0,50,100,150,175},
    grid=major, grid style={popline},
    legend pos=north east,
    legend style={fill=popcream, draw=popdark, font=\small, cells={anchor=west}},
    clip=false,
    axis line style={popdark},
    tick style={popdark},
    label style={popdark},
]
\addplot[popcurve, thick, popblue] {-x^2 + 30*x - 50};
\addlegendentry{$B(q) = -q^2 + 30q - 50$}
\addplot[popcurve, thick, poporange, dashed] {-2*x + 30};
\addlegendentry{$B'(q) = -2q + 30$}
% Point max
\addplot[only marks, mark=*, mark size=2.5pt, popgreen] coordinates {(15,175)};
\node[above right=2pt, popgreen, font=\small\bfseries] at (axis cs:15,175) {Max $B(15)=175$};
% Point B(0)
\addplot[only marks, mark=*, mark size=2.5pt, poppink] coordinates {(0,-50)};
\node[below right=2pt, poppink, font=\small] at (axis cs:0,-50) {$B(0)=-50$};
% Droite verticale q=15
\draw[dashed, popmuted, thick] (axis cs:15,0) -- (axis cs:15,175);
\node[below, popmuted, font=\small\bfseries] at (axis cs:15,0) {$q=15$};
% Zéro de B' (q=15) déjà marqué
% Zéros de B (seuils de rentabilité) : q^2 - 30q + 50 = 0 -> q = 15 +/- 5*sqrt(7) approx 1.77 et 28.2. Seul 1.77 dans [0,22].
\pgfmathsetmacro{\qzero}{15 - 5*sqrt(7)}
\addplot[only marks, mark=*, mark size=2pt, poppurple] coordinates {(\qzero,0)};
\node[above left=2pt, poppurple, font=\tiny] at (axis cs:\qzero,0) {Seuil rentabilité};
\end{axis}
\end{tikzpicture}
\legende{Représentation graphique du bénéfice total $B(q)$ (courbe bleue) et du bénéfice marginal $B'(q)$ (courbe orange pointillée). Le maximum de $B$ correspond à l'annulation de $B'$ en $q=15$. La zone hachurée (non tracée ici mais implicite) sous $B'$ entre 0 et 15 représente le gain net par rapport aux frais fixes.}
\end{popfigure}

\end{exoresolu}

\begin{attention}[Pièges fréquents dans ce type de situation]
\begin{itemize}
\item \textbf{Oublier la constante $k$.} Écrire directement $B(q) = -q^2 + 30q$ revient à affirmer que $B(0) = 0$, ce qui contredit les frais fixes. Sans condition initiale, il y a une infinité de primitives ; la donnée $B(0) = -50$ est ce qui permet de choisir la bonne.
\item \textbf{Se tromper d'unité.} Les résultats sont en \emph{milliers} de FCFA : $B(10) = 150$ signifie \SI{150000}{FCFA}, et non \SI{150}{FCFA}. L'interprétation finale doit toujours revenir à l'unité concrète.
\item \textbf{Confondre maximum de $B$ et zéro de $B$.} Le bénéfice est maximal là où $B'$ s'annule en changeant de signe (ici $q = 15$), et non là où $B$ s'annule. Les points où $B(q) = 0$ (seuils de rentabilité) répondent à une autre question.
\item \textbf{Négliger la vérification par dérivation.} C'est le contrôle le plus simple et le plus rentable : dérivez toujours la primitive obtenue.
\item \textbf{Omettre l'intervalle de définition.} Une primitive est toujours définie \emph{sur un intervalle}. Ici $[0\,;\,+\infty[$ car $q$ représente un nombre de sacs. Ne pas l'écrire est une imprécision sanctionnée au Bac.
\end{itemize}
\end{attention}

\courssub{Bilan de l'activité}

Cette activité a permis de mettre en évidence les points essentiels de la leçon :

\begin{itemize}
\item \textbf{La primitive comme geste inverse de la dérivation.} Dans la réalité économique, on observe souvent la \emph{variation} d'une grandeur (bénéfice marginal, vitesse, débit) et l'on cherche la grandeur elle-même : c'est exactement le problème de la recherche de primitives.
\item \textbf{Le rôle décisif de la condition initiale.} Mathématiquement, une fonction admet une infinité de primitives différant d'une constante ; mais dans une situation concrète, une donnée supplémentaire (ici les frais fixes $B(0) = -50$) permet de sélectionner \textbf{l'unique} primitive adaptée au problème. La constante d'intégration n'est pas un détail technique : elle porte toute l'information sur les frais fixes du commerçant.
\item \textbf{L'enchaînement des savoir-faire.} La résolution a mobilisé successivement la lecture inverse du tableau des dérivées, la linéarité, la détermination de la constante, la vérification par dérivation, puis l'étude du signe de la dérivée pour optimiser : les primitives ne sont pas un chapitre isolé, elles prolongent naturellement l'étude des fonctions.
\item \textbf{Les mathématiques comme aide à la décision.} Le modèle a montré qu'au-delà de 15 sacs, vendre davantage diminue le bénéfice : un résultat contre-intuitif que seul le calcul (signe de $B'$) permet d'établir rigoureusement. De Yaoundé à Bafia, un commerçant bien conseillé par un élève de Terminale peut ainsi fixer son objectif de vente hebdomadaire.
\end{itemize}

\begin{savaistu}[Le bénéfice marginal, un outil des économistes]
La notion de \og marginal \fg{} (coût marginal, bénéfice marginal, utilité marginale) est au cœur de la science économique moderne depuis les travaux des économistes du XIX\textsuperscript{e} siècle. Elle repose entièrement sur le concept mathématique de dérivée : le marginal est la dérivée du total. Inversement, retrouver le total à partir du marginal, comme tu viens de le faire, est un problème de primitives, qui deviendra en classe supérieure un problème de \emph{calcul intégral}. Les grandes écoles de commerce et les banques centrales, y compris la BEAC à Yaoundé, utilisent quotidiennement ces outils pour modéliser les coûts et les revenus.
\end{savaistu}

\newpage

\coursec{Méthodes et exercices résolus}

\coursec{Primitives d'une fonction}

\courssub{Notion de primitive}

\begin{definition}[Primitive d'une fonction sur un intervalle]
Soit $I$ un intervalle de $\mathbb{R}$ et $f$ une fonction définie sur $I$. On appelle \textbf{primitive de $f$ sur $I$} toute fonction $F$ dérivable sur $I$ telle que :
\[ \forall x \in I,\quad F'(x) = f(x). \]
\end{definition}

\begin{exemplebox}[Premiers exemples]
\begin{itemize}
\item La fonction $F(x) = x^3$ est une primitive de $f(x) = 3x^2$ sur $\mathbb{R}$, car $F'(x) = 3x^2 = f(x)$.
\item La fonction $G(x) = \ln x$ est une primitive de $g(x) = \dfrac{1}{x}$ sur $]0\,;\,+\infty[$, car $G'(x) = \dfrac{1}{x} = g(x)$.
\item La fonction $H(x) = \cos x$ n'est \textbf{pas} une primitive de $\sin x$ sur $\mathbb{R}$, car $H'(x) = -\sin x \neq \sin x$. En revanche, $-\cos x$ en est une.
\end{itemize}
\end{exemplebox}

\courssub{Ensemble des primitives d'une fonction}

\begin{propriete}[Caractérisation de l'ensemble des primitives]
Soit $f$ une fonction admettant des primitives sur un intervalle $I$.
\begin{enumerate}
\item Si $F$ est une primitive de $f$ sur $I$, alors pour toute constante $C \in \mathbb{R}$, la fonction $F + C$ est aussi une primitive de $f$ sur $I$.
\item Réciproquement, si $F$ et $G$ sont deux primitives de $f$ sur $I$, alors il existe une constante $C \in \mathbb{R}$ telle que $G = F + C$ sur $I$.
\end{enumerate}
En d'autres termes, l'ensemble des primitives de $f$ sur $I$ est une droite affine de direction $\mathbb{R}$ : dès qu'on en connaît une, on les connaît toutes.
\end{propriete}

\begin{aretenir}[Notation]
On note souvent $\displaystyle \int f(x)\,dx$ l'ensemble des primitives de $f$ (primitive indéfinie). Si $F_0$ est une primitive particulière, on écrit :
\[ \int f(x)\,dx = F_0(x) + C,\quad C \in \mathbb{R}. \]
\end{aretenir}

\courssub{Primitives des fonctions usuelles}

\begin{propriete}[Tableau des primitives usuelles (lecture inverse du tableau des dérivées)]
\begin{center}
\begin{tabularx}{\linewidth}{|l|X|l|}
\hline
\rowcolor{popblueL} \textbf{Fonction $f$} & \textbf{Une primitive $F$} & \textbf{Intervalle de validité} \\
\hline
$a$ (constante) & $ax$ & $\mathbb{R}$ \\
\hline
$x^n$ \ ($n\in\mathbb{N}$) & $\dfrac{x^{n+1}}{n+1}$ & $\mathbb{R}$ \\
\hline
$\dfrac{1}{x^2}$ & $-\dfrac{1}{x}$ & $]0\,;\,+\infty[$ ou $]-\infty\,;\,0[$ \\
\hline
$\dfrac{1}{x}$ & $\ln x$ & $]0\,;\,+\infty[$ \\
\hline
$\dfrac{1}{\sqrt{x}}$ & $2\sqrt{x}$ & $]0\,;\,+\infty[$ \\
\hline
$e^x$ & $e^x$ & $\mathbb{R}$ \\
\hline
$\cos x$ & $\sin x$ & $\mathbb{R}$ \\
\hline
$\sin x$ & $-\cos x$ & $\mathbb{R}$ \\
\hline
\end{tabularx}
\end{center}
\end{propriete}

\begin{methode}[Trouver une primitive d'une fonction usuelle]
Pour déterminer une primitive $F$ d'une fonction $f$ sur un intervalle $I$ :
\begin{enumerate}
\item \textbf{Identifier} la forme de $f$ : est-ce une fonction de référence (puissance, $\frac{1}{x}$, $\frac{1}{x^2}$, $\sqrt{x}$, exponentielle, cos, sin\dots) ou une somme de telles fonctions ?
\item \textbf{Se poser la question clé} : \og quelle fonction, une fois dérivée, donne $f$ ? \fg{} C'est la lecture \emph{inverse} du tableau des dérivées.
\item \textbf{Utiliser la linéarité} : une primitive d'une somme est la somme des primitives ; une primitive de $k f$ est $k$ fois une primitive de $f$ ($k \in \mathbb{R}$).
\item \textbf{Ajuster la constante multiplicative} : si $F_0$ est une primitive de $f_0$, alors $kF_0$ est une primitive de $kf_0$. On cherche le coefficient qui \og compense \fg{} le facteur apparaissant en dérivant.
\item \textbf{Vérifier systématiquement} en dérivant le résultat obtenu : on doit retrouver exactement $f$.
\item \textbf{Conclure} en écrivant toutes les primitives sous la forme $F(x) = F_0(x) + C$, où $C$ est une constante réelle quelconque, en précisant l'intervalle.
\end{enumerate}
\end{methode}

\begin{exoresolu}[Primitives de fonctions polynômes et puissances]
Déterminer toutes les primitives sur $\mathbb{R}$ des fonctions suivantes :
\begin{enumerate}
\item $f(x) = 3x^2 - 5x + 7$ ;
\item $g(x) = x^3 - \dfrac{4}{x^2}$ sur $]0\,;\,+\infty[$.
\end{enumerate}
\tcblower
\textbf{Solution détaillée.}

\textbf{1.} La fonction $f$ est une fonction polynôme, donc elle admet des primitives sur $\mathbb{R}$.

\emph{Étape 1 :} on primitive terme à terme (linéarité).

\emph{Étape 2 :} lecture inverse du tableau.
\begin{itemize}
\item Une primitive de $x^2$ est $\dfrac{x^3}{3}$, donc une primitive de $3x^2$ est $3\times\dfrac{x^3}{3} = x^3$.
\item Une primitive de $x$ est $\dfrac{x^2}{2}$, donc une primitive de $-5x$ est $-\dfrac{5x^2}{2}$.
\item Une primitive de la constante $7$ est $7x$.
\end{itemize}

\emph{Étape 3 :} on obtient
\[ F(x) = x^3 - \frac{5}{2}x^2 + 7x + C, \qquad C \in \mathbb{R}. \]

\emph{Étape 4 : vérification.} $F'(x) = 3x^2 - 5x + 7 = f(x)$. \checkmark

\textbf{Conclusion :} les primitives de $f$ sur $\mathbb{R}$ sont les fonctions $F(x) = x^3 - \dfrac{5}{2}x^2 + 7x + C$, $C\in\mathbb{R}$.

\textbf{2.} Sur $]0\,;\,+\infty[$, la fonction $g$ est continue (somme de fonctions continues), donc admet des primitives.
\begin{itemize}
\item Une primitive de $x^3$ est $\dfrac{x^4}{4}$.
\item Une primitive de $\dfrac{1}{x^2}$ est $-\dfrac{1}{x}$, donc une primitive de $-\dfrac{4}{x^2}$ est $-4\times\left(-\dfrac{1}{x}\right) = \dfrac{4}{x}$.
\end{itemize}
\emph{Vérification :} $\left(\dfrac{x^4}{4} + \dfrac{4}{x}\right)' = x^3 - \dfrac{4}{x^2} = g(x)$. \checkmark

\textbf{Conclusion :} les primitives de $g$ sur $]0\,;\,+\infty[$ sont les fonctions
\[ G(x) = \frac{x^4}{4} + \frac{4}{x} + C, \qquad C \in \mathbb{R}. \]
\end{exoresolu}

\courssub{Opérations sur les primitives}

\begin{propriete}[Linéarité de la primitive]
Soient $f$ et $g$ deux fonctions admettant des primitives sur un intervalle $I$, et $\lambda \in \mathbb{R}$.
\begin{itemize}
\item Une primitive de $f+g$ sur $I$ est la somme d'une primitive de $f$ et d'une primitive de $g$.
\item Une primitive de $\lambda f$ sur $I$ est $\lambda$ fois une primitive de $f$.
\end{itemize}
\end{propriete}

\begin{methode}[Primitives des formes composées usuelles]
Quand $f$ n'est pas directement dans le tableau, on cherche à reconnaître une \textbf{forme composée} $u'(x) \times (\text{fonction de } u(x))$ :
\begin{enumerate}
\item Repérer dans $f$ un \textbf{bloc} $u(x)$ et sa dérivée $u'(x)$ (à un facteur constant près).
\item Identifier la forme : $u'\,u^n$, $\dfrac{u'}{u^2}$, $\dfrac{u'}{u}$, $u'e^u$, $\dfrac{u'}{\sqrt{u}}$\dots
\item Appliquer la formule correspondante, puis \textbf{ajuster le coefficient} pour faire apparaître exactement $u'$.
\item Vérifier en dérivant (dérivation d'une composée : $(F\circ u)' = u'\times(F'\circ u)$).
\end{enumerate}
\end{methode}

\begin{aretenir}[Formes composées à connaître]
\begin{center}
\begin{tabularx}{\linewidth}{|X|X|}
\hline
\rowcolor{popblueL} \textbf{Fonction} & \textbf{Une primitive} \\
\hline
$u'\,u^n$ \ ($n\in\mathbb{N}$) & $\dfrac{u^{n+1}}{n+1}$ \\
\hline
$\dfrac{u'}{u^2}$ & $-\dfrac{1}{u}$ \\
\hline
$\dfrac{u'}{u}$ \ (avec $u>0$ sur l'intervalle) & $\ln u$ \\
\hline
$u'e^{u}$ & $e^{u}$ \\
\hline
$\dfrac{u'}{\sqrt{u}}$ \ (avec $u>0$ sur l'intervalle) & $2\sqrt{u}$ \\
\hline
\end{tabularx}
\end{center}
\end{aretenir}

\begin{exoresolu}[Reconnaître une forme $u'u^n$ et une forme $\frac{u'}{u}$]
\begin{enumerate}
\item Déterminer les primitives de $f(x) = (2x+1)(x^2+x-3)^4$ sur $\mathbb{R}$.
\item Déterminer les primitives de $g(x) = \dfrac{3}{3x-2}$ sur $\left]\dfrac{2}{3}\,;\,+\infty\right[$.
\end{enumerate}
\tcblower
\textbf{Solution détaillée.}

\textbf{1.} Posons $u(x) = x^2 + x - 3$. Alors $u'(x) = 2x + 1$.

On reconnaît exactement la forme $f = u'\times u^4$ : la dérivée $u'$ est présente \emph{sans facteur à ajuster}.

Une primitive de $u'u^4$ est $\dfrac{u^5}{5}$. Donc :
\[ F(x) = \frac{(x^2+x-3)^5}{5} + C, \qquad C\in\mathbb{R}. \]

\emph{Vérification :} par dérivation d'une composée,
\[ F'(x) = \frac{1}{5}\times 5(x^2+x-3)^4\times(2x+1) = (2x+1)(x^2+x-3)^4 = f(x). \checkmark \]

\textbf{Conclusion :} les primitives de $f$ sont les $F(x) = \dfrac{(x^2+x-3)^5}{5} + C$, $C\in\mathbb{R}$.

\textbf{2.} Posons $u(x) = 3x - 2$. Alors $u'(x) = 3$ : la dérivée apparaît exactement au numérateur. On reconnaît la forme $\dfrac{u'}{u}$.

Sur $\left]\dfrac{2}{3}\,;\,+\infty\right[$, on a $3x - 2 > 0$, donc $u > 0$ : la formule $\ln u$ s'applique (condition essentielle !).

Une primitive de $\dfrac{u'}{u}$ est $\ln u$, donc :
\[ G(x) = \ln(3x-2) + C, \qquad C\in\mathbb{R}. \]

\emph{Vérification :} $G'(x) = \dfrac{3}{3x-2} = g(x)$. \checkmark

\textbf{Conclusion :} les primitives de $g$ sur $\left]\dfrac{2}{3}\,;\,+\infty\right[$ sont les $G(x) = \ln(3x-2) + C$, $C\in\mathbb{R}$.
\end{exoresolu}

\begin{exoresolu}[Ajuster le coefficient : forme $u'e^u$]
Déterminer les primitives de $h(x) = x\,e^{x^2}$ sur $\mathbb{R}$.
\tcblower
\textbf{Solution détaillée.}

\emph{Étape 1 :} posons $u(x) = x^2$. Alors $u'(x) = 2x$.

\emph{Étape 2 :} la fonction $h(x) = x\,e^{x^2}$ ressemble à la forme $u'e^u$, mais le facteur devant l'exponentielle est $x$ et non $2x$. On écrit donc :
\[ h(x) = \frac{1}{2}\times 2x\,e^{x^2} = \frac{1}{2}\,u'(x)\,e^{u(x)}. \]

\emph{Étape 3 :} une primitive de $u'e^u$ est $e^u$, donc une primitive de $\dfrac{1}{2}u'e^u$ est $\dfrac{1}{2}e^u$ :
\[ H(x) = \frac{1}{2}e^{x^2} + C, \qquad C\in\mathbb{R}. \]

\emph{Étape 4 : vérification.} $H'(x) = \dfrac{1}{2}\times 2x\,e^{x^2} = x e^{x^2} = h(x)$. \checkmark

\textbf{Conclusion :} les primitives de $h$ sur $\mathbb{R}$ sont les $H(x) = \dfrac{1}{2}e^{x^2} + C$, $C\in\mathbb{R}$.

\textbf{Réflexe à acquérir :} quand la dérivée $u'$ est présente \og à un facteur constant près \fg{}, on divise par ce facteur. Ici $u' = 2x$ et on n'avait que $x$ : on divise par $2$.
\end{exoresolu}

\courssub{Primitive vérifiant une condition initiale}

\begin{methode}[Primitive telle que $F(a) = b$]
Pour déterminer \textbf{la} primitive $F$ de $f$ sur $I$ vérifiant $F(a) = b$ :
\begin{enumerate}
\item Vérifier que $a \in I$ (sinon le problème n'a pas de sens).
\item Déterminer \textbf{toutes} les primitives de $f$ : $F(x) = F_0(x) + C$.
\item Traduire la condition initiale : $F_0(a) + C = b$.
\item Résoudre cette équation d'inconnue $C$ : $C = b - F_0(a)$.
\item Remplacer $C$ par sa valeur et \textbf{conclure par une phrase} donnant l'unique primitive cherchée.
\end{enumerate}
\end{methode}

\begin{aretenir}[Unicité]
Si $f$ admet des primitives sur un \textbf{intervalle} $I$, alors pour tout $a\in I$ et tout réel $b$, il existe une \textbf{unique} primitive $F$ de $f$ sur $I$ telle que $F(a) = b$. C'est la condition initiale qui \og fixe \fg{} la constante.
\end{aretenir}

\begin{exoresolu}[Primitive avec condition initiale (type Bac)]
On considère la fonction $f$ définie sur $\mathbb{R}$ par $f(x) = 6x^2 - 4x + 1$.

Déterminer la primitive $F$ de $f$ sur $\mathbb{R}$ telle que $F(1) = 3$.
\tcblower
\textbf{Solution détaillée.}

\emph{Étape 1 :} $f$ est une fonction polynôme, continue sur $\mathbb{R}$, donc admet des primitives sur $\mathbb{R}$, et $1\in\mathbb{R}$.

\emph{Étape 2 :} détermination de toutes les primitives. On primitive terme à terme :
\begin{itemize}
\item primitive de $6x^2$ : $6\times\dfrac{x^3}{3} = 2x^3$ ;
\item primitive de $-4x$ : $-4\times\dfrac{x^2}{2} = -2x^2$ ;
\item primitive de $1$ : $x$.
\end{itemize}
Donc $F(x) = 2x^3 - 2x^2 + x + C$, où $C\in\mathbb{R}$.

\emph{Étape 3 :} traduction de la condition initiale $F(1) = 3$ :
\[ F(1) = 2(1)^3 - 2(1)^2 + 1 + C = 2 - 2 + 1 + C = 1 + C. \]

\emph{Étape 4 :} résolution : $1 + C = 3 \iff C = 2$.

\emph{Étape 5 : vérification.} $F(x) = 2x^3 - 2x^2 + x + 2$ ; on a bien $F'(x) = 6x^2 - 4x + 1 = f(x)$ et $F(1) = 2 - 2 + 1 + 2 = 3$. \checkmark

\textbf{Conclusion :} l'unique primitive de $f$ sur $\mathbb{R}$ vérifiant $F(1) = 3$ est définie par
\[ F(x) = 2x^3 - 2x^2 + x + 2. \]
\end{exoresolu}

\begin{exoresolu}[Condition initiale avec logarithme]
Déterminer la primitive $G$ de la fonction $g(x) = \dfrac{1}{x}$ sur $]0\,;\,+\infty[$ qui s'annule en $e$.
\tcblower
\textbf{Solution détaillée.}

\emph{Étape 1 :} $g$ est continue sur $]0\,;\,+\infty[$, donc admet des primitives sur cet intervalle, et $e \in\, ]0\,;\,+\infty[$.

\emph{Étape 2 :} par lecture inverse du tableau, les primitives de $\dfrac{1}{x}$ sur $]0\,;\,+\infty[$ sont :
\[ G(x) = \ln x + C, \qquad C\in\mathbb{R}. \]

\emph{Étape 3 :} \og qui s'annule en $e$ \fg{} signifie $G(e) = 0$. Or :
\[ G(e) = \ln e + C = 1 + C. \]

\emph{Étape 4 :} $1 + C = 0 \iff C = -1$.

\emph{Étape 5 : vérification.} $G(x) = \ln x - 1$ ; $G'(x) = \dfrac{1}{x} = g(x)$ et $G(e) = 1 - 1 = 0$. \checkmark

\textbf{Conclusion :} la primitive de $g$ sur $]0\,;\,+\infty[$ qui s'annule en $e$ est
\[ G(x) = \ln x - 1. \]
\end{exoresolu}

\courssub{Applications et synthèse}

\begin{methode}[Modéliser : retrouver une grandeur à partir de son taux de variation]
Dans les situations concrètes (physique, économie, démographie) :
\begin{enumerate}
\item Identifier que la donnée est une \textbf{dérivée} (vitesse, taux d'accroissement, coût marginal\dots) et que la grandeur cherchée en est une \textbf{primitive}.
\item Déterminer toutes les primitives de la fonction donnée.
\item Exploiter la \textbf{donnée initiale} de l'énoncé (position de départ, population initiale, coût fixe\dots) pour calculer la constante.
\item Répondre à la question posée en \textbf{calculant une valeur} de la primitive obtenue, sans oublier l'unité.
\end{enumerate}
\end{methode}

\begin{exoresolu}[Coût total à partir du coût marginal]
Une entreprise de Douala produit des sacs de ciment. Le \textbf{coût marginal} (en milliers de FCFA par tonne) pour une production de $x$ tonnes est modélisé par
\[ C_m(x) = 3x^2 - 12x + 15, \qquad x \in [0\,;\,10]. \]
On rappelle que le coût marginal est la dérivée du coût total $C_T$. Les \textbf{coûts fixes} (coût à production nulle) s'élèvent à $50$ milliers de FCFA.
\begin{enumerate}
\item Déterminer l'expression du coût total $C_T(x)$.
\item Calculer le coût total d'une production de $4$ tonnes.
\end{enumerate}
\tcblower
\textbf{Solution détaillée.}

\textbf{1.} Puisque $C_m = C_T'$, le coût total $C_T$ est une primitive du coût marginal $C_m$ sur $[0\,;\,10]$.

\emph{Étape 1 :} primitives de $C_m$ : on primitive terme à terme :
\[ C_T(x) = 3\times\frac{x^3}{3} - 12\times\frac{x^2}{2} + 15x + K = x^3 - 6x^2 + 15x + K, \quad K\in\mathbb{R}. \]

\emph{Étape 2 :} les coûts fixes traduisent la condition initiale $C_T(0) = 50$ :
\[ C_T(0) = 0 - 0 + 0 + K = K = 50. \]

\textbf{Conclusion :} le coût total est donné par
\[ C_T(x) = x^3 - 6x^2 + 15x + 50 \quad \text{(en milliers de FCFA)}. \]

\textbf{2.} Pour $x = 4$ :
\begin{align*}
C_T(4) &= 4^3 - 6\times 4^2 + 15\times 4 + 50 \\
&= 64 - 96 + 60 + 50 \\
&= 78.
\end{align*}

\textbf{Conclusion :} la production de $4$ tonnes de ciment coûte au total $78$ milliers de FCFA, soit $\num{78000}$ FCFA.
\end{exoresolu}

\begin{attention}[Les 5 erreurs qui coûtent des points au Bac]
\begin{enumerate}
\item \textbf{Oublier la constante $C$.} Quand on demande \emph{les} primitives (ou \emph{une} primitive sans condition initiale), écrire $F(x) = \dots + C$ avec $C\in\mathbb{R}$ est \textbf{obligatoire}. Sans condition initiale, une réponse sans constante est fausse ou incomplète.
\item \textbf{Confondre primitiver et dériver.} Piège classique : la primitive de $x^n$ est $\dfrac{x^{n+1}}{n+1}$ (on \emph{augmente} l'exposant), et non $nx^{n-1}$. En cas de doute, \textbf{dérive ton résultat} : tu dois retomber sur $f$. Cette vérification, écrite sur la copie, est très appréciée du correcteur.
\item \textbf{Oublier d'ajuster le coefficient dans les formes composées.} Une primitive de $e^{2x}$ n'est \textbf{pas} $e^{2x}$ mais $\dfrac{1}{2}e^{2x}$, car $(e^{2x})' = 2e^{2x}$. De même, une primitive de $\dfrac{1}{2x+1}$ est $\dfrac{1}{2}\ln(2x+1)$, pas $\ln(2x+1)$.
\item \textbf{Négliger les conditions d'existence et de signe.} La formule $\dfrac{u'}{u} \to \ln u$ n'est valable que si $u > 0$ sur l'intervalle considéré : il faut le \textbf{justifier} sur la copie. De même, vérifie que le point $a$ de la condition initiale appartient bien à l'intervalle d'étude.
\item \textbf{Se tromper dans le calcul de la constante.} Erreur fréquente : écrire $C = F_0(a) - b$ au lieu de $C = b - F_0(a)$, ou oublier d'évaluer $F_0$ en $a$. Rédige proprement l'équation $F_0(a) + C = b$, résous-la, puis \textbf{vérifie} en recalculant $F(a)$ avec ta réponse finale.
\end{enumerate}
\end{attention}

\newpage

\coursec{Exercices}

\coursec{Primitives d'une fonction}

\courssub{Notion de primitive}

\begin{definition}[Primitive d'une fonction sur un intervalle]
Soit $I$ un intervalle de $\mathbb{R}$ et $f$ une fonction définie sur $I$. On appelle \cle{primitive} de $f$ sur $I$ toute fonction $F$ dérivable sur $I$ telle que :
\[
\forall x \in I,\quad F'(x) = f(x).
\]
\end{definition}

\begin{exemplebox}[Exemples simples]
\begin{itemize}
\item La fonction $F(x) = x^3$ est une primitive de $f(x) = 3x^2$ sur $\mathbb{R}$ car $F'(x) = 3x^2 = f(x)$.
\item La fonction $F(x) = \sin x$ est une primitive de $f(x) = \cos x$ sur $\mathbb{R}$.
\item La fonction $F(x) = \ln x$ est une primitive de $f(x) = \dfrac{1}{x}$ sur $]0\,;+\infty[$.
\end{itemize}
\end{exemplebox}

\begin{propriete}[Existence des primitives]
Toute fonction \textbf{continue} sur un intervalle $I$ admet au moins une primitive sur $I$.
\end{propriete}

\begin{attention}[Condition de continuité]
La continuité sur un \emph{intervalle} est une condition \emph{suffisante} mais non nécessaire. Une fonction discontinue peut avoir des primitives (ex: fonction partie entière n'en a pas, mais $f(x)=2x\mathbf{1}_{x\neq 0}$ en a une). Au Bac, on vérifie systématiquement la continuité sur l'intervalle considéré pour garantir l'existence.
\end{attention}

\begin{popfigure}
\begin{tikzpicture}[scale=0.8, domain=-2:2]
\draw[->] (-2.5,0) -- (2.5,0) node[right] {$x$};
\draw[->] (0,-1.5) -- (0,2.5) node[above] {$y$};
\draw[popcurve, thick, domain=-2:2, samples=100] plot (\x, {2*\x}) node[right, pos=0.9] {$f(x)=2x$};
\draw[poporange, thick, domain=-2:2, samples=100] plot (\x, {\x*\x}) node[right, pos=0.9] {$F(x)=x^2$};
\draw[dashed, popmuted] (1,0) node[below] {1} -- (1,2) -- (0,2) node[left] {2};
\draw[dashed, popmuted] (-1,0) node[below] {-1} -- (-1,-2) -- (0,-2) node[left] {-2};
\node[popblue] at (1.5, 1.5) {$f = F'$};
\node[poporange] at (-1.5, -1.5) {$F$ primitive de $f$};
\end{tikzpicture}
\legende{Illustration : $f(x)=2x$ est la dérivée de $F(x)=x^2$. La pente de $F$ en $x$ vaut $f(x)$.}
\end{popfigure}

\courssub{Ensemble des primitives d'une fonction}

\begin{propriete}[Caractère local des primitives]
Soit $f$ une fonction admettant des primitives sur un intervalle $I$. Si $F$ et $G$ sont deux primitives de $f$ sur $I$, alors il existe une constante $k \in \mathbb{R}$ telle que :
\[
\forall x \in I,\quad G(x) = F(x) + k.
\]
\end{propriete}

\begin{aretenir}[Notation et ensemble des primitives]
L'ensemble de toutes les primitives de $f$ sur $I$ se note $\int f(x)\,\mathrm{d}x$ ou simplement $\int f$. Si $F_0$ est une primitive particulière, alors :
\[
\int f(x)\,\mathrm{d}x = \bigl\{ F_0(x) + k \;\big|\; k \in \mathbb{R} \bigr\}.
\]
On appelle $k$ la \cle{constante d'intégration}.
\end{aretenir}

\begin{exemplebox}[Ensemble des primitives]
Soit $f(x) = 2x$ sur $\mathbb{R}$. Une primitive est $F_0(x) = x^2$.
L'ensemble des primitives est $\int 2x\,\mathrm{d}x = \{ x^2 + k \mid k \in \mathbb{R} \}$.
\end{exemplebox}

\courssub{Primitives des fonctions usuelles}

\begin{propriete}[Tableau des primitives usuelles (à connaître par cœur)]
On obtient ce tableau par lecture inverse du tableau des dérivées. Pour tout intervalle $I$ où les expressions ont un sens :
\begin{center}
\begin{tabularx}{\linewidth}{|c|X|c|}
\hline
\rowcolor{popblueL}
\textbf{Fonction $f(x)$} & \textbf{Primitive $F(x)$ sur $I$} & \textbf{Intervalle $I$} \\
\hline
$k$ ($k \in \mathbb{R}$) & $kx$ & $\mathbb{R}$ \\
$x^n$ ($n \in \mathbb{N}$) & $\dfrac{x^{n+1}}{n+1}$ & $\mathbb{R}$ si $n$ pair, $\mathbb{R}$ si $n$ impair \\
$x^\alpha$ ($\alpha \in \mathbb{Q}\setminus\{-1\}$) & $\dfrac{x^{\alpha+1}}{\alpha+1}$ & $]0\,;+\infty[$ (ou $\mathbb{R}^*$ selon $\alpha$) \\
$\dfrac{1}{x}$ & $\ln |x|$ & $\mathbb{R}^*$ (souvent $]0\,;+\infty[$ en Tle) \\
$\mathrm{e}^x$ & $\mathrm{e}^x$ & $\mathbb{R}$ \\
$\cos x$ & $\sin x$ & $\mathbb{R}$ \\
$\sin x$ & $-\cos x$ & $\mathbb{R}$ \\
$\dfrac{1}{\cos^2 x}$ & $\tan x$ & $]-\frac{\pi}{2}+k\pi\,;\frac{\pi}{2}+k\pi[$ \\
$\dfrac{1}{1+x^2}$ & $\arctan x$ & $\mathbb{R}$ \\
$\dfrac{1}{\sqrt{1-x^2}}$ & $\arcsin x$ & $]-1\,;1[$ \\
\hline
\end{tabularx}
\end{center}
\end{propriete}

\begin{methode}[Déterminer une primitive d'une fonction usuelle]
\begin{enumerate}
\item Identifier la forme de $f(x)$ dans le tableau (puissance, exponentielle, trigonométrique, etc.).
\item Écrire la primitive correspondante $F(x)$ en oubliant la constante $k$ (on cherche \emph{une} primitive).
\item Vérifier en dérivant $F$ qu'on retrouve bien $f$.
\end{enumerate}
\end{methode}

\begin{exoresolu}[Application du tableau]
Déterminer une primitive de $f(x) = 3x^2 - \dfrac{2}{x} + 5\mathrm{e}^x$ sur $]0\,;+\infty[$.
\tcblower
On utilise la linéarité (vue après) et le tableau :
\begin{align*}
\int 3x^2\,\mathrm{d}x &= 3 \cdot \frac{x^3}{3} = x^3, \\
\int -\frac{2}{x}\,\mathrm{d}x &= -2\ln x, \\
\int 5\mathrm{e}^x\,\mathrm{d}x &= 5\mathrm{e}^x.
\end{align*}
Une primitive est $F(x) = x^3 - 2\ln x + 5\mathrm{e}^x$ sur $]0\,;+\infty[$.
Vérification : $F'(x) = 3x^2 - \frac{2}{x} + 5\mathrm{e}^x = f(x)$.
\end{exoresolu}

\courssub{Opérations sur les primitives}

\begin{propriete}[Linéarité de la primitive]
Soient $f$ et $g$ deux fonctions admettant des primitives sur un intervalle $I$, et $\lambda, \mu \in \mathbb{R}$. Alors $\lambda f + \mu g$ admet des primitives sur $I$ et :
\[
\int \bigl( \lambda f(x) + \mu g(x) \bigr)\,\mathrm{d}x = \lambda \int f(x)\,\mathrm{d}x + \mu \int g(x)\,\mathrm{d}x.
\]
En clair : une primitive de $\lambda f + \mu g$ est $\lambda F + \mu G$ où $F,G$ sont des primitives respectives de $f,g$.
\end{propriete}

\begin{attention}[Pièges fréquents]
\begin{itemize}
\item \textbf{Produit :} Une primitive de $f \times g$ n'est \textbf{pas} le produit des primitives. (Ex: $f(x)=x, g(x)=x$. Primitive de $x^2$ est $x^3/3$, mais $(x^2/2)\times(x^2/2)=x^4/4$).
\item \textbf{Quotient :} Une primitive de $f/g$ n'est \textbf{pas} le quotient des primitives.
\item \textbf{Composition :} Une primitive de $f \circ g$ n'est \textbf{pas} $F \circ G$. (C'est le domaine de l'intégration par substitution, programme Spé).
\end{itemize}
\end{attention}

\begin{exemplebox}[Utilisation de la linéarité]
$f(x) = 4x^3 - 3\sqrt{x} + \dfrac{2}{x^2}$ sur $]0\,;+\infty[$.
On écrit $f(x) = 4x^3 - 3x^{1/2} + 2x^{-2}$.
Une primitive :
$F(x) = 4\frac{x^4}{4} - 3\frac{x^{3/2}}{3/2} + 2\frac{x^{-1}}{-1} = x^4 - 2x^{3/2} - \frac{2}{x}$.
\end{exemplebox}

\courssub{Primitive vérifiant une condition initiale}

\begin{methode}[Déterminer la primitive $F$ telle que $F(a)=b$]
\begin{enumerate}
\item Déterminer l'ensemble des primitives de $f$ sur l'intervalle $I$ : $F(x) = F_0(x) + k$, $k \in \mathbb{R}$.
\item Utiliser la condition $F(a) = b$ pour trouver $k$ : $k = b - F_0(a)$.
\item Écrire l'expression finale de $F(x) = F_0(x) + k$.
\item \textbf{Vérifier} : $F'(x) = f(x)$ et $F(a) = b$.
\end{enumerate}
\end{methode}

\begin{exoresolu}[Condition initiale]
Déterminer la primitive $F$ de $f(x) = 3x^2 - 4x$ sur $\mathbb{R}$ telle que $F(1) = 2$.
\tcblower
\begin{enumerate}
\item Primitive générale : $F_0(x) = x^3 - 2x^2$. Donc $F(x) = x^3 - 2x^2 + k$.
\item Condition : $F(1) = 1 - 2 + k = 2 \Rightarrow -1 + k = 2 \Rightarrow k = 3$.
\item Réponse : $F(x) = x^3 - 2x^2 + 3$.
\item Vérification : $F'(x) = 3x^2 - 4x = f(x)$ et $F(1) = 1 - 2 + 3 = 2$. \checkmark
\end{enumerate}
\end{exoresolu}

\courssub{Applications et synthèse}

\begin{experience}[Lien Physique : Vitesse et Position]
Un mobile se déplace sur un axe. Sa vitesse $v(t)$ est la dérivée de sa position $x(t)$ : $x'(t) = v(t)$.
Donc la position $x(t)$ est une \textbf{primitive} de la vitesse $v(t)$.
La condition initiale $x(t_0) = x_0$ (position connue à un instant) permet de déterminer la constante $k$.
\begin{exemplebox}[Chute libre simplifiée]
$v(t) = -gt + v_0$ ($g \approx 10\,\mathrm{m/s^2}$). Position : $x(t) = -\frac{1}{2}gt^2 + v_0 t + x_0$.
\end{exemplebox}
\end{experience}

\begin{savaistu}[Contexte camerounais : Coût marginal]
En économie (ex: gestion d'une menuiserie à Bafoussam ou d'une plantation de cacao au Sud), le \cle{coût marginal} $C_m(q)$ est la dérivée du \cle{coût total} $C(q)$.
Connaître $C_m(q)$ et le \cle{coût fixe} $C(0)$ permet de retrouver la fonction coût total $C(q)$ par recherche de primitive avec condition initiale. C'est un outil décisionnel pour fixer les prix de vente.
\end{savaistu}

% ============================================================
% EXERCICES
% ============================================================

\courssub{Je vérifie mes connaissances}

\begin{exercice}[QCM : une seule réponse exacte]
Pour chaque question, choisir la bonne réponse (aucune justification n'est demandée).
\begin{enumerate}
\item Une primitive de $f(x) = x^2$ sur $\mathbb{R}$ est :
\begin{enumerate}
\item $F(x) = 2x$ \qquad \textbf{b)} $F(x) = \dfrac{x^3}{3}$ \qquad \textbf{c)} $F(x) = x^3$ \qquad \textbf{d)} $F(x) = \dfrac{x^2}{2}$
\end{enumerate}
\item Une primitive de $f(x) = \dfrac{1}{x^2}$ sur $]0\,;+\infty[$ est :
\begin{enumerate}
\item $F(x) = \dfrac{1}{x}$ \qquad \textbf{b)} $F(x) = -\dfrac{1}{x}$ \qquad \textbf{c)} $F(x) = \ln x$ \qquad \textbf{d)} $F(x) = \dfrac{2}{x^3}$
\end{enumerate}
\item Si $F$ est une primitive de $f$ sur $\mathbb{R}$, alors l'ensemble des primitives de $f$ sur $\mathbb{R}$ est :
\begin{enumerate}
\item réduit à $F$ seule \qquad \textbf{b)} $\{F + k \; ; \; k \in \mathbb{R}\}$ \qquad \textbf{c)} $\{kF \; ; \; k \in \mathbb{R}\}$
\end{enumerate}
\item Une primitive de $f(x) = \cos x$ sur $\mathbb{R}$ est :
\begin{enumerate}
\item $F(x) = -\sin x$ \qquad \textbf{b)} $F(x) = \sin x$ \qquad \textbf{c)} $F(x) = -\cos x$ \qquad \textbf{d)} $F(x) = \tan x$
\end{enumerate}
\end{enumerate}
\end{exercice}

\begin{exercice}[Vrai ou faux, en justifiant]
Répondre par \emph{vrai} ou \emph{faux} à chacune des affirmations suivantes, en justifiant soigneusement la réponse.
\begin{enumerate}
\item « La primitive de $x^3$ sur $\mathbb{R}$ est $3x^2$. »
\item « Si $F$ et $G$ sont deux primitives d'une même fonction $f$ sur $\mathbb{R}$, alors $F - G$ est une primitive de la fonction nulle. »
\item « La fonction $F(x) = \sqrt{x}$ est une primitive de $f(x) = \dfrac{1}{2\sqrt{x}}$ sur $]0\,;+\infty[$. »
\item « Deux primitives d'une même fonction sur un intervalle diffèrent toujours d'une constante. »
\end{enumerate}
\end{exercice}

\begin{exercice}[Vérification de primitive]
On considère les fonctions définies sur $\mathbb{R}$ par $F(x) = (2x+1)^2$ et $f(x) = 8x + 4$.
\begin{enumerate}
\item Calculer $F'(x)$ pour tout réel $x$.
\item En déduire que $F$ est une primitive de $f$ sur $\mathbb{R}$.
\item Donner alors l'ensemble de toutes les primitives de $f$ sur $\mathbb{R}$.
\end{enumerate}
\end{exercice}

\begin{exercice}[Calculs directs de primitives]
Déterminer une primitive de chacune des fonctions suivantes sur l'intervalle indiqué.
\begin{enumerate}
\item $f(x) = 5$ sur $\mathbb{R}$ ;
\item $g(x) = 6x^2 - 4x + 1$ sur $\mathbb{R}$ ;
\item $h(x) = \dfrac{3}{x}$ sur $]0\,;+\infty[$ ;
\item $k(x) = \mathrm{e}^{x} + 2x$ sur $\mathbb{R}$.
\end{enumerate}
\end{exercice}

\courssub{Je m'entraîne}

\begin{exercice}[Primitives et condition initiale]
Soit $f$ la fonction définie sur $\mathbb{R}$ par $f(x) = 4x^3 - 3x + 2$.
\begin{enumerate}
\item Déterminer l'ensemble des primitives de $f$ sur $\mathbb{R}$.
\item Déterminer la primitive $F$ de $f$ qui prend la valeur $5$ en $x = 1$.
\end{enumerate}
\end{exercice}

\begin{exercice}[Primitives avec puissances de $x$]
Soit $f$ la fonction définie sur $]0\,;+\infty[$ par $f(x) = \dfrac{3}{x^2} - \dfrac{2}{\sqrt{x}}$.
\begin{enumerate}
\item Justifier que $f$ admet des primitives sur $]0\,;+\infty[$.
\item En écrivant $f$ sous la forme $f(x) = 3x^{-2} - 2x^{-1/2}$, déterminer une primitive $F$ de $f$ sur $]0\,;+\infty[$, en justifiant chaque étape.
\item Vérifier le résultat en calculant $F'(x)$.
\end{enumerate}
\end{exercice}

\begin{exercice}[Condition initiale et trigonométrie]
Déterminer la fonction $f$ définie sur $\mathbb{R}$ telle que :
\[ f'(x) = \cos x - \sin x \quad \text{et} \quad f(0) = 2. \]
\end{exercice}

\begin{exercice}[Mouvement d'un mobile]
Un mobile se déplace sur un axe. Sa vitesse à l'instant $t$ (en secondes) est donnée, en mètres par seconde, par :
\[ v(t) = 6t^2 - 4t. \]
On note $x(t)$ la position du mobile à l'instant $t$. On rappelle que $x'(t) = v(t)$.
\begin{enumerate}
\item Déterminer l'expression de $x(t)$ sachant qu'à l'instant $t = 0$ le mobile se trouve à la position $x = 2$ m.
\item Calculer la position du mobile à l'instant $t = 3$ s.
\item Calculer la distance parcourue par le mobile entre les instants $t = 0$ et $t = 3$ s.
\end{enumerate}
\end{exercice}

\courssub{Je me prépare au Bac}

\begin{exercice}[La menuiserie de Bafoussam]
Une menuiserie de Bafoussam fabrique des tables. Le \cle{coût marginal} de production, exprimé en milliers de francs CFA, est modélisé pour une quantité $q$ de tables produites ($q \geq 0$) par :
\[ C_m(q) = 2q + 5. \]
On rappelle que le coût marginal est la dérivée du \cle{coût total} : $C'(q) = C_m(q)$. Par ailleurs, le coût fixe de l'entreprise (coût lorsque aucune table n'est produite) est de $15\,000$ FCFA, soit $15$ milliers de FCFA : $C(0) = 15$.
\begin{enumerate}
\item Déterminer l'ensemble des primitives de la fonction $C_m$ sur $[0\,;+\infty[$.
\item En utilisant la condition $C(0) = 15$, montrer que le coût total est donné par $C(q) = q^2 + 5q + 15$.
\item Calculer le coût total de production de $20$ tables, en milliers de FCFA puis en francs CFA.
\item Le menuisier vend chaque table $12\,000$ FCFA. On note $R(q) = 12q$ la recette (en milliers de FCFA) pour $q$ tables vendues.
\begin{enumerate}
\item Exprimer le bénéfice $B(q) = R(q) - C(q)$ en fonction de $q$.
\item Étudier le signe de $B(q)$ sur $[0\,;+\infty[$ et en déduire à partir de combien de tables vendues la menuiserie réalise un bénéfice positif.
\end{enumerate}
\end{enumerate}
\end{exercice}

\begin{exercice}[Forage et débit d'eau à Maroua]
Dans le cadre d'un projet d'électrification et d'adduction d'eau en zone rurale près de Maroua, on étudie le remplissage d'un château d'eau. Le débit instantané d'arrivée de l'eau, en mètres cubes par heure, est modélisé à l'instant $t$ (en heures, $t \geq 0$) par :
\[ d(t) = 3t^2 - 12t + 13. \]
On note $V(t)$ le volume d'eau (en m$^3$) arrivé dans le château d'eau à l'instant $t$. On a donc $V'(t) = d(t)$, et au début du pompage le château d'eau contient déjà $5$ m$^3$ : $V(0) = 5$.
\begin{enumerate}
\item Déterminer l'expression de $V(t)$ pour tout $t \geq 0$.
\item Calculer le volume d'eau contenu dans le château d'eau au bout de $4$ heures de pompage.
\item \begin{enumerate}
\item Étudier les variations de la fonction $d$ sur $[0\,;+\infty[$.
\item En déduire que le débit est toujours strictement positif, c'est-à-dire que le pompage ne s'interrompt pas.
\end{enumerate}
\item La capacité maximale du château d'eau est de $70$ m$^3$. Montrer que l'équation $V(t) = 70$ admet une unique solution $\alpha$ dans l'intervalle $[0\,;6]$, puis donner un encadrement de $\alpha$ d'amplitude $0{,}5$. Interpréter le résultat pour les techniciens de la CAMWATER.
\end{enumerate}
\end{exercice}

\begin{exercice}[Étude complète d'une fonction]
On considère la fonction $f$ définie sur $\mathbb{R}$ par :
\[ f(x) = x^3 - 3x + 1. \]
\begin{enumerate}
\item Déterminer l'ensemble des primitives de $f$ sur $\mathbb{R}$.
\item Déterminer la primitive $F$ de $f$ sur $\mathbb{R}$ vérifiant $F(1) = \dfrac{5}{4}$.
\item \begin{enumerate}
\item Calculer $f'(x)$ et étudier son signe sur $\mathbb{R}$.
\item Dresser le tableau de variations complet de $f$ sur $\mathbb{R}$ (on donne $f(-1) = 3$ et $f(1) = -1$).
\end{enumerate}
\item \begin{enumerate}
\item Montrer que l'équation $f(x) = 0$ admet exactement trois solutions réelles, dont une, notée $\alpha$, appartient à l'intervalle $]1\,;2[$.
\item Donner un encadrement de $\alpha$ d'amplitude $10^{-1}$.
\end{enumerate}
\item On considère maintenant la fonction $G$ définie sur $\mathbb{R}$ par $G(x) = F(x) - F(0)$, où $F$ est la fonction obtenue à la question 2.
\begin{enumerate}
\item Justifier que $G$ est encore une primitive de $f$ sur $\mathbb{R}$.
\item Calculer $G(0)$ et $G(2)$.
\end{enumerate}
\end{enumerate}
\end{exercice}

\newpage

\coursec{Corrigés des exercices}

\begin{corrige}[Exercice 1 — QCM : une seule réponse exacte]
\begin{enumerate}
\item \textbf{Réponse b)} : $F(x) = \dfrac{x^3}{3}$. En effet, $F'(x) = \dfrac{3x^2}{3} = x^2 = f(x)$. La réponse a) est la \emph{dérivée} de $f$, pas une primitive.
\item \textbf{Réponse b)} : $F(x) = -\dfrac{1}{x}$. En effet, $F'(x) = \dfrac{1}{x^2} = f(x)$. On pouvait aussi écrire $f(x) = x^{-2}$, dont une primitive est $\dfrac{x^{-1}}{-1} = -\dfrac{1}{x}$.
\item \textbf{Réponse b)} : $\{F + k \; ; \; k \in \mathbb{R}\}$. Deux primitives d'une même fonction sur un intervalle diffèrent d'une constante additive (et non multiplicative).
\item \textbf{Réponse b)} : $F(x) = \sin x$, car $(\sin x)' = \cos x$. Attention au signe : c'est la primitive de $\sin x$ qui fait intervenir $-\cos x$.
\end{enumerate}
\end{corrige}

\begin{corrige}[Exercice 2 — Vrai ou faux, en justifiant]
\begin{enumerate}
\item \textbf{Faux.} $3x^2$ est la \emph{dérivée} de $x^3$, pas une primitive. Une primitive de $x^3$ sur $\mathbb{R}$ est $F(x) = \dfrac{x^4}{4}$, car $F'(x) = \dfrac{4x^3}{4} = x^3$.
\item \textbf{Vrai.} Si $F$ et $G$ sont deux primitives de $f$ sur $\mathbb{R}$, alors pour tout réel $x$ :
\[
(F - G)'(x) = F'(x) - G'(x) = f(x) - f(x) = 0.
\]
Donc $F - G$ a pour dérivée la fonction nulle : c'est bien une primitive de la fonction nulle (c'est d'ailleurs une fonction constante).
\item \textbf{Vrai.} $F(x) = \sqrt{x}$ est dérivable sur $]0\,;+\infty[$ et $F'(x) = \dfrac{1}{2\sqrt{x}} = f(x)$. Donc $F$ est une primitive de $f$ sur $]0\,;+\infty[$. (On pouvait aussi écrire $f(x) = \tfrac{1}{2}x^{-1/2}$, de primitive $x^{1/2} = \sqrt{x}$.)
\item \textbf{Vrai.} C'est la propriété du cours : si $F$ et $G$ sont deux primitives de $f$ sur un \emph{intervalle} $I$, alors $(F-G)' = 0$ sur $I$, donc $F - G$ est constante sur $I$ : il existe $k \in \mathbb{R}$ tel que $G = F + k$. \textbf{Point de vigilance :} l'hypothèse « sur un intervalle » est indispensable ; sur une réunion d'intervalles disjoints, la conclusion est fausse.
\end{enumerate}
\end{corrige}

\begin{corrige}[Exercice 3 — Vérification de primitive]
\begin{enumerate}
\item $F(x) = (2x+1)^2$ est de la forme $u^2$ avec $u(x) = 2x+1$, $u'(x) = 2$. Donc :
\[
F'(x) = 2u'(x)\,u(x) = 2 \times 2 \times (2x+1) = 4(2x+1) = 8x + 4.
\]
\item Pour tout réel $x$, $F'(x) = 8x + 4 = f(x)$, et $F$ est dérivable sur $\mathbb{R}$. Donc $F$ est une primitive de $f$ sur $\mathbb{R}$.
\item L'ensemble de toutes les primitives de $f$ sur $\mathbb{R}$ est :
\[
\left\{ x \mapsto (2x+1)^2 + k \; ; \; k \in \mathbb{R} \right\}.
\]
\emph{Remarque :} en développant, $(2x+1)^2 + k = 4x^2 + 4x + 1 + k = 4x^2 + 4x + k'$ avec $k' = k+1$ ; on retrouve la primitive « attendue » $4x^2 + 4x$ à une constante près, ce qui confirme la cohérence.
\end{enumerate}
\end{corrige}

\begin{corrige}[Exercice 4 — Calculs directs de primitives]
\begin{enumerate}
\item $f(x) = 5$ : une primitive est $F(x) = 5x$ sur $\mathbb{R}$ (primitive de la constante $k$ : $kx$).
\item $g(x) = 6x^2 - 4x + 1$ : par linéarité,
\[
G(x) = 6 \times \frac{x^3}{3} - 4 \times \frac{x^2}{2} + x = 2x^3 - 2x^2 + x \quad \text{sur } \mathbb{R}.
\]
Vérification : $G'(x) = 6x^2 - 4x + 1 = g(x)$. \checkmark
\item $h(x) = \dfrac{3}{x}$ sur $]0\,;+\infty[$ : une primitive de $\frac{1}{x}$ est $\ln x$, donc $H(x) = 3\ln x$.
\item $k(x) = \mathrm{e}^{x} + 2x$ : une primitive de $\mathrm{e}^x$ est $\mathrm{e}^x$, une primitive de $2x$ est $x^2$. Donc $K(x) = \mathrm{e}^{x} + x^2$ sur $\mathbb{R}$.
\end{enumerate}
\end{corrige}

\begin{corrige}[Exercice 5 — Primitives et condition initiale]
\begin{enumerate}
\item $f$ est une fonction polynôme, continue sur $\mathbb{R}$, donc admet des primitives sur $\mathbb{R}$. Par linéarité :
\[
F(x) = 4 \times \frac{x^4}{4} - 3 \times \frac{x^2}{2} + 2x + k = x^4 - \frac{3}{2}x^2 + 2x + k, \quad k \in \mathbb{R}.
\]
\item Condition $F(1) = 5$ :
\[
F(1) = 1 - \frac{3}{2} + 2 + k = \frac{3}{2} + k = 5 \quad \Longrightarrow \quad k = 5 - \frac{3}{2} = \frac{7}{2}.
\]
Donc la primitive cherchée est :
\[
F(x) = x^4 - \frac{3}{2}x^2 + 2x + \frac{7}{2}.
\]
Vérification : $F'(x) = 4x^3 - 3x + 2 = f(x)$ et $F(1) = 1 - \tfrac{3}{2} + 2 + \tfrac{7}{2} = 5$. \checkmark
\end{enumerate}
\end{corrige}

\begin{corrige}[Exercice 6 — Primitives avec puissances de $x$]
\begin{enumerate}
\item Sur $]0\,;+\infty[$, les fonctions $x \mapsto \dfrac{3}{x^2}$ et $x \mapsto \dfrac{2}{\sqrt{x}}$ sont continues (quotients et composées de fonctions continues, dénominateurs non nuls). Donc $f$ est continue sur l'intervalle $]0\,;+\infty[$ : elle y admet des primitives.
\item On écrit $f(x) = 3x^{-2} - 2x^{-1/2}$. On applique la formule : une primitive de $x^{\alpha}$ ($\alpha \neq -1$) est $\dfrac{x^{\alpha+1}}{\alpha+1}$.
\begin{align*}
\text{Pour } 3x^{-2} : \quad & 3 \times \frac{x^{-2+1}}{-2+1} = 3 \times \frac{x^{-1}}{-1} = -\frac{3}{x}. \\
\text{Pour } -2x^{-1/2} : \quad & -2 \times \frac{x^{-1/2+1}}{-\tfrac{1}{2}+1} = -2 \times \frac{x^{1/2}}{1/2} = -4\sqrt{x}.
\end{align*}
Une primitive de $f$ sur $]0\,;+\infty[$ est donc :
\[
F(x) = -\frac{3}{x} - 4\sqrt{x}.
\]
\item Vérification : $F'(x) = -3 \times \left(-\dfrac{1}{x^2}\right) - 4 \times \dfrac{1}{2\sqrt{x}} = \dfrac{3}{x^2} - \dfrac{2}{\sqrt{x}} = f(x)$. \checkmark
\end{enumerate}
\end{corrige}

\begin{corrige}[Exercice 7 — Condition initiale et trigonométrie]
On cherche $f$ telle que $f'(x) = \cos x - \sin x$ : $f$ est une primitive de $x \mapsto \cos x - \sin x$.
\begin{enumerate}
\item Primitive générale : une primitive de $\cos x$ est $\sin x$ ; une primitive de $-\sin x$ est $\cos x$ (car $(\cos x)' = -\sin x$). Donc :
\[
f(x) = \sin x + \cos x + k, \quad k \in \mathbb{R}.
\]
\item Condition $f(0) = 2$ : $f(0) = \sin 0 + \cos 0 + k = 0 + 1 + k = 2$, d'où $k = 1$.
\item Conclusion : $\boxed{f(x) = \sin x + \cos x + 1}$.
\end{enumerate}
Vérification : $f'(x) = \cos x - \sin x$ et $f(0) = 0 + 1 + 1 = 2$. \checkmark
\end{corrige}

\begin{corrige}[Exercice 8 — Mouvement d'un mobile]
\begin{enumerate}
\item $x$ est une primitive de $v$ avec $v(t) = 6t^2 - 4t$ :
\[
x(t) = 6 \times \frac{t^3}{3} - 4 \times \frac{t^2}{2} + k = 2t^3 - 2t^2 + k.
\]
Condition $x(0) = 2$ : $k = 2$. Donc :
\[
x(t) = 2t^3 - 2t^2 + 2 \quad \text{(en mètres)}.
\]
\item À $t = 3$ : $x(3) = 2 \times 27 - 2 \times 9 + 2 = 54 - 18 + 2 = 38$ m.
\item Le \textbf{déplacement} (variation de position) entre $t = 0$ et $t = 3$ est $x(3) - x(0) = 38 - 2 = 36$ m.
La \textbf{distance parcourue} (longueur de la trajectoire) demande d'étudier le signe de la vitesse.
$v(t) = 6t^2 - 4t = 2t(3t-2)$ s'annule en $t = 0$ et $t = \tfrac{2}{3}$.
Sur $]0\,;\tfrac{2}{3}[$, $v(t) < 0$ (le mobile recule), puis $v(t) > 0$ sur $]\tfrac{2}{3}\,;3]$. Le mobile change donc de sens en $t = \frac{2}{3}$.
Calculons la position à cet instant :
\[
x\left(\tfrac{2}{3}\right) = 2 \times \frac{8}{27} - 2 \times \frac{4}{9} + 2 = \frac{16}{27} - \frac{24}{27} + \frac{54}{27} = \frac{46}{27} \approx 1{,}70 \text{ m}.
\]
Distance réellement parcourue :
\[
\left|x\left(\tfrac{2}{3}\right) - x(0)\right| + \left|x(3) - x\left(\tfrac{2}{3}\right)\right| = \left|\frac{46}{27} - 2\right| + \left|38 - \frac{46}{27}\right| = \frac{8}{27} + \frac{980}{27} = \frac{988}{27} \approx 36{,}6 \text{ m}.
\]

\emph{Point de vigilance :} ne pas confondre déplacement et distance parcourue. Si l'énoncé demande la distance parcourue, il faut sommer les valeurs absolues des déplacements sur chaque intervalle de monotonie de $x$ (signe constant de $v$).
\end{enumerate}
\end{corrige}

\begin{corrige}[Exercice 9 — La menuiserie de Bafoussam]
\emph{Barème indicatif : question 1 : 1 pt ; question 2 : 1,5 pt ; question 3 : 1 pt ; question 4 : 2,5 pts.}
\begin{enumerate}
\item $C_m(q) = 2q + 5$ est continue sur $[0\,;+\infty[$, donc admet des primitives. L'ensemble des primitives est :
\[
C(q) = q^2 + 5q + k, \quad k \in \mathbb{R}.
\]
\item Condition $C(0) = 15$ : $0 + 0 + k = 15$, donc $k = 15$. Ainsi :
\[
C(q) = q^2 + 5q + 15 \quad \text{(en milliers de FCFA)}.
\]
\item $C(20) = 400 + 100 + 15 = 515$ milliers de FCFA, soit $\mathbf{515\,000}$ \textbf{FCFA} pour 20 tables.
\item \begin{enumerate}
\item $B(q) = R(q) - C(q) = 12q - (q^2 + 5q + 15) = -q^2 + 7q - 15$.
\item Signe de $B(q) = -q^2 + 7q - 15$ : discriminant $\Delta = 49 - 60 = -11 < 0$. Le trinôme n'a pas de racine réelle et est du signe de $a = -1 < 0$ : $B(q) < 0$ pour tout $q \geq 0$.

\textbf{Conclusion :} avec un prix de vente de $12\,000$ FCFA par table, la menuiserie ne réalise \emph{jamais} de bénéfice positif : $B(q) < 0$ pour toute quantité. Le bénéfice maximal est atteint en $q = \tfrac{7}{2} = 3{,}5$ et vaut $B(3{,}5) = -12{,}25 + 24{,}5 - 15 = -2{,}75$ milliers de FCFA : même au mieux, l'entreprise perd $2\,750$ FCFA. Le menuisier doit augmenter son prix de vente ou réduire ses coûts.
\end{enumerate}
\end{enumerate}
\textbf{Points de vigilance :} citer la continuité pour justifier l'existence des primitives ; ne pas oublier la conversion milliers $\to$ FCFA ; pour le signe du trinôme, justifier par le discriminant.
\end{corrige}

\begin{corrige}[Exercice 10 — Forage et débit d'eau à Maroua]
\emph{Barème indicatif : question 1 : 1,5 pt ; question 2 : 1 pt ; question 3 : 2 pts ; question 4 : 2,5 pts.}
\begin{enumerate}
\item $V$ est une primitive de $d(t) = 3t^2 - 12t + 13$ :
\[
V(t) = t^3 - 6t^2 + 13t + k.
\]
Condition $V(0) = 5$ : $k = 5$. Donc :
\[
V(t) = t^3 - 6t^2 + 13t + 5 \quad (t \geq 0).
\]
\item $V(4) = 64 - 96 + 52 + 5 = 25$ m$^3$. Au bout de 4 heures, le château d'eau contient $25$ m$^3$.
\item \begin{enumerate}
\item $d(t) = 3t^2 - 12t + 13$ est un trinôme ; $d'(t) = 6t - 12$, qui s'annule en $t = 2$. $d'(t) < 0$ sur $[0\,;2[$ et $d'(t) > 0$ sur $]2\,;+\infty[$. Donc $d$ est décroissante sur $[0\,;2]$ puis croissante sur $[2\,;+\infty[$, avec un minimum en $t = 2$ :
\[
d(2) = 12 - 24 + 13 = 1.
\]
\item Le minimum de $d$ sur $[0\,;+\infty[$ est $d(2) = 1 > 0$. Donc pour tout $t \geq 0$, $d(t) \geq 1 > 0$ : le débit est toujours strictement positif, le pompage ne s'interrompt jamais.
\end{enumerate}
\item $V'(t) = d(t) > 0$ sur $[0\,;6]$, donc $V$ est \textbf{continue et strictement croissante} sur $[0\,;6]$. De plus :
\[
V(0) = 5 < 70 \quad \text{et} \quad V(6) = 216 - 216 + 78 + 5 = 83 > 70.
\]
Donc $70 \in [V(0)\,;V(6)]$ : par le \textbf{théorème des valeurs intermédiaires} (cas des fonctions strictement monotones), l'équation $V(t) = 70$ admet une \textbf{unique} solution $\alpha \in [0\,;6]$.

Encadrement : $V(5) = 125 - 150 + 65 + 5 = 45 < 70$ et $V(5{,}5) = 166{,}375 - 181{,}5 + 71{,}5 + 5 = 61{,}375 < 70$ ; $V(6) = 83 > 70$. Donc $5{,}5 < \alpha < 6$ : encadrement d'amplitude $0{,}5$.

\textbf{Interprétation :} les techniciens de la CAMWATER savent que le château d'eau sera plein au bout d'environ 5 h 30 min à 6 h de pompage ; il faudra prévoir l'arrêt ou la régulation du pompage avant ce délai pour éviter le débordement.
\end{enumerate}
\textbf{Points de vigilance :} pour le TVI, citer explicitement la continuité, la stricte monotonie et l'appartenance de 70 à l'intervalle image ; vérifier les calculs numériques de $V$ aux bornes.
\end{corrige}

\begin{corrige}[Exercice 11 — Étude complète d'une fonction]
\emph{Barème indicatif : question 1 : 1 pt ; question 2 : 1,5 pt ; question 3 : 2 pts ; question 4 : 2,5 pts ; question 5 : 1,5 pt.}
\begin{enumerate}
\item $f$ est polynomiale, donc continue sur $\mathbb{R}$. L'ensemble des primitives de $f(x) = x^3 - 3x + 1$ est :
\[
F(x) = \frac{x^4}{4} - \frac{3x^2}{2} + x + k, \quad k \in \mathbb{R}.
\]
\item Condition $F(1) = \dfrac{5}{4}$ :
\[
F(1) = \frac{1}{4} - \frac{3}{2} + 1 + k = -\frac{1}{4} + k = \frac{5}{4} \quad \Longrightarrow \quad k = \frac{6}{4} = \frac{3}{2}.
\]
Donc $F(x) = \dfrac{x^4}{4} - \dfrac{3x^2}{2} + x + \dfrac{3}{2}$.
Vérification : $F(1) = \tfrac{1}{4} - \tfrac{3}{2} + 1 + \tfrac{3}{2} = \tfrac{1}{4} + 1 = \tfrac{5}{4}$. \checkmark
\item \begin{enumerate}
\item $f'(x) = 3x^2 - 3 = 3(x^2 - 1) = 3(x-1)(x+1)$. C'est un trinôme du second degré s'annulant en $-1$ et $1$, positif à l'extérieur des racines : $f'(x) > 0$ sur $]-\infty\,;-1[$ et $]1\,;+\infty[$ ; $f'(x) < 0$ sur $]-1\,;1[$.
\item Limites : $\displaystyle\lim_{x \to -\infty} f(x) = -\infty$ et $\displaystyle\lim_{x \to +\infty} f(x) = +\infty$ (terme dominant $x^3$). Tableau de variations :
\begin{center}
\begin{tabular}{|c|ccccccc|}
\hline
$x$ & $-\infty$ & & $-1$ & & $1$ & & $+\infty$ \\
\hline
$f'(x)$ & & $+$ & $0$ & $-$ & $0$ & $+$ & \\
\hline
$f$ & $-\infty$ & $\nearrow$ & $3$ & $\searrow$ & $-1$ & $\nearrow$ & $+\infty$ \\
\hline
\end{tabular}
\end{center}
\end{enumerate}
\item \begin{enumerate}
\item Sur $]-\infty\,;-1]$ : $f$ est continue, strictement croissante de $-\infty$ vers $3$ ; $0 \in {]-\infty\,;3]}$, donc une unique solution $\alpha_1$.
Sur $[-1\,;1]$ : $f$ est continue, strictement décroissante de $3$ vers $-1$ ; $0 \in [-1\,;3]$, donc une unique solution $\alpha_2$.
Sur $[1\,;+\infty[$ : $f$ est continue, strictement croissante de $-1$ vers $+\infty$ ; $0 \in [-1\,;+\infty[$, donc une unique solution $\alpha$.
Au total, l'équation $f(x) = 0$ admet exactement \textbf{trois} solutions réelles. Comme $f(1) = -1 < 0$ et $f(2) = 8 - 6 + 1 = 3 > 0$, la troisième solution $\alpha$ appartient à $]1\,;2[$.
\item $f(1{,}5) = 3{,}375 - 4{,}5 + 1 = -0{,}125 < 0$ et $f(1{,}6) = 4{,}096 - 4{,}8 + 1 = 0{,}296 > 0$. Donc :
\[
1{,}5 < \alpha < 1{,}6 \quad \text{(amplitude } 10^{-1}\text{)}.
\]
\end{enumerate}
\item \begin{enumerate}
\item $G(x) = F(x) - F(0)$ : $G$ diffère de $F$ d'une constante, donc $G'(x) = F'(x) = f(x)$ pour tout réel $x$, et $G$ est dérivable sur $\mathbb{R}$ : $G$ est bien une primitive de $f$ sur $\mathbb{R}$.
\item $G(0) = F(0) - F(0) = 0$.
$F(0) = \dfrac{3}{2}$ et $F(2) = \dfrac{16}{4} - \dfrac{12}{2} + 2 + \dfrac{3}{2} = 4 - 6 + 2 + \dfrac{3}{2} = \dfrac{3}{2}$.
Donc $G(2) = F(2) - F(0) = \dfrac{3}{2} - \dfrac{3}{2} = 0$.
\end{enumerate}
\end{enumerate}
\textbf{Points de vigilance :} pour la question 4.a), le TVI doit être appliqué \emph{sur chacun des trois intervalles de monotonie} séparément, avec continuité et stricte monotonie citées ; pour 5.a), il suffit d'observer que $G = F + \text{constante}$.
\end{corrige}

\newpage

\coursec{Fiche bilan}

\begin{popfigure}
\begin{tikzpicture}[mindmap, grow cyclic, every node/.style={concept}, concept color=popblue,
  root concept/.append style={concept color=popblue, font=\bfseries, text=white},
  level 1/.append style={level distance=3.6cm, sibling angle=60},
  level 2/.append style={level distance=2.2cm, sibling angle=45, font=\scriptsize}]
\node[root concept]{Primitives d'une fonction}
  child[concept color=poporange]{ node {Définition}
    child { node {$F' = f$ sur $I$} }
    child { node {$f$ continue $\Rightarrow$ primitive existe} }
  }
  child[concept color=popgreen]{ node {Ensemble des primitives}
    child { node {Deux primitives diffèrent d'une constante} }
    child { node {$F + k$, $k \in \mathbb{R}$} }
  }
  child[concept color=poppurple]{ node {Primitives usuelles}
    child { node {Lecture inverse du tableau des dérivées} }
    child { node {$x^n$, $\frac{1}{x}$, $e^x$, $\cos$, $\sin$, $\frac{1}{1+x^2}$, $\frac{1}{\sqrt{1-x^2}}$} }
  }
  child[concept color=popgold]{ node {Opérations}
    child { node {Linéarité : somme, multiple} }
    child { node {Formes composées : $u'u^n$, $\frac{u'}{u}$, $u'e^u$, \ldots} }
  }
  child[concept color=poppink]{ node {Condition initiale}
    child { node {$F(a) = b$ fixe $k$} }
    child { node {Unicité de la primitive} }
  }
  child[concept color=popmuted]{ node {Applications}
    child { node {Vérifier par dérivation} }
    child { node {Problèmes concrets (croissance, population, \ldots)} }
  };
\end{tikzpicture}
\legende{Carte mentale de la leçon : les six idées forces autour de la notion de primitive.}
\end{popfigure}

\coursec{Notion de primitive}

\begin{definition}[Primitive d'une fonction sur un intervalle]
Soit $I$ un intervalle de $\mathbb{R}$ et $f$ une fonction définie sur $I$.
On appelle \cle{primitive} de $f$ sur $I$ toute fonction $F$ \textbf{dérivable} sur $I$ telle que :
\[
\forall x \in I,\quad F'(x) = f(x).
\]
On note souvent $F' = f$ sur $I$.
\end{definition}

\begin{propriete}[Existence des primitives (théorème admis)]
Toute fonction \cle{continue} sur un intervalle $I$ admet au moins une primitive sur $I$.
\end{propriete}

\begin{attention}[Conditions d'application]
\begin{itemize}
  \item L'intervalle $I$ doit être \textbf{explicitement précisé} (ex : $]0;+\infty[$, $\mathbb{R}$, $[-1;1]$).
  \item La fonction $f$ doit être \textbf{continue} sur $I$ pour garantir l'existence (condition suffisante, non nécessaire).
  \item Une primitive est \textbf{toujours continue} sur $I$ (car dérivable).
\end{itemize}
\end{attention}

\begin{exemplebox}[Premiers exemples]
\begin{enumerate}
  \item $f(x) = 2x$ sur $\mathbb{R}$. Une primitive est $F(x) = x^2$ car $F'(x) = 2x = f(x)$.
  \item $f(x) = \cos x$ sur $\mathbb{R}$. Une primitive est $F(x) = \sin x$ car $F'(x) = \cos x$.
  \item $f(x) = \frac{1}{x}$ sur $]0;+\infty[$. Une primitive est $F(x) = \ln x$ car $F'(x) = \frac{1}{x}$.
  \item $f(x) = \frac{1}{x}$ sur $]-\infty;0[$. Une primitive est $F(x) = \ln(-x)$ car $F'(x) = \frac{-1}{-x} = \frac{1}{x}$.
\end{enumerate}
\end{exemplebox}

\coursec{Ensemble des primitives d'une fonction}

\begin{propriete}[Caractérisation de l'ensemble des primitives]
Soit $f$ une fonction continue sur un intervalle $I$ et $F_0$ une primitive de $f$ sur $I$.
Alors l'ensemble de \textbf{toutes} les primitives de $f$ sur $I$ est :
\[
\{ F_0 + k \mid k \in \mathbb{R} \}.
\]
En d'autres termes, deux primitives d'une même fonction sur un intervalle ne diffèrent que d'une \cle{constante}.
\end{propriete}

\begin{aretenir}[Conséquence fondamentale]
\begin{itemize}
  \item Si on demande « \emph{une} primitive », on donne une fonction particulière (souvent celle avec $k=0$).
  \item Si on demande « \emph{les} primitives » ou « \emph{l'ensemble des} primitives », on \textbf{doit} écrire $F(x) = F_0(x) + k,\quad k \in \mathbb{R}$.
  \item L'oubli de la constante $k$ (ou de la mention $k \in \mathbb{R}$) est une erreur de rédaction sanctionnée au Baccalauréat.
\end{itemize}
\end{aretenir}

\begin{exemplebox}[Ensemble des primitives]
Soit $f(x) = 3x^2$ sur $\mathbb{R}$.
Une primitive est $F_0(x) = x^3$.
L'ensemble des primitives est $\{ x \mapsto x^3 + k \mid k \in \mathbb{R} \}$.
\end{exemplebox}

\coursec{Primitives des fonctions usuelles}

\begin{propriete}[Tableau des primitives usuelles (à connaître par cœur)]
On obtient ce tableau par \textbf{lecture inverse} du tableau des dérivées.
\begin{center}
\begin{tabularx}{\linewidth}{|c|X|c|}
\hline
\rowcolor{popblue!40} \textbf{Fonction $f(x)$} & \textbf{Une primitive $F(x)$} & \textbf{Intervalle de validité} \\
\hline
$k$ (constante) & $kx$ & $\mathbb{R}$ \\
\hline
$x^n$ \ ($n \in \mathbb{Z}$, $n \neq -1$) & $\dfrac{x^{n+1}}{n+1}$ & $\mathbb{R}$ si $n \ge 0$ ; $]-\infty;0[ \cup ]0;+\infty[$ si $n < 0$ \\
\hline
$\dfrac{1}{x}$ & $\ln|x|$ & $\mathbb{R}^*$ (ou $\ln x$ sur $]0;+\infty[$) \\
\hline
$\dfrac{1}{\sqrt{x}}$ & $2\sqrt{x}$ & $]0;+\infty[$ \\
\hline
$e^x$ & $e^x$ & $\mathbb{R}$ \\
\hline
$\cos x$ & $\sin x$ & $\mathbb{R}$ \\
\hline
$\sin x$ & $-\cos x$ & $\mathbb{R}$ \\
\hline
$\dfrac{1}{1+x^2}$ & $\arctan x$ & $\mathbb{R}$ \\
\hline
$\dfrac{1}{\sqrt{1-x^2}}$ & $\arcsin x$ & $]-1;1[$ \\
\hline
\end{tabularx}
\end{center}
\end{propriete}

\begin{methode}[Lecture inverse du tableau des dérivées]
\begin{enumerate}
  \item Identifier la forme de $f(x)$ dans la première colonne du tableau.
  \item Écrire la primitive correspondante $F(x)$ de la deuxième colonne.
  \item \textbf{Vérifier systématiquement} en dérivant $F$ : on doit retrouver exactement $f$.
  \item Préciser l'intervalle $I$ (celui de la 3ème colonne, ou un sous-intervalle).
\end{enumerate}
\end{methode}

\begin{attention}[Pièges classiques sur les domaines]
\begin{itemize}
  \item $\dfrac{1}{x}$ : primitive $\ln|x|$ sur $\mathbb{R}^*$ (deux branches distinctes). Sur $]0;+\infty[$, on écrit $\ln x$.
  \item $x^n$ avec $n < 0$ : la fonction n'est pas définie en $0$. L'intervalle ne doit \textbf{pas} contenir $0$ (ex : $]0;+\infty[$ ou $]-\infty;0[$).
  \item $\dfrac{1}{\sqrt{1-x^2}}$ : définie seulement sur $]-1;1[$.
  \item Ne jamais écrire $\ln x$ pour une primitive sur un intervalle contenant des négatifs.
\end{itemize}
\end{attention}

\begin{exemplebox}[Application du tableau]
\begin{enumerate}
  \item $f(x) = 5x^4$ sur $\mathbb{R}$. Forme $x^n$ avec $n=4$. Primitive : $F(x) = \frac{5x^5}{5} = x^5$. Vérif : $F'(x) = 5x^4 = f(x)$.
  \item $f(x) = \frac{3}{x}$ sur $]0;+\infty[$. Forme $\frac{1}{x}$. Primitive : $F(x) = 3\ln x$. Vérif : $F'(x) = 3\cdot\frac{1}{x} = f(x)$.
  \item $f(x) = \frac{2}{1+x^2}$ sur $\mathbb{R}$. Forme $\frac{1}{1+x^2}$. Primitive : $F(x) = 2\arctan x$. Vérif : $F'(x) = 2\cdot\frac{1}{1+x^2} = f(x)$.
  \item $f(x) = \frac{1}{\sqrt{x}}$ sur $]0;+\infty[$. Primitive : $F(x) = 2\sqrt{x}$. Vérif : $F'(x) = 2\cdot\frac{1}{2\sqrt{x}} = \frac{1}{\sqrt{x}}$.
\end{enumerate}
\end{exemplebox}

\coursec{Opérations sur les primitives}

\begin{propriete}[Linéarité de la primitive]
Soit $f$ et $g$ deux fonctions continues sur un intervalle $I$, admettant respectivement $F$ et $G$ comme primitives sur $I$. Soit $\lambda \in \mathbb{R}$.
\begin{itemize}
  \item \textbf{Somme} : $F + G$ est une primitive de $f + g$ sur $I$.
  \item \textbf{Multiple par un scalaire} : $\lambda F$ est une primitive de $\lambda f$ sur $I$.
\end{itemize}
En combinant les deux : une primitive de $\lambda f + \mu g$ est $\lambda F + \mu G$.
\end{propriete}

\begin{propriete}[Formes composées usuelles (règle de la chaîne inversée)]
Soit $u$ une fonction dérivable sur un intervalle $I$, et $n \in \mathbb{Z}$, $n \neq -1$.
\begin{center}
\begin{tabular}{|c|c|c|}
\hline
\rowcolor{popblue!40} \textbf{Forme de $f(x)$} & \textbf{Une primitive $F(x)$} & \textbf{Condition} \\
\hline
$u'(x) \cdot [u(x)]^n$ & $\dfrac{[u(x)]^{n+1}}{n+1}$ & $n \neq -1$ \\
\hline
$\dfrac{u'(x)}{u(x)}$ & $\ln|u(x)|$ & $u(x) \neq 0$ sur $I$ \\
\hline
$u'(x) \cdot e^{u(x)}$ & $e^{u(x)}$ & -- \\
\hline
$\dfrac{u'(x)}{1+[u(x)]^2}$ & $\arctan(u(x))$ & -- \\
\hline
$\dfrac{u'(x)}{\sqrt{1-[u(x)]^2}}$ & $\arcsin(u(x))$ & $|u(x)| < 1$ sur $I$ \\
\hline
\end{tabular}
\end{center}
\end{propriete}

\begin{methode}[Reconnaissance des formes composées]
\begin{enumerate}
  \item Repérer une fonction $u(x)$ et sa dérivée $u'(x)$ dans l'expression de $f(x)$.
  \item Identifier quelle forme du tableau ci-dessus correspond à $f(x)$.
  \item Appliquer la formule correspondante.
  \item \textbf{Vérifier en dérivant} (règle de la chaîne).
\end{enumerate}
\end{methode}

\begin{exemplebox}[Exemples de formes composées]
\begin{enumerate}
  \item $f(x) = 2x \cdot (x^2+1)^3$ sur $\mathbb{R}$.
  $u(x) = x^2+1$, $u'(x) = 2x$. Forme $u' \cdot u^3$ avec $n=3$.
  Primitive : $F(x) = \frac{(x^2+1)^4}{4}$. Vérif : $F'(x) = \frac{4(x^2+1)^3 \cdot 2x}{4} = 2x(x^2+1)^3$.
  
  \item $f(x) = \frac{2x}{x^2+3}$ sur $\mathbb{R}$.
  $u(x) = x^2+3$, $u'(x) = 2x$. Forme $\frac{u'}{u}$.
  Primitive : $F(x) = \ln(x^2+3)$ (car $x^2+3 > 0$). Vérif : $F'(x) = \frac{2x}{x^2+3}$.
  
  \item $f(x) = \frac{\cos x}{\sin x}$ sur $]0;\pi[$.
  $u(x) = \sin x$, $u'(x) = \cos x$. Forme $\frac{u'}{u}$.
  Primitive : $F(x) = \ln(\sin x)$ (car $\sin x > 0$ sur $]0;\pi[$). Vérif : $F'(x) = \frac{\cos x}{\sin x}$.
  
  \item $f(x) = \frac{3x^2}{1+x^6}$ sur $\mathbb{R}$.
  $u(x) = x^3$, $u'(x) = 3x^2$. Forme $\frac{u'}{1+u^2}$.
  Primitive : $F(x) = \arctan(x^3)$. Vérif : $F'(x) = \frac{3x^2}{1+(x^3)^2} = \frac{3x^2}{1+x^6}$.
\end{enumerate}
\end{exemplebox}

\coursec{Primitive vérifiant une condition initiale}

\begin{propriete}[Existence et unicité]
Soit $f$ continue sur un intervalle $I$, $a \in I$ et $b \in \mathbb{R}$.
Il existe une \textbf{unique} primitive $F$ de $f$ sur $I$ telle que $F(a) = b$.
\end{propriete}

\begin{methode}[Détermination de la primitive avec condition initiale $F(a)=b$]
\begin{enumerate}
  \item Trouver une primitive $F_0$ de $f$ sur $I$ (par tableau ou formes composées).
  \item Écrire la forme générale : $F(x) = F_0(x) + k$, avec $k \in \mathbb{R}$.
  \item Utiliser la condition : $F(a) = F_0(a) + k = b \implies k = b - F_0(a)$.
  \item Écrire la primitive cherchée : $F(x) = F_0(x) + b - F_0(a)$.
  \item \textbf{Vérifier} : $F'(x) = f(x)$ et $F(a) = b$.
\end{enumerate}
\end{methode}

\begin{exoresolu}[Primitive avec condition initiale]
\textbf{Énoncé :} Déterminer la primitive $F$ de la fonction $f(x) = 3x^2 - 2\sin x$ sur $\mathbb{R}$ telle que $F(0) = 5$.

\textbf{Solution détaillée :}
\begin{enumerate}
  \item \textbf{Primitive de $f$ :} Par linéarité et tableau :
  $f(x) = 3x^2 - 2\sin x$.
  Une primitive de $3x^2$ est $x^3$.
  Une primitive de $-2\sin x$ est $2\cos x$ (car $(\cos x)' = -\sin x$).
  Donc $F_0(x) = x^3 + 2\cos x$ est une primitive de $f$ sur $\mathbb{R}$.
  
  \item \textbf{Forme générale :} $F(x) = x^3 + 2\cos x + k$, $k \in \mathbb{R}$.
  
  \item \textbf{Condition initiale :} $F(0) = 0^3 + 2\cos 0 + k = 2 + k$.
  On impose $F(0) = 5 \implies 2 + k = 5 \implies k = 3$.
  
  \item \textbf{Conclusion :} La primitive cherchée est $F(x) = x^3 + 2\cos x + 3$ sur $\mathbb{R}$.
  
  \item \textbf{Vérification :} $F'(x) = 3x^2 - 2\sin x = f(x)$ et $F(0) = 0 + 2 + 3 = 5$. \checkmark
\end{enumerate}
\end{exoresolu}

\coursec{Applications et synthèse}

\begin{exoresolu}[Exercice type Bac : Primitive et forme composée]
\textbf{Énoncé :} Soit la fonction $f$ définie sur $]0;+\infty[$ par $f(x) = \frac{2\ln x}{x}$.
\begin{enumerate}
  \item Démontrer que $f$ admet des primitives sur $]0;+\infty[$.
  \item Déterminer l'ensemble des primitives de $f$ sur $]0;+\infty[$.
  \item Déterminer la primitive $F$ de $f$ telle que $F(1) = 4$.
\end{enumerate}

\textbf{Solution détaillée :}
\begin{enumerate}
  \item $f$ est continue sur $]0;+\infty[$ (quotient de fonctions continues, dénominateur non nul). Donc $f$ admet des primitives sur $]0;+\infty[$ (théorème d'existence).
  
  \item \textbf{Recherche d'une primitive $F_0$ :}
  Posons $u(x) = \ln x$. Alors $u'(x) = \frac{1}{x}$.
  On a $f(x) = 2 \ln x \cdot \frac{1}{x} = 2 u(x) u'(x) = u'(x) \cdot 2u(x)$.
  C'est de la forme $u' \cdot u^n$ avec $n=1$ (coefficient 2).
  Une primitive est $F_0(x) = 2 \cdot \frac{(\ln x)^2}{2} = (\ln x)^2$.
  \textit{Vérification :} $F_0'(x) = 2\ln x \cdot \frac{1}{x} = f(x)$. \checkmark
  
  L'ensemble des primitives est $\{ x \mapsto (\ln x)^2 + k \mid k \in \mathbb{R} \}$.
  
  \item \textbf{Condition initiale $F(1)=4$ :}
  $F(x) = (\ln x)^2 + k$.
  $F(1) = (\ln 1)^2 + k = 0 + k = k$.
  On veut $F(1)=4 \implies k = 4$.
  La primitive cherchée est $F(x) = (\ln x)^2 + 4$ sur $]0;+\infty[$.
\end{enumerate}
\end{exoresolu}

\begin{exoresolu}[Exercice type Bac : Problème concret - Croissance d'une population]
\textbf{Énoncé :} La population $P(t)$ (en milliers d'habitants) d'une ville du Cameroun évolue selon le modèle :
$P'(t) = 0,05 \sqrt{t}$ pour $t \in [0; 100]$ (où $t$ est le temps en années depuis 2020).
On sait qu'en 2020 ($t=0$), la population était de $200\,000$ habitants.
\begin{enumerate}
  \item Exprimer $P(t)$ en fonction de $t$.
  \item Quelle sera la population en 2025 ? (Arrondir à l'entier le plus proche).
\end{enumerate}

\textbf{Solution détaillée :}
\begin{enumerate}
  \item $P$ est une primitive de $f(t) = 0,05 \sqrt{t} = 0,05 t^{1/2}$ sur $[0;100]$.
  Une primitive de $t^{1/2}$ est $\frac{t^{3/2}}{3/2} = \frac{2}{3} t^{3/2}$.
  Donc $P_0(t) = 0,05 \times \frac{2}{3} t^{3/2} = \frac{0,1}{3} t^{3/2} = \frac{1}{30} t^{3/2}$.
  Forme générale : $P(t) = \frac{1}{30} t^{3/2} + k$.
  Condition initiale : $P(0) = 200$ (en milliers).
  $P(0) = 0 + k = 200 \implies k = 200$.
  \textbf{Réponse :} $P(t) = \frac{1}{30} t^{3/2} + 200$ (population en milliers d'habitants).
  
  \item En 2025, $t = 5$.
  $P(5) = \frac{1}{30} \cdot 5^{3/2} + 200 = \frac{1}{30} \cdot 5\sqrt{5} + 200 \approx \frac{11,18}{30} + 200 \approx 0,373 + 200 = 200,373$ milliers.
  Soit environ $200\,373$ habitants.
\end{enumerate}
\end{exoresolu}

\begin{exercice}[Exercices d'entraînement]
\begin{enumerate}
  \item Déterminer une primitive sur l'intervalle indiqué :
  \begin{enumerate}
    \item $f(x) = 4x^3 - 5x + 2$ sur $\mathbb{R}$.
    \item $g(x) = \frac{3}{x} - 2\sin x$ sur $]0;\pi[$.
    \item $h(x) = \frac{2x}{x^2+1} + \frac{1}{\sqrt{1-x^2}}$ sur $]0;1[$.
    \item $k(x) = 3x^2 e^{x^3}$ sur $\mathbb{R}$.
  \end{enumerate}
  
  \item Déterminer l'ensemble des primitives sur l'intervalle indiqué :
  \begin{enumerate}
    \item $f(x) = \frac{1}{x^2}$ sur $]0;+\infty[$.
    \item $g(x) = \cos(2x)$ sur $\mathbb{R}$.
  \end{enumerate}
  
  \item Trouver la primitive $F$ de $f(x) = 2x + \frac{1}{x}$ sur $]0;+\infty[$ telle que $F(1) = 3$.
  
  \item Soit $f(x) = \frac{2x+1}{x^2+x+1}$ sur $\mathbb{R}$.
  \begin{enumerate}
    \item Montrer que $f$ s'écrit sous la forme $\frac{u'}{u}$.
    \item En déduire une primitive de $f$ sur $\mathbb{R}$.
    \item Déterminer la primitive $F$ telle que $F(0) = 0$.
  \end{enumerate}
\end{enumerate}
\end{exercice}

\begin{corrige}[Exercices d'entraînement]
\begin{enumerate}
  \item \textbf{Primitives :}
  \begin{enumerate}
    \item $F(x) = x^4 - \frac{5}{2}x^2 + 2x$ sur $\mathbb{R}$.
    \item $G(x) = 3\ln x + 2\cos x$ sur $]0;\pi[$.
    \item $H(x) = \ln(x^2+1) + \arcsin x$ sur $]0;1[$.
    \item $K(x) = e^{x^3}$ sur $\mathbb{R}$ (forme $u'e^u$ avec $u=x^3$).
  \end{enumerate}
  
  \item \textbf{Ensembles des primitives :}
  \begin{enumerate}
    \item $f(x) = x^{-2}$. Primitive $F_0(x) = -x^{-1} = -\frac{1}{x}$.
    Ensemble : $\{ x \mapsto -\frac{1}{x} + k \mid k \in \mathbb{R} \}$ sur $]0;+\infty[$.
    \item $g(x) = \cos(2x)$. Avec $u=2x$, $u'=2$ : $g(x) = \frac{1}{2} \cdot 2\cos(2x) = \frac{1}{2}\,u'\cos u$.
    Primitive $G_0(x) = \frac{1}{2}\sin(2x)$.
    Ensemble : $\{ x \mapsto \frac{1}{2}\sin(2x) + k \mid k \in \mathbb{R} \}$ sur $\mathbb{R}$.
  \end{enumerate}
  
  \item $f(x) = 2x + \frac{1}{x}$ sur $]0;+\infty[$.
  $F_0(x) = x^2 + \ln x$.
  $F(x) = x^2 + \ln x + k$.
  $F(1) = 1 + 0 + k = 3 \implies k = 2$.
  \textbf{Réponse :} $F(x) = x^2 + \ln x + 2$ sur $]0;+\infty[$.
  
  \item \begin{enumerate}
    \item $u(x) = x^2+x+1$, $u'(x) = 2x+1$. $u(x) = (x+\frac{1}{2})^2 + \frac{3}{4} > 0$ sur $\mathbb{R}$.
    $f(x) = \frac{u'(x)}{u(x)}$.
    \item Primitive : $F_0(x) = \ln(x^2+x+1)$ sur $\mathbb{R}$.
    \item $F(x) = \ln(x^2+x+1) + k$.
    $F(0) = \ln 1 + k = k = 0 \implies k=0$.
    \textbf{Réponse :} $F(x) = \ln(x^2+x+1)$ sur $\mathbb{R}$.
  \end{enumerate}
\end{enumerate}
\end{corrige}

\begin{aretenir}[Checklist avant le Bac]
\begin{itemize}
  \item $\square$ Je sais énoncer la définition d'une primitive : $F' = f$ sur un intervalle $I$.
  \item $\square$ Je sais que la continuité sur $I$ garantit l'existence de primitives (théorème admis).
  \item $\square$ Je connais par cœur le tableau des primitives usuelles (y compris $e^x$, $\arctan$, $\arcsin$).
  \item $\square$ Je n'oublie \textbf{jamais} la constante $k$ (et $k \in \mathbb{R}$) quand on demande \emph{toutes} les primitives.
  \item $\square$ Je sais que deux primitives d'une même fonction sur un intervalle diffèrent d'une constante.
  \item $\square$ Je sais déterminer $k$ grâce à une condition initiale $F(a) = b$.
  \item $\square$ Je sais primitiver une somme terme à terme et sortir un coefficient constant (linéarité).
  \item $\square$ Je reconnais les formes composées : $u'u^n$, $\frac{u'}{u}$, $u'e^u$, $\frac{u'}{1+u^2}$, $\frac{u'}{\sqrt{1-u^2}}$.
  \item $\square$ Je vérifie \textbf{toujours} mon résultat en dérivant : je dois retrouver exactement $f$.
  \item $\square$ Je fais attention au domaine : $\frac{1}{x} \to \ln|x|$ sur $\mathbb{R}^*$, $\frac{1}{\sqrt{x}}$ sur $]0;+\infty[$, $\arcsin$ sur $]-1;1[$.
  \item $\square$ Je rédige proprement : « Une primitive de $f$ sur $I$ est $F$ définie par $F(x) = \ldots$ ».
  \item $\square$ Je ne confonds pas « dériver » et « primitiver » : ce sont deux lectures opposées du même tableau.
\end{itemize}
\end{aretenir}

\begin{savaistu}[Le saviez-vous ?]
\begin{itemize}
  \item Le symbole $\int f(x)\,dx$ (intégrale indéfinie) désigne précisément l'ensemble des primitives de $f$. On écrira plus tard $\int f(x)\,dx = F(x) + C$.
  \item Au Cameroun, les ingénieurs de la SONEL (électricité) ou de CAMWATER (eau) utilisent les primitives pour calculer la quantité totale d'énergie ou d'eau consommée à partir de la puissance ou du débit instantané (qui sont des dérivées).
  \item Historiquement, le calcul intégral (recherche de primitives) est né avant le calcul différentiel (dérivées) pour résoudre des problèmes d'aires et de volumes (Archimède, puis Newton/Leibniz).
\end{itemize}
\end{savaistu}

\findecours

\end{document}
